\documentclass[11pt,letterpaper]{article}
\usepackage[T1]{fontenc}
\usepackage[utf8]{inputenc}
\usepackage{lmodern}
\usepackage[margin=1in,headheight=15pt]{geometry}
\usepackage{amsmath,amssymb,amsthm,mathtools,mathrsfs}
\usepackage{microtype,booktabs,longtable,array,enumitem,aliascnt}
\usepackage{xcolor}
\usepackage{fancyhdr}
\usepackage[colorlinks=true,linkcolor=blue!40!black,citecolor=green!30!black,urlcolor=blue!40!black]{hyperref}
\usepackage[nameinlink,noabbrev]{cleveref}
\hypersetup{pdftitle={Trigonometry on the Surcomplex Plane},%
pdfsubject={Canonical phases, arbitrary-scale geometry, and infinitesimal degeneration}}
\pagestyle{fancy}\fancyhf{}
\fancyhead[L]{\small Trigonometry on the surcomplex plane}
\fancyhead[R]{\small\thepage}
\renewcommand{\headrulewidth}{0.3pt}
\setlength{\parskip}{0.25em}
\setlength{\emergencystretch}{2em}
\setlist{itemsep=0.2em,topsep=0.3em}
\setcounter{tocdepth}{2}
\numberwithin{equation}{section}
\newtheorem{theorem}{Theorem}[section]
% Shared numbering through alias counters, so that \cref names each result by its own kind.
\newaliascnt{lemma}{theorem}\newtheorem{lemma}[lemma]{Lemma}\aliascntresetthe{lemma}
\newaliascnt{proposition}{theorem}\newtheorem{proposition}[proposition]{Proposition}
\aliascntresetthe{proposition}
\newaliascnt{corollary}{theorem}\newtheorem{corollary}[corollary]{Corollary}
\aliascntresetthe{corollary}
\theoremstyle{definition}
\newaliascnt{definition}{theorem}\newtheorem{definition}[definition]{Definition}
\aliascntresetthe{definition}
\newaliascnt{example}{theorem}\newtheorem{example}[example]{Example}\aliascntresetthe{example}
\theoremstyle{remark}
\newaliascnt{remark}{theorem}\newtheorem{remark}[remark]{Remark}\aliascntresetthe{remark}
\crefname{lemma}{lemma}{lemmas}
\crefname{proposition}{proposition}{propositions}
\crefname{corollary}{corollary}{corollaries}
\crefname{definition}{definition}{definitions}
\crefname{example}{example}{examples}
\crefname{remark}{remark}{remarks}
\newcommand{\No}{\mathbf{No}}
\newcommand{\SC}{\mathbf{SC}}
\newcommand{\Oz}{\mathbf{Oz}}
\newcommand{\R}{\mathbb R}
\newcommand{\C}{\mathbb C}
\newcommand{\Q}{\mathbb Q}
\newcommand{\Z}{\mathbb Z}
\newcommand{\N}{\mathbb N}
\newcommand{\OO}{\mathcal O}
\newcommand{\mm}{\mathfrak m}
\newcommand{\TT}{\mathbb T}
\newcommand{\PP}{\mathbb P}
\newcommand{\Dd}{\mathbb D}
\newcommand{\Hy}{\mathbb H}
\newcommand{\eps}{\varepsilon}
\newcommand{\st}{\operatorname{st}}
\newcommand{\fin}{\operatorname{fin}}
\newcommand{\supp}{\operatorname{supp}}
\newcommand{\lc}{\operatorname{lc}}
\newcommand{\re}{\operatorname{Re}}
\newcommand{\im}{\operatorname{Im}}
\newcommand{\cis}{\operatorname{cis}}
\newcommand{\Arg}{\operatorname{Arg}}
\newcommand{\arcosh}{\operatorname{arcosh}}
\newcommand{\artanh}{\operatorname{artanh}}
\newcommand{\Exp}{\operatorname{Exp}}
\newcommand{\Ex}{\operatorname{Exp}_{\mm}}
\newcommand{\Logm}{\operatorname{Log}_{\mm}}
\newcommand{\Sin}{\operatorname{Sin}}
\newcommand{\Cos}{\operatorname{Cos}}
\newcommand{\Tan}{\operatorname{Tan}}
\newcommand{\SO}{\operatorname{SO}}
\newcommand{\Orth}{\operatorname{O}}
\newcommand{\GL}{\operatorname{GL}}
\newcommand{\End}{\operatorname{End}}
\newcommand{\Rot}{\operatorname{Rot}}
\newcommand{\Gm}{\mathbb G_{\mathrm m}}
\newcommand{\Ga}{\mathbb G_{\mathrm a}}
\newcommand{\ellm}{\ell_{\mm}}
\newcommand{\Tm}{\TT_{\mm}}
\newcommand{\Hom}{\operatorname{Hom}}
\newcommand{\Mult}{\operatorname{Mult}}
\newcommand{\pr}{\operatorname{pr}}
\newcommand{\res}{\operatorname{res}}
\newcommand{\pd}{\operatorname{pd}}
\newcommand{\Sh}{\operatorname{Sh}}
\newcommand{\cl}{\operatorname{cl}}
\newcommand{\Zhat}{\widehat{\Z}}
\newcommand{\Qv}{\Q^{\vee}}
\newcommand{\cI}{\mathcal I}
\newcommand{\cL}{\mathcal L}
\newcommand{\sh}{^{\sharp}}
\newcommand{\HInt}{\mathop{\overset{\scriptscriptstyle H}{\int}}\nolimits}
\newcommand{\abs}[1]{\left|#1\right|}
\newcommand{\ang}[2]{\angle(#1,#2)}
\newcolumntype{L}[1]{>{\raggedright\arraybackslash}p{#1}}

\begin{document}
\begin{titlepage}
\centering
\vspace*{0.2cm}
{\large\scshape A merged theorem-driven development\par}
\vspace{0.3cm}
{\LARGE\bfseries Trigonometry on the\\[0.2em] Surcomplex Plane\par}
\vspace{0.3cm}
{\Large Canonical Phases, Arbitrary-Scale Geometry,\\[0.15em]
and Infinitesimal Degeneration\par}
\vspace{0.3cm}
{\large September 21, 2026; extended September 22, 2026\par}
\vspace{0.3cm}
\begin{minipage}{0.94\textwidth}\small
\begin{abstract}
The surcomplex plane is the algebraically closed class field $\SC=\No[i]$ with Euclidean norm
$|x+iy|=\sqrt{x^2+y^2}$. Its lengths may be infinite or infinitesimal, but every direction has a
\emph{finite} surreal angle. We make that observation into a trigonometric theory. Hahn summation
gives canonical sine and cosine on finite surreal arguments, and identifies the full unit circle
$\TT(\No)$ with the additive group of finite surreals modulo $2\pi\Z$; the circle splits
canonically as an ordinary phase times an additive infinitesimal phase, and its angular and
chordal distances have the same valuation.

We prove exact analogues of the angle-sum, cosine, sine, tangent, half-angle, Heron, Euler,
angle-bisector, trigonometric Ceva, inscribed-angle and Ptolemy theorems at arbitrary surreal
scale, together with the spherical sine and cosine laws and hyperbolic triangle laws in the full
surreal Poincar\'e disk. Valuation formulas convert each of these into an exact scale dictionary;
a relative flatness threshold and a quadratic triangle-inequality defect formula quantify
degeneration; and near-tangent intersections exhibit exact square-root splitting with a quantified
loss of angular precision. Three different stability certificates are proved and kept apart: a
sharp $\sigma>2\kappa$ root-stability bound with its sharpness proposition, a conditioned
inverse-cosine expansion with second-order term, and a support-controlled exact cluster
multiplicity. A canonical complex extension exists on the strip of arguments with finite real
part and arbitrary surreal imaginary part; there the sine is surjective with completely determined
fibres, trigonometric polynomials algebraize to ordinary-degree polynomials, and one obtains sharp
root counts, finite Fourier inversion and Parseval identities, and a Fej\'er--Riesz factorization
over the real closed surreal field, with a degree-two polynomial that is positive at every ordinary
angle and negative at a finite nonordinary one.

Finally we classify the freedom in extending circular trigonometry to infinite real arguments.
The Ehrlich--Kaplan integer-part normalization is one member of a character family; every member
necessarily acquires infinite periods, and none admits a bounded common-domain Hahn-analytic
description at infinite frequency. Two further sections treat the
rotation group $\SO(2,\No)$ and the arithmetic of the periods, where exponentials with the initial
kernel-multiplier ring $\Z$ give a partial negative answer to a robustness question of Ehrlich and
Kaplan, which stays open.
\end{abstract}
\end{minipage}
\vfill
\begin{minipage}{0.94\textwidth}\small
\textbf{Status and provenance.} This is one merged report assembled from six independently
written manuscripts of September 21--22, 2026 (\cref{trigonometry:sec:provenance}). It develops three supplied unpublished manuscripts
rather than presenting their results as published literature. Classical identities and established
surreal analytic machinery are identified as such. The literature search behind it was targeted
rather than exhaustive; no priority claim and no claim to have solved a named published open
problem is made; nothing here is formally verified, and the accompanying symbolic checks validate
displayed examples rather than machine-checking the general proofs. Per-theorem non-claims are
stated where they arise and collected in \cref{trigonometry:sec:notproved}.
\end{minipage}
\end{titlepage}
\setcounter{page}{2}
\begingroup
\small
\setlength{\parskip}{0pt}
\tableofcontents
\endgroup
\clearpage

\section{The organizing distinction: lengths versus angles}\label{trigonometry:sec:intro}

Let $z=x+iy$ with $x,y\in\No$ and $i^2=-1$; this is the meaning of a surcomplex number
throughout. The field structure and real closedness of $\No$, and the algebraic closedness of
$\SC$, are established parts of surreal number theory \cite{Conway,Gonshor,Alling,vdDE,Poonen}.
They already give distances with nonnegative square roots, two-dimensional determinants measuring
oriented area, and exact solutions of finite linear systems.

The analytic difficulty is different. The expression
\[
 \sum_{n\ge0}\frac{(-1)^nx^{2n+1}}{(2n+1)!}
\]
cannot simply be declared to define $\sin x$ for every surreal $x$. At a nonzero ordinary real
argument, infinitely many summands already share one Hahn exponent. At a positive infinite
argument the monomial supports run in the wrong direction and their union is not well ordered.
Ordinary real convergence repairs neither defect in the surreal fine topology.

There is nevertheless a canonical and sufficient theory of geometric angles. If $z\ne0$, its
normalized direction $z/|z|$ has coordinates in $[-1,1]$, hence an ordinary standard part on the
usual unit circle together with an infinitesimal angular correction. Consequently the trigonometry
of \emph{all} surcomplex triangles requires no choice of $\sin\omega$. The construction applies to
$1+i\omega$, to $\omega+i$, and to $\omega^{-\omega}+i\omega^{\omega}$ exactly as it applies
to $1+i$.

Three assertions must be kept apart, and this article keeps them apart everywhere.
\begin{enumerate}
\item An angle has a cosine and a sine defined as the coordinates of a unit vector. This is
 algebra over a real closed field; it needs no transcendental evaluation.
\item A \emph{finite} surreal number has a canonical cosine and sine, obtained from the ordinary
 analytic germs by strongly summable infinitesimal Taylor evaluation.
\item Every surreal number, $\omega$ included, has a uniquely specified oscillatory phase
 compatible with those finite values.
\end{enumerate}
We prove that (1) and (2) agree and suffice for all plane, spherical and hyperbolic geometry
below. Assertion (3) requires additional data; \cref{trigonometry:sec:global} says exactly which
data, and what is lost whichever choice is made, and \cref{trigonometry:per:sec:main} studies the
arithmetic of the periods that every such choice creates.

\begin{center}
\small
\begin{tabular}{L{0.23\textwidth}L{0.32\textwidth}L{0.35\textwidth}}
\toprule
Object & Domain used here & Governing construction\\
\midrule
Lengths and coordinates & All of $\No$ and $\SC$ & Ordered-field algebra and square roots\\
Circular angles & Finite $\theta\in\No$, modulo $2\pi\Z$ & Ordinary value plus infinitesimal Taylor series\\
Inverse tangent & Every $x\in\No$ & Direction of $1+ix$\\
Complex sine and cosine & $\{x+iy:x\text{ finite},\ y\in\No\}$ & Finite phase and Gonshor exponential\\
Trigonometric polynomials & Ordinary finite degree, arbitrary surcomplex coefficients & Laurent-polynomial algebra on the unit circle\\
Sine at infinite real inputs & Only after an extension is specified & An additional additive-group character\\
\bottomrule
\end{tabular}
\end{center}

\subsection{What comes from the supplied manuscripts}\label{trigonometry:sec:provenance}
The supplied analysis manuscript \cite{SourceAnalysis} provides the notation
$t^\gamma=\omega^{-\gamma}$ for Conway normal-form monomials, the set-sized support convention,
strong summability, standard parts, Taylor evaluation on infinitesimal thickenings, coefficientwise
contour and integral functionals, the fine-local logarithm, and examples of nonunique global
fine-analytic exponentials. We retain those conventions and reprove the elementary parts needed
here.

The two further supplied manuscripts \cite{SourceGeometry,SourceIntersections} concern finite
\emph{zero algebras} and analytic intersections rather than Euclidean triangle geometry. Their
relevant themes are conservation of multiplicity, positive-support coefficient recursion,
collision-stable residues, and loss of precision near a singular Jacobian; the threshold
$\sigma>2\kappa$ of \cite[\S9]{SourceIntersections} is the direct antecedent of
\cref{trigonometry:thm:stability} and \cref{trigonometry:thm:conditioned}. We do not assume their
several-variable Noetherianity, division, Nullstellensatz or residue theorems to prove the
geometric results below; each trigonometric instance is proved directly.

Ehrlich and Kaplan construct canonical global trigonometric functions and a canonical surcomplex
exponential from the initial integer part \cite[\S11]{EK}; that construction is recorded in
\cref{trigonometry:def:globaltrig} and \cref{trigonometry:thm:globalexp} and is \emph{not} claimed
as a new definition. Restricted analytic evaluation on surreal fields is likewise established
\cite[Theorem 2.1]{vdDE}. Classical trigonometric identities and inverse-function expansions are
recorded in \cite{DLMF}; algebraic approaches to trigonometry over general fields are in
\cite{Wildberger,Gunn}; the classical scalar Fej\'er--Riesz context is in \cite{DR}. A proof of an
analogue is not by itself evidence that its exact formulation is absent from all previous
literature.

\emph{The six component manuscripts.} The report itself is built from six manuscripts, numbered
locally 16--21; the numbers prefix their code and data files. Sources 16, 17 and 18 --- \emph{Finite
Radians, Arbitrary-Scale Geometry, and Infinitesimal Angular Phenomena}; \emph{Canonical Phases,
Arbitrary-Scale Triangles, and Infinitesimal Degeneration}; \emph{Canonical Angles, Infinite
Triangles, Infinitesimal Contact, and Hahn-Analytic Oscillation}, all of September 21, 2026 --- are
three drafts of one article on the trigonometry of $\SC$; 18 is the base text of
Sections~\ref{trigonometry:sec:intro}--\ref{trigonometry:sec:global}, and 16 and 17 contribute the results
the directory's README lists. They were written from the three supplied manuscripts above and record
no repository commit. Source 19, \emph{Rotations of the Surreal Plane: The Norm-One Torus,
Infinitesimal Angles, and Topology at Every Scale} (September 22, 2026), was written against
repository commit \texttt{dcf86662b574}, at which this article stood unchanged; it contributes
\cref{trigonometry:rot:sec:main}. Its structure theorem re-derives \cref{trigonometry:prop:rotation}
and \cref{trigonometry:thm:polar}, and \cref{trigonometry:rot:sec:overlap} lists every result it
shares with sources 16--18, each printed once with both credited. Sources 20 and 21, \emph{Exponential
Kernels and Arithmetic Rigidity in the Surcomplex Field} and \emph{Period Arithmetic and
Noncanonical Surcomplex Exponentials} (both September 22, 2026), were written independently against
repository commit \texttt{465a54b479a1}, at which this article stood unchanged; they build on
\cref{trigonometry:thm:characters,trigonometry:thm:infiniteperiods} and together contribute
\cref{trigonometry:per:sec:main}, with source 21 as its base text. \Cref{trigonometry:per:sec:overlap}
lists what they share, the credits they owe, and where the merge had to choose.

\subsection{Conventions: one symbol per concept}\label{trigonometry:sec:conventions}
Because sources 16--18 are three drafts with three notations, the following choices are fixed once
and used without exception thereafter. Synonyms appearing in those manuscripts (``the sources''
below) are recorded here and nowhere else; those of source 19 are recorded in
\cref{trigonometry:rot:sec:overlap}, and those of sources 20 and 21 in \cref{trigonometry:per:sec:overlap}.

\begin{itemize}
\item \textbf{Unit circle.} $\TT(\No)=\{u\in\SC:u\bar u=1\}$, with $\TT(\R)$ its ordinary complex
 counterpart. The sources also write $\mathbb U$ and $S^1(\C)$ for these; we do not.
\item \textbf{Finite phase.} $\cis\theta=\cos\theta+i\sin\theta$ for finite real $\theta$. The
 sources also write $E(\theta)$ and $\Phi(\theta)$ for this map; the letter $E$ is reserved below
 for the strip exponential of \cref{trigonometry:thm:stripexp}.
\item \textbf{Finite and infinitesimal elements.} $\OO_{\R},\mm_{\R}\subset\No$ and
 $\OO_{\C},\mm_{\C}\subset\SC$. The sources also write $\OO,\mm$ without subscripts.
\item \textbf{Standard part.} The map $\st$ to $\C$ on finite elements, and the topology it
 induces, are called the \emph{standard part} and the \emph{standard-part topology}. The source
 manuscripts variously call this map or topology the reduction, the shadow and the residue; those
 three words are not used again. The genuinely different \emph{fine} topology --- generated by all
 balls of positive surreal radius --- keeps its own name throughout.
\item \textbf{Triangles.} Vertices $A,B,C$; opposite sides $a=|B-C|$, $b=|C-A|$, $c=|A-B|$;
 interior angles $\alpha,\beta,\gamma\in(0,\pi)$; area $\Delta$, circumradius $R$, inradius $r$,
 semiperimeter $s$. One source uses the letters $A,B,C$ for both vertices and angles; we do not.
\item \textbf{Angle defect and angular metric.} $d(\alpha)=\min\{\alpha,\pi-\alpha\}$ is the
 \emph{defect} of an interior angle (a source writes $e_A$, another $d_\alpha$).
 $\ang uv\in[0,\pi]$ is the \emph{angular metric} on $\TT(\No)$ (a source writes $d_\theta$).
\item \textbf{Small parameters.} $h$ denotes a generic nonzero infinitesimal increment or angle;
 $\eps$ a generic positive infinitesimal; $\tau$ the ramification parameter of a cosine extremum
 or a near tangency, as in $\arccos(1-\tau)$; $\delta$ the side gap $b+c-a$ of
 \cref{trigonometry:thm:flatrelative}. The sources overload $\delta$ across all four roles.
\item \textbf{Asymptotics.} For surreal or surcomplex $a,b$ with $b\ne0$,
 $a\prec b$ means $a/b\in\mm_{\C}$, and $a\sim b$ means
 $a/b-1\in\mm_{\C}$. These are exact algebraic relations, never limits
 along a sequence. For a fixed nonzero infinitesimal $h$, $O(h^m)$ denotes a remainder whose
 quotient by $h^m$ is finite; it is an algebraic size statement, not a statement about an ordinary
 integer tending to infinity.
\item \textbf{Rotation-group symbols.} Source 19 uses its own symbols for the objects above and
 for the new ones of \cref{trigonometry:rot:sec:main}; they are converted once, in
 \cref{trigonometry:rot:sec:overlap}, and not used elsewhere.
\item \textbf{Period-arithmetic symbols.} Angles are radians everywhere, including
 \cref{trigonometry:per:sec:main}. Sources 20 and 21 measure phases in turns; their objects are
 converted once, in \cref{trigonometry:per:sec:overlap}, where periods divided by $2\pi$ form the
 normalized period group $\cI$.
\item \textbf{Ordinary finiteness.} Every integer $n$, polynomial degree $N$ or $n$, Fourier grid
 size $M$, and polygon size in this article is an \emph{ordinary} finite integer unless the
 contrary is said explicitly. ``Finite coefficients'' of a trigonometric polynomial means instead
 that each Laurent coefficient lies in $\OO_{\C}$, i.e.\ is finite \emph{in magnitude}; this second
 sense is used only in \cref{trigonometry:thm:stability} and is flagged there.
\end{itemize}

\subsection{Conventions: which coefficient ring each statement lives over}
\label{trigonometry:sec:rings}
Several distinct rings of Hahn-analytic data circulate in this subject under nearly identical
notation, and theorems about one of them are not theorems about another. Fix a set-sized divisible
ordered additive subgroup $\Gamma\subset\No$ and write $t^\gamma=\omega^{-\gamma}$, so that
$F_\Gamma=\R((t^\Gamma))$ is real closed and $K_\Gamma=\C((t^\Gamma))$ is algebraically closed
\cite{Poonen,vdDE}. In one variable $z$ the four rings are:
\begin{description}
\item[(R1) the fixed-domain ring $\mathcal O(J)((t^\Gamma))$.] A datum is
 $F=\sum_{\gamma\in S}f_\gamma t^\gamma$ with $S\subset\Gamma$ one well-ordered support and
 \emph{every} coefficient $f_\gamma$ analytic on one fixed ordinary domain $J$. In the real case
 $J$ is an ordinary real interval or neighbourhood; this is also the \emph{uniform-support} ring
 used for integration along an ordinary parameter interval.
\item[(R2) the common-domain germ ring.] The same, with $J$ allowed to shrink: the coefficients
 are germs at a point admitting \emph{some} common domain of definition.
\item[(R3) the radius-free ring $\C\{z\}((t^\Gamma))$.] Each $f_\gamma$ is an individually
 convergent germ, with \emph{no} common radius of convergence required.
\item[{(R4) the formal ring $\C[[z]]((t^\Gamma))$.}] Each $f_\gamma$ is a formal power series.
\end{description}
The inclusions $\text{(R1)}\subset\text{(R2)}\subset\text{(R3)}\subset\text{(R4)}$ are all strict; the
separating witnesses are not reproduced here, no statement below depending on them, and are
proved in the companion report on infinitesimal analytic geometry.
\emph{Every coherence statement in this article names its ring.} Concretely: the obstruction
theorems \cref{trigonometry:thm:infinitefrequency} and \cref{trigonometry:thm:allscale} are stated
over (R1)--(R2) and over nothing else; the coefficientwise integral of
\cref{trigonometry:sec:arcs} uses the uniform-support avatar of (R1) along an ordinary real
interval; and no theorem below is stated over the radius-free ring (R3) or the formal ring (R4),
which appear here only to locate the two that are used. No statement is transferred from one of
these rings to another anywhere in this article.

By contrast, the \emph{scalar} Taylor lift of \cref{trigonometry:eq:lift}, which carries the bulk
of the article, is not a statement about any of these rings: it evaluates one ordinary germ at one
finite surreal point, and the only hypothesis is that the germ be defined near the standard part of
that point.

\subsection{The workspace, and what it does not authorize}
Every normal-form support, coefficient family and recursion used here is a \emph{set}; $\No$ is a
proper class and a class theory such as NBG may be used. Any finite tuple of surcomplex numbers
lies in some set-sized divisible workspace $K_\Gamma$, obtained as the divisible hull of the group
generated by its normal-form supports, and any fixed calculation may be placed there. This
enlargement rule never authorizes a proper-class sum or a proper-class polynomial ring, and no
sum below has a proper-class index set. Workspace enlargement does not change any formula: the
coefficients and strong sums have the same normal forms after embedding into a larger Hahn field,
and the algebraic factorization and root statements are interpreted in the full field $\SC$, so
none of them is tied to a rank-one subfield.

\section{Surreal scalars and support-controlled evaluation}\label{trigonometry:sec:foundations}

\subsection{Norm, valuation, and standard part}
Conjugation, norm, inner product and oriented determinant are
\begin{align}
 \overline{x+iy}&=x-iy,& |z|&=\sqrt{z\bar z},\\
 \langle z,w\rangle&=\re(\bar zw),&[z,w]&=\im(\bar zw).
 \label{trigonometry:eq:dotcross}
\end{align}
Expanding coordinates gives the Gram identity
\begin{equation}\label{trigonometry:eq:gram}
 \langle z,w\rangle^2+[z,w]^2=|z|^2|w|^2,
\end{equation}
whence $|zw|=|z||w|$, the Cauchy--Schwarz inequality $|\langle z,w\rangle|\le|z||w|$, and the
triangle inequality $|z+w|\le|z|+|w|$. Equality in the last, for nonzero $z,w$, holds exactly when
$w/z$ is a positive element of $\No$; squaring the equality gives
$\langle z,w\rangle=|z||w|$, which selects the nonnegative proportionality factor. Here and below,
``real'' as a coordinate condition means belonging to $\No$, while ``ordinary real'' means
belonging to $\R$.

For a nonzero normal form $a=\sum_{\gamma\in S}a_\gamma t^\gamma$ with $S$ well ordered, put
\[
 v(a)=\min S,\qquad \lc(a)=a_{v(a)},\qquad v(0)=+\infty .
\]
Then $v(ab)=v(a)+v(b)$, $v(a+b)\ge\min(v(a),v(b))$, and $v(|a|)=v(a)$. The value group is not
assumed Archimedean; every comparison of valuations below uses only its ordered-group operations.
Write
\[
 \OO_{\R}=\{x:|x|<n\text{ for some }n\in\N_{>0}\},\qquad
 \mm_{\R}=\{x:|x|<1/n\text{ for every }n\in\N_{>0}\},
\]
and $\OO_{\C}=\OO_{\R}+i\OO_{\R}$, $\mm_{\C}=\mm_{\R}+i\mm_{\R}$. Every finite element decomposes
uniquely as
\begin{equation}\label{trigonometry:eq:standard}
 z=c+\eta,\qquad c=\st(z)\in\C,\quad\eta\in\mm_{\C},
\end{equation}
with $c\in\R$ for real elements. Standard part is a ring homomorphism on finite elements, and
equals the coefficient at exponent zero. A finite element is \emph{appreciable} when its standard
part is nonzero; a finite nonzero element has valuation zero exactly when it is appreciable.

\subsection{Strong summability and the support lemma}
\begin{definition}[Strong summability]\label{trigonometry:def:strong}
A set-indexed Hahn family $(A_j)_{j\in J}$ is \emph{strongly summable} when the union of its
normal-form supports is well ordered \emph{and} every exponent occurs in the supports of only
finitely many members. Its sum is formed by finite addition at each exponent.
\end{definition}
The two clauses are independent, and both are needed.

\begin{lemma}[Hahn--Neumann support calculus]\label{trigonometry:lem:support}
If $S,T$ are well-ordered subsets of an ordered abelian group, then $S+T$ is well ordered and each
exponent has only finitely many representations $s+t$. If $T\subset\Gamma_{>0}$ is well ordered,
the monoid
\[
 T^*=\{\tau_1+\cdots+\tau_k:k\ge0,\ \tau_j\in T\}
\]
is well ordered, and each exponent has only finitely many representations as a finite word in $T$,
including only finitely many possible word lengths. In particular, for finitely many infinitesimals
$\eta_1,\dots,\eta_d\in\mm_{\C}$ and \emph{arbitrary} ordinary complex coefficients $c_\alpha$, the
family
\[
 \sum_{\alpha\in\N^d}c_\alpha\eta_1^{\alpha_1}\cdots\eta_d^{\alpha_d}
\]
is strongly summable. Products and formal substitutions with zero constant term may be regrouped
coefficientwise, and $S+T^*$ is a support certificate for the result.
\end{lemma}
\begin{proof}
The first assertion follows by extracting a nondecreasing subsequence in one coordinate: either a
decreasing sequence of sums, or infinitely many representations of one sum, would force the other
coordinate to decrease. The positive-monoid assertion is Neumann's lemma \cite{Neumann}. One
convenient proof uses the well-quasi-ordering of finite words in a well-ordered alphabet under
subsequence embedding with coordinatewise increase: such an embedding cannot decrease a sum of
positive letters, and equality forces identical letters and identical length. This excludes both
descending sums and infinitely many distinct nondecreasing words of a fixed sum, and each word has
only finitely many permutations. Applying this to the union of the positive supports of the
$\eta_j$ gives the last assertions.
\end{proof}

Three negative companions to \cref{trigonometry:lem:support} are load-bearing and are stated here
once, because every proof below tacitly relies on them.

\begin{remark}[What strong summability is not]\label{trigonometry:rem:notlimit}
\emph{(i) A strong sum is not a valuation limit.} With $\Gamma=\Q+\Q\omega$, the strong sum of
$t^n$ over $n\ge0$ is $(1-t)^{-1}$, but the tails never reach valuation $\omega$. No argument of
the form ``the valuation improves by a fixed positive amount at each iteration'' justifies any
construction in this article.

\emph{(ii) ``Finitely many dependencies at a given exponent'' is not ``finitely many exponents
below it''.} The latter already fails below $\omega$ in $\Q+\Q\omega$. The distinction is exactly
what makes finite-parameter replacement in the proofs of
\cref{trigonometry:thm:conditioned,trigonometry:thm:cluster} legitimate.

\emph{(iii) Well-ordering cannot be weakened to positivity}: the set $\{1/n:n\ge1\}$ is positive
and not well ordered. The lemma applies even when the value group has arbitrary rank: with
$0<\lambda\ll\mu$ one may have $n\lambda<\mu$ for every ordinary $n$ while
$\sum_{n\ge0}t^{n\lambda}$ remains a perfectly valid Hahn series whose partial sums converge in no
relevant sense.
\end{remark}

\begin{proposition}[Degeneracy of fine convergence]\label{trigonometry:prop:topology}
In the fine topology --- generated by all balls $|z-a|<\rho$ with $\rho$ a positive surreal ---
every set-sized subset of $\SC$ is closed and discrete. A net with a set-sized directed index set
which converges in that topology is eventually equal to its limit. In particular a nonconstant
ordinary real-parameter path cannot be fine-continuous on a connected real interval.
\end{proposition}
\begin{proof}
For a set $A\subseteq\SC$ and a point $p$, the cut $\{0\mid\{|p-a|:a\in A,\ a\ne p\}\}$ is a
positive surreal smaller than every nonzero distance from $p$ to $A$; its ball meets $A$ in at most
$\{p\}$. This proves closedness and discreteness. For a convergent net, a positive lower bound for
all its nonzero distances from the alleged limit shows that eventually entering the corresponding
ball forces eventual equality. The image of an ordinary real interval is a set, hence discrete; a
continuous image of a connected space is connected, hence a singleton.
\end{proof}
Thus sequence limits are rejected as a definition device once and for all. For example the strong
series $\sum_{n\ge0}\omega^{-n}$ exists, while its finite partial sums do not converge to it. This
proposition is used again in \cref{trigonometry:prop:perimeters} and in
\cref{trigonometry:sec:arcs}, where it forces the arc-length functional to be \emph{defined}
coefficientwise rather than obtained as a limit of inscribed polygons.

\subsection{Lifting ordinary analytic germs}
For an ordinary real-analytic $f$ near $r\in\R$ define
\begin{equation}\label{trigonometry:eq:lift}
 f\sh(r+\eps)=\sum_{n\ge0}\frac{f^{(n)}(r)}{n!}\eps^n,\qquad\eps\in\mm_{\R},
\end{equation}
and identically for ordinary holomorphic germs and complex infinitesimals. The ordinary centre is
fixed by the standard part, so the definition is unambiguous: \emph{the ordinary value $f(r)$ is
taken first in ordinary analysis, and only the infinitesimal correction is summed.} This is the
restricted-analytic evaluation of \cite{SourceAnalysis,vdDE}. Formal identities among germs lift
to exact identities, because both sides have identical Taylor coefficients and
\cref{trigonometry:lem:support} permits the substitutions.

\begin{proposition}[Local calculus and one-variable sign lifting]\label{trigonometry:prop:lift}
Taylor lifting preserves sums, products, composition where defined, and differentiation, and it is
the unique extension specified by the Taylor rule \cref{trigonometry:eq:lift}. If an ordinary
analytic $g$ is nonnegative on an ordinary closed interval $[a,b]$ and analytic on a neighbourhood
of it, then $g\sh(x)\ge0$ for every surreal $x\in[a,b]$, including points infinitesimally close to
an endpoint.
\end{proposition}
\begin{proof}
The algebraic and compositional statements are formal Taylor identities justified by
\cref{trigonometry:lem:support}. Re-expansion at $x$ gives
$f\sh(x+h)=f\sh(x)+(f')\sh(x)h+h^2R_x(h)$ for infinitesimal $h$, with $R_x(h)$ finite on a
sufficiently small neighbourhood, either from the convergent ordinary germ before substitution or
directly from its Taylor coefficients and their standard parts. The difference quotient therefore
tends to $(f')\sh(x)$ in the fine $\eps$--$\delta$ sense: for any positive surreal tolerance,
choose $h$ smaller than both an infinitesimal radius and that tolerance divided by a fixed real
bound for $R_x$.

For the sign assertion write $x=r+\eps$ with $r\in[a,b]$. If $g(r)>0$ the standard part decides.
If $g(r)=0$ and the germ is not zero, factor it as $(X-r)^mu(X)$ with $u(r)\ne0$; at an interior
zero nonnegativity forces even $m$ and $u(r)>0$, and at an endpoint the leading term has the
correct sign on the side allowed by $x\in[a,b]$. Substitution preserves these signs; a zero germ
causes no difficulty.
\end{proof}
The uniqueness assertion is deliberately about the \emph{Taylor rule}, not about arbitrary
functions satisfying the same differential equation in the fine topology. Moreover, \emph{no
general transfer of arbitrary infinite families of inequalities is asserted}: in particular,
coefficients that are themselves infinitesimal are not automatically covered by an ordinary
real-coefficient sign argument. \Cref{trigonometry:ex:invisible} makes that distinction concrete.

\begin{lemma}[Leading-term test for a lifted germ]\label{trigonometry:lem:leading}
Suppose $f$ has an ordinary zero of order $m$ at $c$, with $a=f^{(m)}(c)/m!\ne0$. Then for every
nonzero infinitesimal $\eta$,
\[
 f\sh(c+\eta)=a\eta^m(1+\xi),\qquad\xi\in\mm_{\SC}.
\]
Hence $v(f\sh(c+\eta))=m\,v(\eta)$ and $\lc(f\sh(c+\eta))=\lc(a\eta^m)$.
\end{lemma}
\begin{proof}
Factor the ordinary germ as $f(c+X)=aX^m(1+Xg(X))$; the evaluated series $\eta g\sh(\eta)$ has
strictly positive support by \cref{trigonometry:lem:support}.
\end{proof}
The two mechanisms used for every inequality below are exactly these: a strict ordinary inequality
with a nonzero ordinary margin survives by taking standard parts, and an ordinary equality case is
resolved by the first nonzero Taylor coefficient.

Finally, a disambiguation that is needed in every later section: \emph{a derivative in this article
is the derivative of a function of a surreal or surcomplex variable}, with limits understood using
all positive surreal error tolerances. It is not a derivation on $\No$ differentiating the
monomials $t^\gamma$; the Berarducci--Mantova differential-field structure \cite{BM} is important
background but is not an input to any argument here.

\section{Canonical sine and cosine on finite angles}\label{trigonometry:sec:finite}

\subsection{Definition without a divergent global series}
For $\theta=r+\eps\in\OO_{\R}$ with $r\in\R$ and $\eps\in\mm_{\R}$ define
\begin{align}
 \sin\theta&=\sin r\sum_{n\ge0}\frac{(-1)^n\eps^{2n}}{(2n)!}
 +\cos r\sum_{n\ge0}\frac{(-1)^n\eps^{2n+1}}{(2n+1)!},\label{trigonometry:eq:sinfinite}\\
 \cos\theta&=\cos r\sum_{n\ge0}\frac{(-1)^n\eps^{2n}}{(2n)!}
 -\sin r\sum_{n\ge0}\frac{(-1)^n\eps^{2n+1}}{(2n+1)!}.\label{trigonometry:eq:cosfinite}
\end{align}
The ordinary trigonometric functions on the right are evaluated \emph{only} at the ordinary real
number $r$. These are exactly the scalar lifts \cref{trigonometry:eq:lift}, and
\begin{equation}\label{trigonometry:eq:finiteEuler}
 \cis\theta=\cos\theta+i\sin\theta=e^{ir}\sum_{n\ge0}\frac{(i\eps)^n}{n!}.
\end{equation}
Unless an extension is explicitly named, lowercase sine, cosine and tangent below have finite
surreal real arguments.

\begin{theorem}[Elementary trigonometry on finite surreal angles]\label{trigonometry:thm:identities}
For $\alpha,\beta\in\OO_{\R}$,
\begin{align}
 \cos^2\alpha+\sin^2\alpha&=1,\\
 \sin(\alpha+\beta)&=\sin\alpha\cos\beta+\cos\alpha\sin\beta,\\
 \cos(\alpha+\beta)&=\cos\alpha\cos\beta-\sin\alpha\sin\beta,\\
 \cis(\alpha+\beta)&=\cis\alpha\,\cis\beta.\label{trigonometry:eq:addition}
\end{align}
Sine is odd, cosine even, both have period $2\pi$, and $(\sin)'=\cos$, $(\cos)'=-\sin$,
$(\cis)'=i\cis$. For every ordinary integer $n$,
\begin{equation}\label{trigonometry:eq:demoivre}
 (\cis\alpha)^n=\cis(n\alpha).
\end{equation}
In particular the usual double-angle, half-angle, sum-to-product and product-to-sum identities
hold, with every quotient restricted to a nonzero denominator.
\end{theorem}
\begin{proof}
Write $\alpha=r+\eps$, $\beta=s+\eta$. The ordinary addition identities at $r,s$ and the formal
identities for the power series in $\eps,\eta$ give \cref{trigonometry:eq:addition}; all
reorganizations are justified by \cref{trigonometry:lem:support}. Conjugating
\cref{trigonometry:eq:finiteEuler} gives $\cis(-\alpha)=\overline{\cis\alpha}$, so the modulus is
one. The remaining identities follow by substitution, conjugation, differentiation and ordinary
finite induction; negative integer powers use $\overline{\cis\alpha}=(\cis\alpha)^{-1}$.
\end{proof}

For reference, several consequences are
\begin{align}
 \tan(\alpha+\beta)&=\frac{\tan\alpha+\tan\beta}{1-\tan\alpha\tan\beta},&
 \sin\alpha+\sin\beta&=2\sin\tfrac{\alpha+\beta}2\cos\tfrac{\alpha-\beta}2,\\
 \cos\alpha-\cos\beta&=-2\sin\tfrac{\alpha+\beta}2\sin\tfrac{\alpha-\beta}2,&
 2\sin\alpha\sin\beta&=\cos(\alpha-\beta)-\cos(\alpha+\beta),\\
 \cos(n\alpha)&=T_n(\cos\alpha),&
 \sin(n\alpha)&=\sin\alpha\,U_{n-1}(\cos\alpha)\quad(n\ge1),\label{trigonometry:eq:chebyshev}
\end{align}
with $T_0=1$, $T_1=X$, $T_{n+1}=2XT_n-T_{n-1}$ and $U_0=1$, $U_1=2X$, $U_{n+1}=2XU_n-U_{n-1}$.
The last two follow from the same recurrence, so no new summability issue enters; the index $n$ is
an ordinary finite integer. The conventional forms are recorded in \cite[\S4.21]{DLMF}; the point
here is their well-defined evaluation and exact validity in a non-Archimedean field.

\subsection{Order, monotonicity, and exact leading orders}
\begin{proposition}[Signs and monotonicity]\label{trigonometry:prop:order}
Sine is positive on $(0,\pi)$ and negative on $(\pi,2\pi)$; cosine is strictly decreasing on
$[0,\pi]$, and sine is strictly increasing on $[-\pi/2,\pi/2]$. Jordan's inequality
\begin{equation}\label{trigonometry:eq:jordan}
 \frac{2x}\pi\le\sin x\le x\qquad(0\le x\le\pi/2)
\end{equation}
holds, and $|\sin x|\le|x|$ for every finite $x$. Moreover $0<\sin x<x<\tan x$ for finite
$0<x<\pi/2$.
\end{proposition}
\begin{proof}
An ordinary standard part strictly between $0$ and $\pi$ decides the positive sign; near $0$,
$\sin\eps=\eps(1+\text{infinitesimal})$, and near $\pi$, $\sin(\pi-\eps)=\sin\eps$. These cover
the remaining points of $(0,\pi)$, and translation by $\pi$ gives the other sign. For
$0\le a<b\le\pi$ the identity $\cos a-\cos b=2\sin\frac{a+b}2\sin\frac{b-a}2>0$ proves strict
decrease, and the analogous sine difference uses a positive cosine of the midpoint on
$[-\pi/2,\pi/2]$. The ordinary inequalities in \cref{trigonometry:eq:jordan} lift by
\cref{trigonometry:prop:lift}; for $|x|\ge1$ use $|\sin x|\le1$. For the last chain, standard parts
in $(0,\pi/2)$ have ordinary margins; at positive infinitesimal $x$ the leading terms
$x-\sin x=x^3/6+\cdots$ and $\tan x-x=x^3/3+\cdots$ are positive by
\cref{trigonometry:lem:leading}; and infinitesimally below $\pi/2$ the tangent is positive
infinite.
\end{proof}
This proof does not infer global monotonicity from the sign of a fine derivative. Such an inference
would require additional hypotheses on a general surreal function.

\begin{proposition}[Exact infinitesimal leading orders]\label{trigonometry:prop:leading}
For $0\ne h\in\mm_{\R}$,
\begin{align}
 \sin h&=h-\frac{h^3}6+\frac{h^5}{120}+O(h^7),&
 1-\cos h&=\frac{h^2}2-\frac{h^4}{24}+O(h^6),\\
 \tan h&=h+\frac{h^3}3+\frac{2h^5}{15}+O(h^7),
\end{align}
and, exactly,
\begin{equation}\label{trigonometry:eq:sinval}
 v(\sin h)=v(\tan h)=v(h),\qquad v(1-\cos h)=2v(h),
\end{equation}
with $\sin h\sim h$, $\tan h\sim h$ and $1-\cos h\sim h^2/2$.
\end{proposition}
\begin{proof}
The truncations come from the strongly summable formal series. Dividing the three expressions by
$h,h^2,h$ gives finite elements with nonzero standard parts $1,1/2,1$, and such units have
valuation zero.
\end{proof}
These asymptotic equivalences are statements about \emph{a single arbitrary infinitesimal}. They
are strictly stronger than standard-part equalities and do not mean that a sequence of real
arguments converges in the fine topology.

\subsection{What normalizes the finite pair}
\begin{proposition}[Normalization of the finite trigonometric pair]
\label{trigonometry:prop:normalization}
Suppose two functions $S,C$ on $\OO_{\R}$ agree with the ordinary sine and cosine at every
ordinary real point, and satisfy, for every $r\in\R$ and every $\eps\in\mm_{\R}$, the Taylor
identities \cref{trigonometry:eq:sinfinite}--\cref{trigonometry:eq:cosfinite} with the ordinary
derivatives at $r$. Then $S=\sin$ and $C=\cos$ on $\OO_{\R}$. Equivalently, the prescribed ordinary
germs together with their strong evaluations determine the pair uniquely.
\end{proposition}
\begin{proof}
Every $x\in\OO_{\R}$ has exactly one decomposition $r+\eps$, so the hypothesis
reproduces \cref{trigonometry:eq:sinfinite} and \cref{trigonometry:eq:cosfinite}
verbatim.
\end{proof}
This normalization is deliberately stronger than a differential equation with one initial value.
Fine-local analyticity alone permits unrelated germs on disjoint monads \cite{SourceAnalysis}, and
\cref{trigonometry:thm:characters} will show that even a global addition law does not determine
the infinite-argument values.

\section{The full unit circle and the group of angles}\label{trigonometry:sec:circle}

\subsection{Infinitesimal exponential and logarithm}
For $h,\eta\in\mm_{\C}$ put
\begin{equation}\label{trigonometry:eq:explog}
 \Ex(h)=\sum_{n\ge0}\frac{h^n}{n!},\qquad
 \Logm(1+\eta)=\sum_{n\ge1}\frac{(-1)^{n+1}\eta^n}{n}.
\end{equation}
Formal power-series identities and \cref{trigonometry:lem:support} prove that these are mutually
inverse bijections between $\mm_{\C}$ and $1+\mm_{\C}$, that $\Ex(h+k)=\Ex(h)\Ex(k)$, and that the
logarithm turns products into sums; both commute with conjugation. In particular
\begin{equation}\label{trigonometry:eq:logleading}
 v(\Logm(1+\eta))=v(\eta)\quad(\eta\ne0),\qquad
 \cis h=\Ex(ih)\quad(h\in\mm_{\R}).
\end{equation}
Define
\[
 \TT(\No)=\{u\in\SC:u\bar u=1\},
\]
whose elements all have finite coordinates even though $\TT(\No)$ is a proper class, and let
$\TT(\R)$ be the ordinary complex unit circle. For nonzero $z,w$ the \emph{oriented angle} from
$z$ to $w$ is, before any radian parameter, the group element
\begin{equation}\label{trigonometry:eq:angleunit}
 u(z,w)=\frac{\bar zw}{|z||w|}\in\TT(\No),
\end{equation}
whose real and imaginary coordinates are its cosine and sine:
$\cos(z,w)=\langle z,w\rangle/(|z||w|)$ and $\sin(z,w)=[z,w]/(|z||w|)$. Changing either vector's
length by a positive surreal factor does not change its angle. The \emph{unoriented} angle is the
unique $\alpha\in[0,\pi]$ with $\cos\alpha=\langle z,w\rangle/(|z||w|)$ and
$\sin\alpha=|[z,w]|/(|z||w|)$; by \cref{trigonometry:eq:gram} this is meaningful even when one
vector has infinitesimal length and the other infinite length.

\begin{proposition}[Algebraic angle addition and rotations]\label{trigonometry:prop:rotation}
Angle composition is multiplication: $u(z,w)u(w,q)=u(z,q)$, and if the angle coordinates are
$(c,s)$ and $(d,t)$ the composite has coordinates $(cd-st,\ sd+ct)$. The map
$u=c+is\mapsto\left(\begin{smallmatrix}c&-s\\ s&c\end{smallmatrix}\right)$ identifies $\TT(\No)$
with $\SO(2,\No)$. Multiplication by an element of $\TT(\No)$ preserves lengths, dot products,
determinants and oriented angles; every orientation-preserving $\No$-linear isometry of $\No^2$ is
of this form; every orientation-preserving affine similarity is $z\mapsto a+\lambda uz$ with
$a\in\SC$, $\lambda\in\No_{>0}$, $u\in\TT(\No)$, and the orientation-reversing ones are obtained by
composing with conjugation. Finally every nonzero $z\in\SC$ is uniquely $z=\rho u$ with
$\rho=|z|\in\No_{>0}$ and $u\in\TT(\No)$, and multiplication gives a group isomorphism
$\No_{>0}\times\TT(\No)\to\SC^\times$.
\end{proposition}
\begin{proof}
The angle identity follows by cancellation, and multiplying $c+is$ by $d+it$ gives the displayed
coordinates. The indicated matrix preserves the dot product and has determinant $c^2+s^2=1$;
conversely the first column of an orthogonal determinant-one matrix is some $(c,s)$ of norm one,
and orthogonality fixes the second column as $(-s,c)$. For a unit $q$ one has
$\overline{qz}\,qw=\bar zw$, which gives the preservation statements. Dividing a similarity's
linear part by its positive scale factor reduces it to the isometry case. For the polar statement
take $\rho=|z|$ and $u=z/|z|$; taking norms proves uniqueness and multiplicativity.
\end{proof}
Thus addition, double-angle and finite multiple-angle identities already have an algebraic meaning
over any real closed field. The linearity hypothesis on isometries can be dropped
(\cref{trigonometry:rot:thm:isometry}, from source 19, which also proves the matrix identification). Surreal-specific analysis enters only when these unit directions are
identified with numerical radians, which is the content of the next theorem.

\subsection{The finite phase theorem}
\begin{theorem}[Polar decomposition and the angle group]\label{trigonometry:thm:polar}
The map $\cis:(\OO_{\R},+)\to(\TT(\No),\cdot)$ is surjective with kernel exactly $2\pi\Z$, so that
\begin{equation}\label{trigonometry:eq:anglegroup}
 \OO_{\R}/2\pi\Z\ \cong\ \TT(\No),
 \qquad\text{and}\qquad
 \SC^\times\ \cong\ \No_{>0}\times\bigl(\OO_{\R}/2\pi\Z\bigr).
\end{equation}
More precisely, there is a canonical group isomorphism
\begin{equation}\label{trigonometry:eq:circlesplit}
 \TT(\No)\longrightarrow\TT(\R)\times(\mm_{\R},+),\qquad
 u\longmapsto\bigl(u_0,\ -i\Logm(u/u_0)\bigr),\quad u_0=\st(u),
\end{equation}
with inverse $(u_0,\delta)\mapsto u_0\Ex(i\delta)$. Consequently every $z\ne0$ has a unique
modulus $\rho=|z|>0$ and a unique angle class $[\theta]\in\OO_{\R}/2\pi\Z$ with $z=\rho\cis\theta$;
the modulus may be infinitesimal, finite or infinite, while the angle is always finite.
\end{theorem}
\begin{proof}
Taking standard parts in $u\bar u=1$ gives $u_0\in\TT(\R)$. Put $q=u/u_0\in1+\mm_{\C}$; since
$\bar q=q^{-1}$,
\[
 \overline{\Logm q}=\Logm\bar q=\Logm(q^{-1})=-\Logm q,
\]
so $\Logm q=i\delta$ for a unique $\delta\in\mm_{\R}$. Choosing any ordinary argument $\theta_0$
of $u_0$ gives $u=u_0\Ex(i\delta)=\cis(\theta_0+\delta)$, proving surjectivity. The standard-part
map and the logarithmic group law prove that \cref{trigonometry:eq:circlesplit} is a homomorphism
with the stated inverse. If $\cis(r+\eps)=1$, the standard part gives $r\in2\pi\Z$, hence
$\Ex(i\eps)=1$, and logarithmic inversion gives $\eps=0$. Applying the circle result to $u=z/|z|$
and \cref{trigonometry:prop:rotation} gives the polar statement.
\end{proof}

This theorem is also the structure theorem of source 19 for $\SO(2,\No)$, proved there by the same
argument; \cref{trigonometry:rot:sec:main} develops its group-theoretic consequences.
The infinitesimal component in \cref{trigonometry:eq:circlesplit} is independent of an ordinary
argument branch. A numerical representative of an angle class, however, requires a cut.

\begin{corollary}[Representatives, roots, and torsion]\label{trigonometry:cor:representatives}
Every direction has exactly one representative in $[0,2\pi)$ and exactly one in $(-\pi,\pi]$,
obtained by adding or subtracting an ordinary multiple of $2\pi$ and then checking the
\emph{actual} surreal inequalities; near the endpoints the infinitesimal correction matters, so
that the representative of $-\eps$ with $\eps>0$ infinitesimal is $2\pi-\eps$ and not an ordinary
angle. The upper open semicircle corresponds exactly to $(0,\pi)$. For $u=\cis\theta$ and an
ordinary $n\ge1$, the equation $v^n=u$ has exactly the $n$ distinct solutions
$\cis((\theta+2\pi k)/n)$, $0\le k<n$; in particular the $n$th roots of unity in $\SC$ are exactly
the ordinary complex roots of unity, and $\TT(\No)$ has no nontrivial infinitesimal torsion. For
$w=\rho\cis\theta\ne0$ the $n$th roots are $\rho^{1/n}\cis((\theta+2\pi k)/n)$, and for $n\ge3$ a
regular $n$-gon of radius $\rho\in\No_{>0}$ has side length $2\rho\sin(\pi/n)$.
\end{corollary}
\begin{proof}
A finite number has an ordinary real standard part, so subtracting an ordinary multiple of $2\pi$,
with at most one endpoint adjustment, lands in either half-open interval; the kernel in
\cref{trigonometry:thm:polar} gives uniqueness. On $(0,\pi)$ sine is positive by
\cref{trigonometry:prop:order}, including the infinitesimal endpoint neighbourhoods, and parity
treats the lower semicircle. The listed roots are distinct by the kernel and exhaust the roots of
$X^n-u$ by degree. In the splitting \cref{trigonometry:eq:circlesplit}, $n\delta=0$ forces
$\delta=0$. The chord formula follows from \cref{trigonometry:eq:gram} or the cosine double-angle
identity.
\end{proof}
Source 19 proves the root count by a second route through the splitting alone
(\cref{trigonometry:rot:rem:rootsroute}), and adds the classification of finite subgroups
(\cref{trigonometry:rot:cor:torsion}).

\begin{remark}[Two branch conventions, kept apart]\label{trigonometry:rem:cut}
A representative constrained to have \emph{actual value} in $(-\pi,\pi]$ has a branch jump at the
negative real axis, including across infinitesimal perturbations: an ordinary principal argument of
$\st(u)$ plus its infinitesimal correction need not already lie in that interval --- it may be
$\pi+\eps$ with $\eps>0$ --- and the representative then subtracts $2\pi$. One may instead fix an
argument $\beta$ of the \emph{standard} direction and retain its infinitesimal correction without
reducing again,
\begin{equation}\label{trigonometry:eq:unwrap}
 \arg_\beta z=\beta-i\Logm\!\left(\frac{z/|z|}{\st(z/|z|)}\right),
\end{equation}
which is fine-locally analytic but may lie infinitesimally outside $(-\pi,\pi]$. The global
logarithm of \cite{SourceAnalysis} uses the second, standard-part-pinned convention; the geometric
representative convention of \cref{trigonometry:cor:representatives} is the first. They give the
same direction and must not be identified as the same single-valued argument function.
\emph{Angle classes avoid the ambiguity entirely}, and are what all geometric theorems below use.
When $h/z$ is infinitesimal the unwrapped branch gives the local differential
\begin{equation}\label{trigonometry:eq:argdiff}
 \arg_\beta(z+h)-\arg_\beta z=\im\Logm(1+h/z),
\end{equation}
so that $d\arg z=\im(dz/z)$ and $d\log|z|=\re(dz/z)$.
\end{remark}

\subsection{Rational coordinates and angle division}
\begin{theorem}[Projective half-angle parametrization of every surreal direction]
\label{trigonometry:thm:cayley}
The map
\begin{equation}\label{trigonometry:eq:cayley}
 C(t)=\frac{1+it}{1-it}=\frac{1-t^2}{1+t^2}+i\frac{2t}{1+t^2},\qquad t\in\No,
 \qquad C(\infty)=-1,
\end{equation}
is a bijection $\PP^1(\No)\to\TT(\No)$; here $\infty$ is the projective point, not an infinitely
large surreal. Restricted to $\No$ it is a bijection onto $\TT(\No)\setminus\{-1\}$, with inverse
$t=\im u/(1+\re u)$. If $u=\cis\theta\ne-1$ then $t=\tan(\theta/2)$ for the representative
$-\pi<\theta<\pi$. In homogeneous coordinates $t=[p:q]$, multiplication of directions is
\begin{equation}\label{trigonometry:eq:projectiveaddition}
 [p:q]\oplus[r:s]=[ps+qr:\ qs-pr],
\end{equation}
which in affine coordinates is $(t+u)/(1-tu)$ whenever the denominator is nonzero. Infinite $t$
are allowed and describe directions infinitesimally close to $-1$.
\end{theorem}
\begin{proof}
The displayed rational expression has norm one and cannot equal $-1$, because
$1+\re C(t)=2/(1+t^2)>0$. Conversely if $u=c+is$ has $c^2+s^2=1$ and $c\ne-1$ then $1+c>0$ and
substituting $t=s/(1+c)$ recovers $c+is$; the only omitted direction is $-1$, supplied by the
projective point. The product formula follows by multiplying $(q+ip)/(q-ip)$ and $(s+ir)/(s-ir)$,
and its two homogeneous outputs cannot both vanish because
$(ps+qr)^2+(qs-pr)^2=(p^2+q^2)(r^2+s^2)>0$. The half-angle identity is
\cref{trigonometry:thm:identities} on the chosen interval; the exceptional angle is handled
projectively.
\end{proof}
The homogeneous rule resolves the cases involving $t=\infty$: in particular
$\infty\oplus\infty=0$, corresponding to two half-turns, and one must not evaluate the affine
formula by informal arithmetic with a symbol $\infty$. This rational coordinate is valid for
infinite $t$: $t=\omega$ gives a direction infinitesimally close to $-1$, not an unspecified
direction associated with $\omega$ radians. Algebraic half-slope coordinates have published
predecessors in vector and rational trigonometry \cite{Wildberger,Gunn}; their use here avoids
unnecessary transcendental evaluation, and \cref{trigonometry:thm:fejer} and
\cref{trigonometry:thm:stability} both run through this chart.

\section{Inverse functions and the angular metric}\label{trigonometry:sec:inverse}

\subsection{Inverse trigonometry on its full natural surreal domains}
\begin{theorem}[Inverse tangent, sine, and cosine]\label{trigonometry:thm:inverse}
There are unique inverse functions
\[
 \arctan:\No\to(-\pi/2,\pi/2),\qquad
 \arcsin:[-1,1]\to[-\pi/2,\pi/2],\qquad
 \arccos:[-1,1]\to[0,\pi],
\]
with the usual inverse identities; \emph{all values are finite surreal angles}, and inverse tangent
is an increasing bijection defined at every surreal slope, infinite ones included. They satisfy
\begin{align}
 \cos(\arctan t)&=\frac1{\sqrt{1+t^2}},&
 \sin(\arctan t)&=\frac t{\sqrt{1+t^2}},\\
 \arcsin x&=2\arctan\frac x{1+\sqrt{1-x^2}},&
 \arccos t&=\frac\pi2-\arcsin t,
\end{align}
and at every point of their open domains --- including infinite $t$ for inverse tangent and points
infinitesimally close to $\pm1$ for the other two ---
\begin{equation}\label{trigonometry:eq:inversederivatives}
 (\arctan)'(t)=\frac1{1+t^2},\quad
 (\arcsin)'(t)=\frac1{\sqrt{1-t^2}},\quad
 (\arccos)'(t)=-\frac1{\sqrt{1-t^2}} .
\end{equation}
\end{theorem}
\begin{proof}
The unit direction $(1+it)/\sqrt{1+t^2}$ has positive real part, so
\cref{trigonometry:thm:polar} and the sign rules of \cref{trigonometry:prop:order} give a unique
angle in $(-\pi/2,\pi/2)$; this defines inverse tangent and proves its two coordinate formulas. If
$x<y$ the determinant of $(1,x)$ and $(1,y)$ is $y-x>0$, so the angle difference has positive sine
and, both angles lying in an interval of length $\pi$, is positive: inverse tangent is increasing.
Similarly the directions $\sqrt{1-t^2}+it$ and $t+i\sqrt{1-t^2}$ give inverse sine and inverse
cosine on their stated ranges, whose uniqueness follows from strict monotonicity. The half-angle
identity proves the inverse-sine formula, whose denominator is positive even at $x=\pm1$.

For the derivatives: locally, square roots of nonzero positive elements have their binomial
expansions after division by the base value, and a local argument difference is the imaginary part
of the infinitesimal logarithm of a direction ratio (\cref{trigonometry:eq:argdiff}). These give
fine-local analytic expansions at every point under discussion, and differentiating the inverse
identities gives \cref{trigonometry:eq:inversederivatives}; the relevant cosine or sine is nonzero.
At an infinite inverse-tangent input the same computation is valid after taking increments small
relative to $1+it$, and no finite bound on $t$ is used.
\end{proof}

For positive infinite $t$, hence infinitesimal $1/t$,
\begin{equation}\label{trigonometry:eq:ataninfty}
 \arctan t=\frac\pi2-\arctan(1/t)
 =\frac\pi2-\frac1t+\frac1{3t^3}-\frac1{5t^5}+O(t^{-7}),
\end{equation}
and for negative $t$ one uses oddness; an infinite slope corresponds to a finite angle
infinitesimally short of a right angle. For example
\[
 \Arg(1+i\omega)=\frac\pi2-\omega^{-1}+\frac{\omega^{-3}}3-\frac{\omega^{-5}}5+\cdots,
\]
a finite value describing an infinite slope without evaluating sine or cosine at an infinite
angular argument. For infinitesimal $h$ the strong identities
\begin{equation}\label{trigonometry:eq:smallseries}
 \arctan h=h-\frac{h^3}3+\frac{h^5}5-\cdots,\qquad
 \arcsin h=h+\frac{h^3}6+\frac{3h^5}{40}+\cdots
\end{equation}
hold, with $(\tan x)'=1+\tan^2x$ wherever cosine is nonzero.

\begin{remark}[An endpoint at which no surreal derivative exists]
\label{trigonometry:rem:noendpointderivative}
The endpoint derivative of inverse sine does not exist as a surreal-valued derivative: its
difference quotient grows beyond every fixed surreal bound as the increment becomes sufficiently
small. Allowing infinitely large elements as candidate derivative values does not repair this,
because there is always a still smaller increment. This is a genuine failure, not an artefact of
the choice of tolerance class, and it is the analytic shadow of the square-root ramification of
\cref{trigonometry:prop:acosendpoint}.
\end{remark}

\begin{proposition}[Endpoint ramification of inverse cosine]\label{trigonometry:prop:acosendpoint}
For $0<\tau\in\mm_{\R}$,
\begin{align}
 \arccos(1-\tau)&=2\arcsin\sqrt{\tau/2}\label{trigonometry:eq:acosseries}\\
 &=\sqrt{2\tau}\sum_{n\ge0}\frac{\binom{2n}n}{8^n(2n+1)}\tau^n\notag\\
 &=\sqrt{2\tau}\left(1+\frac\tau{12}+\frac{3\tau^2}{160}+\frac{5\tau^3}{896}+O(\tau^4)\right),
 \notag
\end{align}
the series being strongly summable after adjoining $\sqrt\tau$, and
\begin{equation}\label{trigonometry:eq:acosval}
 v\bigl(\arccos(1-\tau)\bigr)=\tfrac12v(\tau).
\end{equation}
\end{proposition}
\begin{proof}
The double-angle identity gives $\cos(2\arcsin\sqrt{\tau/2})=1-\tau$; the angle is a positive
infinitesimal, so inverse uniqueness gives the first equality. Substituting the ordinary formal
inverse-sine series gives the rest, every remaining parenthesized factor having standard part one.
\end{proof}
Inverse sine and inverse tangent are valuation-preserving at zero, whereas inverse cosine at $1$ is
a square-root branch. This single distinction governs every degeneration result in
\cref{trigonometry:sec:valuationgeometry} and \cref{trigonometry:sec:contact}.

\subsection{Chord distance and infinitesimal angular accuracy}
For $u,v\in\TT(\No)$ let $\ang uv\in[0,\pi]$ be the least absolute representative of their angle
difference, equivalently $\ang uv=\arccos\re(\bar uv)$. This is a metric: choose shortest
representatives for two successive differences, add them and reduce modulo $2\pi$; reduction cannot
increase the least absolute value, which is the triangle inequality.

\begin{theorem}[Angular and chordal comparison]\label{trigonometry:thm:metric}
For $u,v\in\TT(\No)$, with $d=\ang uv$,
\begin{equation}\label{trigonometry:eq:metricbounds}
 |u-v|=2\sin(d/2),\qquad |u-v|\le d\le\frac\pi2|u-v|,
\end{equation}
and for $u\ne v$,
\begin{equation}\label{trigonometry:eq:metricval}
 v\bigl(\ang uv\bigr)=v(u-v)\quad\text{exactly.}
\end{equation}
If $\st(u)=\st(v)$, their unique infinitesimal oriented angle difference is
\begin{equation}\label{trigonometry:eq:localangle}
 \delta=-i\Logm(v/u),\qquad v(\delta)=v(v-u),\qquad \ang uv=|\delta| .
\end{equation}
Moreover, for infinitesimal $d$,
\begin{equation}\label{trigonometry:eq:chordseries}
 |u-v|=d\left(1-\frac{d^2}{24}+\frac{d^4}{1920}-\cdots\right),\qquad
 1-\cos d=\frac{d^2}2\left(1-\frac{d^2}{12}+\frac{d^4}{360}-\cdots\right).
\end{equation}
\end{theorem}
\begin{proof}
Expanding $|u-v|^2=2-2\re(\bar uv)$ and using the half-angle identity proves the first equality in
\cref{trigonometry:eq:metricbounds}; Jordan's inequality \cref{trigonometry:eq:jordan} proves both
bounds. Bounding a positive quantity above and below by nonzero ordinary real multiples of another
forces \cref{trigonometry:eq:metricval}. Finally $v/u\in1+\mm_{\C}$, and
\cref{trigonometry:eq:logleading} gives \cref{trigonometry:eq:localangle}. The series come from
substitution into the sine and cosine Taylor series, the parenthesized factors having standard part
one.
\end{proof}
Thus \emph{angular precision is exactly the precision of normalized directions}, not merely
comparable to it up to an unspecified non-Archimedean error. Two consequences deserve to be
isolated.

\begin{corollary}[Local valuation isometry of phase and rotation displacement]
\label{trigonometry:cor:phaseisometry}
If $\theta,\phi$ are finite real angles with $0\ne\theta-\phi\in\mm_{\R}$ then
\begin{equation}\label{trigonometry:eq:phaseisometry}
 |\cis\theta-\cis\phi|\sim|\theta-\phi|,\qquad
 v(\cis\theta-\cis\phi)=v(\theta-\phi),\qquad
 \frac{\cis\theta-\cis\phi}{i\cis\phi\,(\theta-\phi)}=1+\text{an infinitesimal}.
\end{equation}
For $z\ne0$ and a nonzero infinitesimal rotation angle $\delta$,
\begin{equation}\label{trigonometry:eq:rotdisplacement}
 |\cis\delta\,z-z|=2|z||\sin(\delta/2)|\sim|z||\delta| .
\end{equation}
\end{corollary}
\begin{proof}
Factor $\cis\theta-\cis\phi=\cis\phi(\Ex(i(\theta-\phi))-1)$; the first nonzero term of the last
factor is $i(\theta-\phi)$, with infinitesimal relative remainder. Equation
\cref{trigonometry:eq:rotdisplacement} is \cref{trigonometry:eq:metricbounds} followed by
\cref{trigonometry:prop:leading}.
\end{proof}
A rotation that is infinitesimal \emph{as an angle} need not be infinitesimal \emph{as a
displacement}: rotating a vector of length $\omega$ by $\omega^{-1}$ moves its endpoint by a
distance relatively equivalent to $1$, and replacing $\omega$ by $\omega^2$ makes that displacement
infinite.

\begin{corollary}[Stability of direction under a relative perturbation]
\label{trigonometry:cor:directionstability}
If $z\ne0$ and $h/z=\eta\in\mm_{\C}$, then the local angle difference between $z+h$ and $z$ is
\[
 \Delta\theta=\im\Logm(1+\eta)=\im\eta-\tfrac12\im(\eta^2)+O(|\eta|^3).
\]
In particular $v(\Delta\theta)\ge v(\eta)$ when the difference is nonzero, with equality when the
leading coefficient of $\eta$ has nonzero imaginary part.
\end{corollary}
\begin{proof}
The real part of the logarithm accounts for the change in modulus, its imaginary part for the change
in direction; taking imaginary parts cannot lower a Hahn valuation, and a first coefficient with
nonzero imaginary part cannot cancel against higher powers.
\end{proof}

Although cosine loses first-order information near zero --- $v(1-\cos d)=2v(d)$ --- the full unit
direction does not, because its sine coordinate retains the angular scale. At other angles one
should use both coordinates, or a nonsingular tangent chart, rather than infer angular error solely
from cosine error.

\section{Triangles at arbitrary surreal scales}\label{trigonometry:sec:triangles}

\subsection{Angles and the exact angle sum}
Let $A,B,C\in\SC$ be noncollinear, with side lengths $a=|B-C|$, $b=|C-A|$, $c=|A-B|$ and positive
area $\Delta=\frac12|[B-A,C-A]|$. All four quantities are positive surreals; \emph{no finiteness
assumption is made}. Define the interior angle $\alpha$ at $A$ by
\begin{equation}\label{trigonometry:eq:interiorangle}
 \cos\alpha=\frac{\langle B-A,C-A\rangle}{bc},\qquad
 \sin\alpha=\frac{2\Delta}{bc},\qquad 0<\alpha<\pi,
\end{equation}
and $\beta,\gamma$ cyclically. Equation \cref{trigonometry:eq:gram} places the right-hand sides on
the upper unit semicircle, so \cref{trigonometry:thm:polar} gives existence and uniqueness.

\begin{theorem}[Euclidean angle sum over the surreal field]\label{trigonometry:thm:anglesum}
Every noncollinear surcomplex triangle satisfies $\alpha+\beta+\gamma=\pi$.
\end{theorem}
\begin{proof}
A translation, rotation and possibly a conjugation put $A=0$, $B=c>0$, $C=x+iy$ with $y>0$. Let
$\alpha=\Arg C\in(0,\pi)$ and $\varphi=\Arg(C-c)\in(0,\pi)$. Since
$\im(\bar C(C-c))=cy>0$ we get $\sin(\varphi-\alpha)>0$; the difference lies in $(-\pi,\pi)$, so
$\varphi>\alpha$. The angle at $B$ is $\pi-\varphi$ and the angle at $C$ is $\varphi-\alpha$, both
in $(0,\pi)$, and their sum with $\alpha$ is $\pi$.
\end{proof}
The proof uses order, angle classes and determinants only. It applies equally to an infinitesimal
triangle, an infinite triangle, or a triangle whose sides lie on three different scales, and the
conclusion is an equality of finite surreals including their infinitesimal parts, not merely of
standard parts. An equivalent argument multiplies the three interior-angle unit elements,
\[
 \frac{C-A}{B-A}\cdot\frac{A-B}{C-B}\cdot\frac{B-C}{A-C}=-1,
\]
so that $\cis(\alpha+\beta+\gamma)=-1$ with the sum in $(0,3\pi)$; the kernel of
\cref{trigonometry:thm:polar} then forces $\pi$. Neither argument uses compactness, a topological
hypothesis, or an Archimedean parallel-postulate step.

\subsection{Cosine, sine and tangent laws}
\begin{theorem}[Triangle laws]\label{trigonometry:thm:trianglelaws}
For every noncollinear surcomplex triangle,
\begin{align}
 a^2&=b^2+c^2-2bc\cos\alpha,\\
 \Delta&=\tfrac12bc\sin\alpha=\tfrac12ca\sin\beta=\tfrac12ab\sin\gamma,\\
 \frac a{\sin\alpha}&=\frac b{\sin\beta}=\frac c{\sin\gamma}=\frac{abc}{2\Delta}=2R,\\
 \frac{a-b}{a+b}&=\frac{\tan((\alpha-\beta)/2)}{\tan((\alpha+\beta)/2)},
 \label{trigonometry:eq:trianglelaws}
\end{align}
together with the cyclic versions of the cosine law. Here $R$ is the radius of the unique
circumcircle, and $R=abc/(4\Delta)$. All denominators are nonzero, and the sides may be infinite,
infinitesimal, or of mutually different scales, as may the area and the circumradius.  Those
scales are not independent of one another: \cref{trigonometry:thm:valuationtriangle} relates
them, and in particular the least of $v(a),v(b),v(c)$ is always attained at least twice.
\end{theorem}
\begin{proof}
Expanding $|(C-A)-(B-A)|^2$ gives the cosine law, and \cref{trigonometry:eq:interiorangle} with its
cyclic versions gives the area identity and the common sine-law value $abc/(2\Delta)$. For the
circumcircle set $A=0$, $B=c$, $C=x+iy$ with $y>0$; a centre $O=X+iY$ equidistant from the three
vertices satisfies $2cX=c^2$ and $2xX+2yY=x^2+y^2=b^2$, two independent linear equations with the
unique solution $X=c/2$, $Y=(b^2-cx)/(2y)$. A direct expansion using $a^2=(x-c)^2+y^2$ gives
$|O|^2=a^2b^2/(4y^2)$, so $R=ab/(2y)=abc/(4\Delta)$; existence is thereby proved without a
compactness argument. Finally the sine law converts $(a-b)/(a+b)$ into
$(\sin\alpha-\sin\beta)/(\sin\alpha+\sin\beta)$, and the sum-to-product identities give the tangent
law; since $0<\alpha+\beta<\pi$ and $|\alpha-\beta|<\pi$, every cancellation and tangent is
legitimate.
\end{proof}

\begin{corollary}[Right triangles and the full range of slopes]\label{trigonometry:cor:right}
A triangle is right at $C$ exactly when $c^2=a^2+b^2$. In a right triangle with positive legs
$p,q$, the acute angle opposite $q$ is $\arctan(q/p)$, and $q/p$ may be \emph{any} positive
surreal, including an infinite one.
\end{corollary}
\begin{proof}
The cosine law and the unique zero of cosine on $(0,\pi)$ give the first assertion; the second is
the coordinate formula of \cref{trigonometry:thm:inverse}.
\end{proof}

\subsection{Reconstruction from side or angle data}
\begin{theorem}[Existence, congruence, and similarity]\label{trigonometry:thm:sss}
Positive surreals $a,b,c$ are the side lengths of a noncollinear surcomplex triangle if and only if
all three strict triangle inequalities $a<b+c$, $b<c+a$, $c<a+b$ hold; the triangle is then unique
up to a Euclidean isometry, the two orientations being interchanged by reflection, and its angles
are recovered by the cosine law and $\arccos$. Conversely, positive finite $\alpha,\beta,\gamma$
with $\alpha+\beta+\gamma=\pi$ and an arbitrary positive scale $k\in\No$ determine a triangle with
\[
 a=k\sin\alpha,\qquad b=k\sin\beta,\qquad c=k\sin\gamma,
\]
unique up to similarity; specifying one positive side fixes $k$. Triangles with the same three
angles have proportional corresponding sides.
\end{theorem}
\begin{proof}
Necessity is the norm triangle inequality with equality excluded by noncollinearity. For
sufficiency place $A=0$, $B=c$ and set
\[
 x=\frac{b^2+c^2-a^2}{2c},\qquad y=\sqrt{b^2-x^2} .
\]
The Heron factorization
\begin{equation}\label{trigonometry:eq:heronfactor}
 4c^2(b^2-x^2)=(a+b+c)(-a+b+c)(a-b+c)(a+b-c)
\end{equation}
shows $y>0$, so $C=x+iy$ has the required distances; the same equations force $x$ and $|y|$, giving
uniqueness. The cosine-law quotient lies strictly between $-1$ and $1$, so
\cref{trigonometry:thm:inverse} recovers each interior angle.

For the converse, note that for $\alpha,\beta>0$ with $\alpha+\beta<\pi$,
\[
 \sin\alpha+\sin\beta=2\sin\tfrac{\alpha+\beta}2\cos\tfrac{\alpha-\beta}2
 >2\sin\tfrac{\alpha+\beta}2\cos\tfrac{\alpha+\beta}2=\sin(\alpha+\beta),
\]
the strict comparison following from cosine monotonicity on $[0,\pi]$ and
$|\alpha-\beta|<\alpha+\beta$. Since $\sin\gamma=\sin(\alpha+\beta)$, the proposed sides satisfy
the strict triangle inequalities cyclically, so a triangle exists. The identity
$\sin^2\beta+\sin^2\gamma-\sin^2\alpha=2\sin\beta\sin\gamma\cos\alpha$, which follows from
$\alpha=\pi-\beta-\gamma$ and the addition formulas, shows via the cosine law that the constructed
angle at $A$ is exactly $\alpha$, cosine being injective on $[0,\pi]$; similarly at the other
vertices. Changing $k$ is a similarity, and the sine law gives proportional sides for equal angles.
\end{proof}
This is the usual side--side--side construction without any Archimedean assumption;
side--angle--side is obtained from $C=b\cis\alpha$ with $0<\alpha<\pi$, and the sine law gives
angle--side--angle. \emph{The usual side--side--angle ambiguity remains}: a ray--circle
intersection is a quadratic problem and may have zero, one or two admissible solutions, and no
additional infinite family of solutions appears merely because the coordinates are surreal. This
last point is made quantitative by \cref{trigonometry:thm:amplitude}.

\subsection{Area, half-angles, and the two radii}
Put $s=(a+b+c)/2$; the strict triangle inequalities make $s-a$, $s-b$, $s-c$ positive.

\begin{theorem}[Heron, half-angle, incircle and bisector formulas]\label{trigonometry:thm:heron}
One has
\begin{align}
 \Delta^2&=s(s-a)(s-b)(s-c),& r&=\frac\Delta s,\\
 \sin\frac\alpha2&=\sqrt{\frac{(s-b)(s-c)}{bc}},&
 \cos\frac\alpha2&=\sqrt{\frac{s(s-a)}{bc}},\\
 \tan\frac\alpha2&=\sqrt{\frac{(s-b)(s-c)}{s(s-a)}}=\frac r{s-a},
 \label{trigonometry:eq:halfangletriangle}
\end{align}
all square roots being the positive ones. The incentre is
\begin{equation}\label{trigonometry:eq:incenter}
 I=\frac{aA+bB+cC}{2s},
\end{equation}
and the internal bisector from $A$ meets $BC$ in the ratio $c:b$ and has length
\begin{equation}\label{trigonometry:eq:bisector}
 \ell_A=\frac{2bc}{b+c}\cos\frac\alpha2 .
\end{equation}
\end{theorem}
\begin{proof}
With the coordinates of \cref{trigonometry:thm:sss}, $\Delta=cy/2$, so
\cref{trigonometry:eq:heronfactor} proves Heron's formula. Substituting the cosine law into
$\sin^2(\alpha/2)=(1-\cos\alpha)/2$ and $\cos^2(\alpha/2)=(1+\cos\alpha)/2$ gives the half-angle
squares; both half-angle functions are positive, which fixes the roots. Their quotient together
with Heron's identity gives \cref{trigonometry:eq:halfangletriangle}.

The point \cref{trigonometry:eq:incenter} has positive barycentric weights, hence lies strictly
inside; its distance to $BC$, computed as an absolute determinant divided by $a$, is
$\frac a{2s}\cdot\frac{2\Delta}a=\Delta/s$, and the other two distances agree cyclically. Thus it
is the incentre with $r=\Delta/s$. For the bisector, put $A=0$, $B=c$, $C=b\cis\alpha$; then
\[
 D=\frac{bB+cC}{b+c}=\frac{bc}{b+c}(1+\cis\alpha)
 =\frac{2bc}{b+c}\cos(\alpha/2)\,\cis(\alpha/2)
\]
lies in the interior of $BC$, has $|B-D|/|D-C|=c/b$, and its direction from $A$ bisects the angle;
its modulus is \cref{trigonometry:eq:bisector}. Projection of an internal-bisector point onto either
ray lies at positive distance along that ray, which places each contact point of the incircle on
its own side segment.
\end{proof}

\begin{theorem}[Euler's incentre--circumcentre identity]\label{trigonometry:thm:euler}
If $O$ is the circumcentre and $I$ the incentre then
\begin{equation}\label{trigonometry:eq:euler}
 |O-I|^2=R(R-2r).
\end{equation}
Consequently $R\ge2r$, with equality exactly for an equilateral triangle.
\end{theorem}
\begin{proof}
Translate $O$ to zero, so $|A|=|B|=|C|=R$. For arbitrary positive weights $w_j$ and vectors $Z_j$,
expansion of the inner products gives the weighted variance identity
\[
 \left|\frac{\sum_jw_jZ_j}{\sum_jw_j}\right|^2
 =\frac{\sum_jw_j|Z_j|^2}{\sum_jw_j}
 -\frac{\sum_{j<k}w_jw_k|Z_j-Z_k|^2}{(\sum_jw_j)^2}.
\]
Apply it with $(w_1,w_2,w_3)=(a,b,c)$; the numerator of the second term is
$abc^2+acb^2+bca^2=2sabc$, so $|I|^2=R^2-abc/(2s)=R^2-2Rr$. Positivity gives $R\ge2r$. If equality
holds then $I=O$, so every side has the same positive distance $r$ from the centre and therefore
length $2\sqrt{R^2-r^2}$; the triangle is equilateral, and the converse is immediate.
\end{proof}
No limiting or Archimedean argument is hidden in any of these identities.

\begin{corollary}[A scale-independent area inequality]\label{trigonometry:cor:areabound}
Every nondegenerate surcomplex triangle satisfies $\Delta\le s^2/(3\sqrt3)$, with equality exactly
for an equilateral triangle.
\end{corollary}
\begin{proof}
The positive quantities $s-a,s-b,s-c$ sum to $s$. The arithmetic--geometric mean inequality for
three terms, proved algebraically over any real closed field, bounds their product by $(s/3)^3$;
Heron's formula converts this into the stated bound, and equality forces the three quantities to
coincide.
\end{proof}

\section{Bisectors, concurrency, chords, and cyclic quadrilaterals}
\label{trigonometry:sec:incidence}

\subsection{The cevian sine ratio and trigonometric Ceva}
\begin{theorem}[Sine ratios for a cevian]\label{trigonometry:thm:cevian}
Let $D$ lie strictly inside the side $BC$ of a noncollinear triangle. Then
\begin{equation}\label{trigonometry:eq:cevian}
 \frac{|B-D|}{|D-C|}=\frac{c\sin\angle BAD}{b\sin\angle DAC}.
\end{equation}
In particular the internal bisector of $\alpha$ meets $BC$ at the unique point with
$|B-D|/|D-C|=c/b$.
\end{theorem}
\begin{proof}
The triangles $ABD$ and $ADC$ share an altitude to the line $BC$, so their areas are in the ratio
$|B-D|/|D-C|$; computing the same areas from the common side $AD$ and the included angles at $A$
gives \cref{trigonometry:eq:cevian}. For the bisector, equal sines of equal positive half-angles
cancel. Uniqueness follows from the strict increase of $\sin x/\sin(\alpha-x)$ on $(0,\alpha)$,
which follows in turn from the sine subtraction identity, whose numerator is
$\sin\alpha\sin(y-x)>0$ for $x<y$.
\end{proof}

\begin{theorem}[Trigonometric Ceva]\label{trigonometry:thm:ceva}
Let $D,E,F$ lie strictly inside $BC$, $CA$, $AB$ respectively. The lines $AD,BE,CF$ are concurrent
if and only if
\begin{equation}\label{trigonometry:eq:trigceva}
 \frac{\sin\angle BAD}{\sin\angle DAC}\cdot
 \frac{\sin\angle CBE}{\sin\angle EBA}\cdot
 \frac{\sin\angle ACF}{\sin\angle FCB}=1.
\end{equation}
\end{theorem}
\begin{proof}
All six split angles lie in $(0,\pi)$, so their sines are positive. For an interior point with
positive barycentric coordinates $(p_A,p_B,p_C)$, its trace on $BC$ has ratio
$|B-D|/|D-C|=p_C/p_B$, the two cyclic ratios being $p_A/p_C$ and $p_B/p_A$, so their product is
one; conversely three prescribed positive ratios with product one determine such positive
coordinates, hence a common point. This proves the ordinary ratio form of Ceva by field algebra
alone, over $\No$, with no topological hypothesis. Substituting \cref{trigonometry:eq:cevian} and
its cyclic versions, the side factors $c/b$, $a/c$, $b/a$ cancel and give
\cref{trigonometry:eq:trigceva}.
\end{proof}
The restriction to interior traces avoids directed-length conventions. \emph{A directed Menelaus or
Ceva formula can be obtained in the same way after fixing signs, but is not needed for the present
development.} Even when a cevian divides a side at an infinitesimal ratio, the sine-ratio criterion
remains exact.

\subsection{Chords, inscribed angles, and Ptolemy}
A circle of centre $O\in\SC$ and radius $R\in\No_{>0}$ is \emph{exactly}
$\{O+R\cis\theta:\theta\in\OO_{\R}\}$: by \cref{trigonometry:thm:polar} the finite parameter
$\theta$ already reaches all of its points, including directions differing by arbitrarily small
infinitesimals. A circle of enormous radius therefore does \emph{not} require the function
$w\mapsto\Sin(Rw)$ to possess one Hahn support over an ordinary angular domain --- a point made
sharp in \cref{trigonometry:thm:infinitefrequency}. If $z_j=O+R\cis\theta_j$, the addition
identities give the exact factorization
\begin{equation}\label{trigonometry:eq:chordfactor}
 z_2-z_1=2iR\,\cis\frac{\theta_1+\theta_2}2\,\sin\frac{\theta_2-\theta_1}2,
 \qquad
 |z_2-z_1|=2R\left|\sin\frac{\theta_2-\theta_1}2\right|,
\end{equation}
at every surreal radius.

\begin{proposition}[Inscribed-angle theorem]\label{trigonometry:prop:inscribed}
For three distinct points $z_1,z_2,z_3$ on a circle, the directed angle from $z_1-z_3$ to
$z_2-z_3$, taken modulo $\pi$, is half the directed central angle from $z_1-O$ to $z_2-O$, halving
sending a central angle modulo $2\pi$ to an angle modulo $\pi$. In particular points on a common
arc subtend a fixed chord under equal interior angles, and an angle subtending a diameter is right.
\end{proposition}
\begin{proof}
Apply \cref{trigonometry:eq:chordfactor} to both vectors based at $z_3$: their quotient has phase
$(\theta_2-\theta_1)/2$ times a nonzero real scalar, and a real sign affects an argument only by
$\pi$. Choosing the actual interior-angle ranges gives the stated consequences. The argument works
equally for infinitesimally separated circle points.
\end{proof}

\begin{theorem}[Ptolemy inequality and its cyclic equality case]\label{trigonometry:thm:ptolemy}
For \emph{arbitrary} four surcomplex points,
\begin{equation}\label{trigonometry:eq:ptolemy}
 |z_1-z_3||z_2-z_4|\le|z_1-z_2||z_3-z_4|+|z_1-z_4||z_2-z_3|,
\end{equation}
and if the four points are distinct and cyclically ordered on a circle, equality holds.
\end{theorem}
\begin{proof}
The algebraic identity
\[
 (z_1-z_3)(z_2-z_4)=(z_1-z_2)(z_3-z_4)+(z_1-z_4)(z_2-z_3)
\]
and the norm triangle inequality prove \cref{trigonometry:eq:ptolemy}. For cyclic points choose
finite lifts $\theta_1<\theta_2<\theta_3<\theta_4<\theta_1+2\pi$. By
\cref{trigonometry:eq:chordfactor} the two products on the right have the same unit phase and
strictly positive real multiples of it, since all the half-difference sines involved have the
prescribed sign; equality therefore holds in the norm triangle inequality. Equivalently, with
$u=(\theta_2-\theta_1)/2$, $v=(\theta_3-\theta_2)/2$, $w=(\theta_4-\theta_3)/2$, all positive with
$u+v+w<\pi$, the equality reduces to the addition-formula identity
$\sin(u+v)\sin(v+w)=\sin u\sin w+\sin v\sin(u+v+w)$, whose two sides differ by
$\sin u\sin w(\cos^2v+\sin^2v-1)=0$. This includes the case in which some consecutive arcs are
infinitesimal.
\end{proof}
Here \emph{cyclic order is the order induced by finite representatives around the circle}; it is
\emph{not} defined by fine-continuous motion of a real parameter, which
\cref{trigonometry:prop:topology} has already shown to be degenerate. This illustrates a general
principle: algebraic identities among finitely many vectors often survive unchanged, while any
analytic interpretation of the associated angles must still be supplied.

\begin{corollary}[Finite regular polygons at surreal scale]\label{trigonometry:cor:polygon}
For an ordinary integer $n\ge3$ the regular $n$-gon with vertices $O+R\cis(2\pi k/n)$ has side
$2R\sin(\pi/n)$, perimeter $2nR\sin(\pi/n)$ and area $\frac n2R^2\sin(2\pi/n)$, for \emph{every}
positive surreal radius $R$.
\end{corollary}
\begin{proof}
The chord formula \cref{trigonometry:eq:chordfactor} gives the side; sum the finitely many equal
sides, and the finitely many positively oriented centre triangles, each of determinant area
$R^2\sin(2\pi/n)/2$.
\end{proof}
There is no literal polygon with ``$\omega$ vertices'' defined by this finite construction, and
\cref{trigonometry:prop:perimeters} shows that the sequence of these perimeters must not be used to
\emph{define} the circumference.

\subsection{Quadrance and spread}
The \emph{quadrance} of a side is its squared length and the \emph{spread} of two lines is the
square of the sine of either angle between them; these are the standard algebraic invariants of
rational trigonometry \cite{Wildberger,Gunn}. \emph{They forget orientation and distinguish neither
an angle from its supplement nor a directed angle from its negative. That loss is useful for
rational formulas but must be remembered.} Writing $Q_a=a^2$ and $\sigma_\alpha=\sin^2\alpha$
cyclically, squaring the sine and cosine laws gives
\begin{equation}\label{trigonometry:eq:spreadlaws}
 \frac{\sigma_\alpha}{Q_a}=\frac{\sigma_\beta}{Q_b}=\frac{\sigma_\gamma}{Q_c},
 \qquad
 (Q_b+Q_c-Q_a)^2=4Q_bQ_c(1-\sigma_\alpha),
\end{equation}
and the triple-spread relation is
\begin{equation}\label{trigonometry:eq:triplespread}
 (\sigma_\alpha+\sigma_\beta+\sigma_\gamma)^2
 =2(\sigma_\alpha^2+\sigma_\beta^2+\sigma_\gamma^2)+4\sigma_\alpha\sigma_\beta\sigma_\gamma .
\end{equation}
To prove the latter, use $\sin\gamma=\sin(\alpha+\beta)$ to obtain
$\sigma_\gamma-\sigma_\alpha-\sigma_\beta+2\sigma_\alpha\sigma_\beta
=2\sin\alpha\sin\beta\cos\alpha\cos\beta$, square both sides, replace each squared cosine by one
minus the corresponding spread, and expand. These formulas require no transcendental operation once
the coordinates are known and remain meaningful at every surreal scale.

\section{Valuation geometry and extremely thin triangles}
\label{trigonometry:sec:valuationgeometry}

\subsection{An exact dictionary between sides, angles and area}
For an interior angle $\alpha\in(0,\pi)$ its \emph{defect} is
$d(\alpha)=\min\{\alpha,\pi-\alpha\}\in(0,\pi/2]$, measuring proximity to either a zero angle or a
straight angle. The defect rather than $\alpha$ itself is the right variable: an angle
infinitesimally below $\pi$ has infinitesimal sine but is not infinitesimal as a number.

\begin{theorem}[Valuation form of the triangle laws]\label{trigonometry:thm:valuationtriangle}
For a noncollinear surcomplex triangle,
\begin{align}
 v(\sin\alpha)&=v(d(\alpha)),&
 v(a)-v(b)&=v(d(\alpha))-v(d(\beta)),\\
 v(\Delta)&=v(b)+v(c)+v(d(\alpha)),&
 2v(\Delta)&=v(s)+v(s-a)+v(s-b)+v(s-c),\\
 v(R)&=v(a)+v(b)+v(c)-v(\Delta),&
 v(r)&=v(\Delta)-v(s),
 \label{trigonometry:eq:valuationtriangle}
\end{align}
together with the cyclic versions; equivalently
$v(a)-v(d(\alpha))=v(b)-v(d(\beta))=v(c)-v(d(\gamma))=v(2R)$. The minimum of $v(a),v(b),v(c)$ is
attained \emph{at least twice}. If $\alpha$ and $\beta$ are both infinitesimal then
$\alpha/\beta\sim a/b$.
\end{theorem}
\begin{proof}
Jordan's inequality compares $\sin d(\alpha)$ with $d(\alpha)$ by nonzero ordinary constants, so
their valuations agree; since $\sin\alpha=\sin d(\alpha)$, the first assertion follows. If
$d(\alpha)$ is appreciable both quantities have positive ordinary standard part and valuation
zero. Applying $v$ to the sine law, the area formula, Heron's identity and the radius formulas
gives the rest, nonzero ordinary constants having valuation zero. If, say, $v(a)<v(b),v(c)$, then
$b/a$ and $c/a$ are positive infinitesimals whose sum is less than one, contradicting $a<b+c$; so
no side can be the unique dominant valuation scale. Finally, when the two angles are infinitesimal,
dividing the exact sine law by $\sin\alpha\sim\alpha$ and $\sin\beta\sim\beta$ gives the last
relation.
\end{proof}
The last statement is a non-Archimedean form of the triangle inequality. \emph{It does not say that
all three sides have equal valuation: a very short third side is possible.} Nor need one of
$\sin\alpha,\sin\beta,\sin\gamma$ be appreciable: in a nearly flat triangle one angle is
infinitesimally close to $\pi$ and the other two are infinitesimal, so all three sines are
infinitesimal.

Choosing a longest side and scaling it to length one, the remaining vertex is $x+iy$ with $y>0$ and
both coordinates finite, $0<x<1$; the normalized standard-part shadow is noncollinear exactly when
$y$ is appreciable. If $y$ is infinitesimal the shadow loses information that remains visible in
\cref{trigonometry:eq:valuationtriangle}. Two further theorems quantify what it loses.

\subsection{The triangle-inequality defect}
\begin{theorem}[A quadratic defect formula]\label{trigonometry:thm:defect}
Let $u,v\ne0$ with $a=|u|$, $b=|v|$, and let their unoriented angle be a nonzero infinitesimal
$\delta$. Set $D=a+b-|u+v|$. Then, exactly,
\begin{equation}\label{trigonometry:eq:defectexact}
 D=\frac{4ab\sin^2(\delta/2)}{a+b+|u+v|},
\end{equation}
and
\begin{equation}\label{trigonometry:eq:defectasym}
 D\sim\frac{ab}{2(a+b)}\delta^2,\qquad
 v(D)=v(a)+v(b)+2v(\delta)-v(a+b).
\end{equation}
\emph{No comparability assumption on $a$ and $b$ is required.}
\end{theorem}
\begin{proof}
Squaring and subtracting, $(a+b)^2-|u+v|^2=2ab(1-\cos\delta)=4ab\sin^2(\delta/2)$, and dividing by
the positive sum of the square roots gives \cref{trigonometry:eq:defectexact}. Also
\[
 \frac{|u+v|^2}{(a+b)^2}=1-\frac{4ab}{(a+b)^2}\sin^2(\delta/2),
\]
and since $0<4ab/(a+b)^2\le1$ the subtracted term is infinitesimal, so the positive square root is
$1$ plus an infinitesimal. Hence the denominator in \cref{trigonometry:eq:defectexact} is
relatively equivalent to $2(a+b)$, while $4\sin^2(\delta/2)\sim\delta^2$.
\end{proof}
Unlike a bound that replaces $a+b$ by an ordinary real constant, this expression \emph{detects the
effective smaller scale} when $a$ and $b$ differ by many orders of magnitude: the factor
$ab/(a+b)$ is comparable to the smaller of $a$ and $b$.

\subsection{Flat triangles: the normalized form and the relative threshold}
The next two theorems are genuinely different statements and both are kept. The first normalizes
coordinates and reads off every leading constant; the second imposes no normalization and isolates
the correct \emph{relative} flatness threshold.

\begin{theorem}[Quadratic slack and reciprocal circumradius in normalized coordinates]
\label{trigonometry:thm:flat}
Let $A=0$, $B=1$, $C=x+iy$ where $x$ is finite with $0<\st(x)<1$ and $0<y\in\mm_{\R}$, and put
$a=\sqrt{(1-x)^2+y^2}$, $b=\sqrt{x^2+y^2}$, $c=1$. Then, exactly,
\begin{equation}\label{trigonometry:eq:flatangles}
 \alpha=\arctan(y/x),\qquad \beta=\arctan\bigl(y/(1-x)\bigr),\qquad
 \gamma=\pi-\alpha-\beta,
\end{equation}
and
\begin{align}
 \alpha&=\frac yx-\frac{y^3}{3x^3}+O(y^5),&
 \beta&=\frac y{1-x}-\frac{y^3}{3(1-x)^3}+O(y^5),\\
 a+b-1&=\frac{y^2}{2x(1-x)}-\frac{y^4}8\left(\frac1{x^3}+\frac1{(1-x)^3}\right)+O(y^6),
 \label{trigonometry:eq:flatslack}\\
 R&=\frac{x(1-x)}{2y}\left[1+\frac{y^2}2\left(\frac1{x^2}+\frac1{(1-x)^2}\right)+O(y^4)\right],&
 r&=\frac y2\left[1-\frac{y^2}{4x(1-x)}+O(y^4)\right].
\end{align}
Consequently
\begin{equation}\label{trigonometry:eq:flatvals}
 v(\alpha)=v(\beta)=v(\pi-\gamma)=v(y),\qquad
 v(a+b-1)=2v(y),\qquad v(R)=-v(y),\qquad v(r)=v(y).
\end{equation}
\end{theorem}
\begin{proof}
Both $x$ and $1-x$ are positive appreciable, so the base angles are acute and their slopes give
\cref{trigonometry:eq:flatangles}; the inverse-tangent series
\cref{trigonometry:eq:smallseries} gives their expansions. Expanding
$b=x\sqrt{1+(y/x)^2}$ and $a=(1-x)\sqrt{1+(y/(1-x))^2}$ by the binomial series --- legitimate at
these infinitesimal arguments by \cref{trigonometry:lem:support} --- gives
\cref{trigonometry:eq:flatslack}. Since $\Delta=y/2$, the exact radius formulas are $R=ab/(2y)$ and
$r=y/(a+b+1)$, and multiplication and geometric inversion yield the displayed expansions. Each
indicated leading coefficient is positive appreciable, including the sum of the two leading angle
coefficients, so taking valuations gives \cref{trigonometry:eq:flatvals}.
\end{proof}
The theorem quantifies a phenomenon invisible in the ordinary standard-part shadow: the two small
angles are first order in the height, the excess in the triangle inequality is second order, and
the circumradius is of inverse order.

\begin{theorem}[Exact gap identity and the relative flatness threshold]
\label{trigonometry:thm:flatrelative}
Let $a$ be a largest side of a nondegenerate triangle and put $B=b+c$, $\delta=B-a>0$ and
$h=\pi-\alpha$. Then $0<\delta\le\min(b,c)$ and
\begin{equation}\label{trigonometry:eq:gapidentity}
 4bc\sin^2(h/2)=\delta(2B-\delta),
 \qquad\text{equivalently}\qquad
 h=2\arcsin\sqrt q,\quad q=\frac{\delta(2B-\delta)}{4bc}\in(0,1).
\end{equation}
The defect $h$ is infinitesimal \emph{if and only if}
\begin{equation}\label{trigonometry:eq:flatcriterion}
 \delta\prec\frac{bc}B .
\end{equation}
Under that criterion,
\begin{align}
 h&=2\sqrt q\left(1+\frac q6+\frac{3q^2}{40}+\frac{5q^3}{112}+\cdots\right),
 \label{trigonometry:eq:flatexpansion}\\
 h&\sim\sqrt{\frac{2B\delta}{bc}},\qquad
 \Delta\sim\sqrt{\frac{bcB\delta}2},\qquad
 R\sim\sqrt{\frac{Bbc}{8\delta}},\label{trigonometry:eq:flatpreview}
\end{align}
and
\begin{align}
 v(h)&=\tfrac12\bigl(v(\delta)+v(B)-v(b)-v(c)\bigr),\nonumber\\
 v(\Delta)&=\tfrac12\bigl(v(\delta)+v(B)+v(b)+v(c)\bigr),\label{trigonometry:eq:flatvaluations}\\
 v(R)&=\tfrac12\bigl(v(B)+v(b)+v(c)-v(\delta)\bigr).\nonumber
\end{align}
\end{theorem}
\begin{proof}
Since $a\ge\max(b,c)$, subtraction gives $\delta\le\min(b,c)$. The cosine law with $\alpha=\pi-h$
gives $B^2-a^2=2bc(1-\cos h)=4bc\sin^2(h/2)$, and $a=B-\delta$ turns this into
\cref{trigonometry:eq:gapidentity}; the half-angle formula
$\sin^2(h/2)=\cos^2(\alpha/2)=s(s-a)/(bc)$ of \cref{trigonometry:thm:heron} identifies the same
quantity with $q$, and $h/2\in(0,\pi/2)$ gives $h=2\arcsin\sqrt q$. The factor $2-\delta/B$ lies
between $3/2$ and $2$, because $\delta\le\min(b,c)\le B/2$; hence $\sin^2(h/2)$ is infinitesimal
exactly when $\delta B/(bc)$ is, and since $0<h/2<\pi/2$ its sine is infinitesimal exactly when the
angle is. This is the criterion.

Under the criterion $\delta/B$ is infinitesimal, because
$\delta/B=(\delta B/(bc))\cdot(bc/B^2)$ with $0<bc/B^2\le1/4$. The inverse-sine series gives
\cref{trigonometry:eq:flatexpansion}. From \cref{trigonometry:eq:gapidentity} and
$\sin(h/2)\sim h/2$ we get $bc\,h^2\sim2B\delta$, and positive square roots give the first
equivalence in \cref{trigonometry:eq:flatpreview}. The area law gives
$\Delta=\frac{bc}2\sin h\sim\frac{bc}2h$, proving the second; and $a\sim B$ with
$R=a/(2\sin h)\sim B/(2h)$ gives the third. Taking valuations, and noting that $2B-\delta$ and $B$
have the same valuation, proves \cref{trigonometry:eq:flatvaluations}.
\end{proof}
\emph{The denominator $bc/B$ is essential}, and the hypothesis is deliberately \emph{relative}
rather than an absolute smallness assumption: since $bc/(b+c)$ is comparable to the smaller of $b$
and $c$, the gap must be small relative to the \emph{shorter adjacent side}, not merely relative to
the total length. A defect that is small compared with the perimeter need not make the opposite
angle close to $\pi$ when one adjacent side is much shorter than the other. The square root of the
side gap controls the angular defect while its reciprocal controls the circumradius; this is the
same ramification as in \cref{trigonometry:prop:acosendpoint}.

\begin{example}[A symmetric infinitesimal-height triangle]\label{trigonometry:ex:flat}
Take $x=1/2$, $y=\eps>0$ infinitesimal in \cref{trigonometry:thm:flat}. Then
\[
 a=b=\tfrac12\sqrt{1+4\eps^2},\qquad \alpha=\beta=\arctan(2\eps),\qquad
 R=\frac1{8\eps}+\frac\eps2,
\]
\[
 r=\frac\eps{1+\sqrt{1+4\eps^2}}=\frac\eps2-\frac{\eps^3}2+\eps^5+O(\eps^7).
\]
The circumradius is infinite although all vertices lie within a bounded ordinary region. Its
existence is not contradicted by the fact that the standard-part points are collinear.
\end{example}

\begin{example}[A symmetric flattening, seen through the gap]\label{trigonometry:ex:flatgap}
Let $0<\tau\in\mm_{\R}$, $b=c=1$, $a=2-\tau$, so $\delta=\tau$ and $B=2$ in
\cref{trigonometry:thm:flatrelative}. Exactly,
\[
 H=\sqrt{\tau-\tau^2/4},\qquad \Delta=(1-\tau/2)H,\qquad R=\frac1{2H},
 \qquad h=2\arccos(1-\tau/2),\qquad \beta=\gamma=h/2,
\]
with the strongly summable expansions
\begin{align*}
 h&=2\sqrt\tau\left(1+\frac\tau{24}+\frac{3\tau^2}{640}+\frac{5\tau^3}{7168}+\cdots\right),\\
 H&=\sqrt\tau\left(1-\frac\tau8-\frac{\tau^2}{128}-\frac{\tau^3}{1024}+\cdots\right),\\
 R&=\frac1{2\sqrt\tau}\left(1+\frac\tau8+\frac{3\tau^2}{128}+\frac{5\tau^3}{1024}+\cdots\right).
\end{align*}
Here the side lengths have ordinary scale, the altitude and the area have square-root infinitesimal
scale, and the circumradius is infinite.
\end{example}

\begin{example}[An appreciable side gap can still be a flat degeneration]
\label{trigonometry:ex:appreciablegap}
Take $b=c=\omega^2$ and $\delta=1$, so $a=2\omega^2-1$. The criterion
\cref{trigonometry:eq:flatcriterion} holds, since $bc/B=\omega^2/2$. Then
$h\sim2/\omega$, $\Delta\sim\omega^3$ and $R\sim\omega^3/2$: the triangle is angularly nearly flat
although its area is infinite and its gap is the ordinary number $1$. This is precisely why
degeneration hypotheses must be stated relatively.
\end{example}

\begin{example}[An infinite triangle of ordinary finite area]\label{trigonometry:ex:infinite}
Let $L>0$ be infinite and take $A=0$, $B=L$, $C=L+i/L$. The legs at $B$ are $L$ and $L^{-1}$, and
\[
 \Delta=\tfrac12,\qquad \alpha=\arctan(L^{-2}),\qquad\beta=\frac\pi2,\qquad
 \gamma=\frac\pi2-\alpha,
\]
with circumradius $\frac12\sqrt{L^2+L^{-2}}$ and inradius
\[
 r=\frac{L+L^{-1}-\sqrt{L^2+L^{-2}}}2=\frac1{2L}-\frac1{4L^3}+O(L^{-7}).
\]
The special case $L=\omega$ is the right triangle with vertices $0,\omega,i/\omega$, of legs
$\omega$ and $\omega^{-1}$, area exactly $1/2$, hypotenuse $\omega\sqrt{1+\omega^{-4}}$, nonright
angles $\arctan(\omega^{-2})$ and $\pi/2-\arctan(\omega^{-2})$, and circumradius $\sim\omega/2$.
Thus an infinite side, an infinitesimal side, an ordinary area and an infinitesimal angle coexist
without disturbing a single triangle identity; standard parts of the normalized small angle alone
would erase information essential to both the height and the area.
\end{example}

\begin{example}[Two thin triangles on different scale hierarchies]
\label{trigonometry:ex:hierarchies}
Take $A=0$, $B=L$, $C=x+ih$ with $0<x<L$ and $0<h\prec\min\{x,L-x\}$. Then, exactly,
$\alpha=\arctan(h/x)$, $\beta=\arctan(h/(L-x))$, $\gamma=\pi-\alpha-\beta$, $\Delta=Lh/2$, and
\[
 R=\frac{\sqrt{x^2+h^2}\sqrt{(L-x)^2+h^2}}{2h}\sim\frac{x(L-x)}{2h},\qquad r\sim h/2 .
\]
For $L=\omega$, $x=\omega/2$, $h=1$, both base angles equal
$\arctan(2/\omega)=2/\omega-8/(3\omega^3)+\cdots$, while $R\sim\omega^2/8$, $r\sim1/2$ and
$\Delta=\omega/2$: infinite area, infinite circumradius, finite inradius, infinitesimal base
angles. For $L=1$, $x=\eps$, $h=\eps^2$ with $\eps>0$ infinitesimal, the two base angles sit on
\emph{different} scales, $\alpha\sim\eps$ and $\beta\sim\eps^2$, with $R\sim1/(2\eps)$ and
$r\sim\eps^2/2$. Both families obey exactly the same identities as ordinary triangles, but their
scale information cannot be represented by ordinary positive real lengths and angles alone.
\end{example}

\section{Trigonometric intersections, contact, and conditioning}
\label{trigonometry:sec:contact}

This section specializes the finite-intersection viewpoint of
\cite{SourceGeometry,SourceIntersections} to angular equations. Rather than invoke their general
division and residue theorems, every instance is computed directly; the geometric theory above is
therefore independent of them.

\subsection{A complete first-harmonic intersection theorem}
\begin{theorem}[Amplitude--phase reduction and intersection count]
\label{trigonometry:thm:amplitude}
Let $A,B,D\in\No$ with $(A,B)\ne(0,0)$ and $\rho=\sqrt{A^2+B^2}>0$, and choose the finite angle
class $\phi$ with $\cos\phi=A/\rho$ and $\sin\phi=B/\rho$. Then
\begin{equation}\label{trigonometry:eq:amplitude}
 A\cos\theta+B\sin\theta=\rho\cos(\theta-\phi),
\end{equation}
and the equation $A\cos\theta+B\sin\theta=D$ has, modulo $2\pi$ in finite angles, no solution if
$|D|>\rho$, exactly one if $|D|=\rho$, and exactly two if $|D|<\rho$, namely
$\theta=\phi\pm\arccos(D/\rho)$. The single solution at $|D|=\rho$ has multiplicity two: the first
angular derivative vanishes there and the second equals $\mp\rho\ne0$. Writing
$\mathcal D=\rho^2-D^2$, the two unit points at $\mathcal D>0$ are
\begin{equation}\label{trigonometry:eq:intersectionpoints}
 (X,Y)=\frac D{\rho^2}(A,B)\pm\frac{\sqrt{\mathcal D}}{\rho^2}(-B,A),
\end{equation}
the absolute angular derivative of the left-hand side at either simple solution is
$\sqrt{\mathcal D}$, and the Jacobian determinant of $X^2+Y^2-1$ and $AX+BY-D$ has absolute value
$2\sqrt{\mathcal D}$.
\end{theorem}
\begin{proof}
Polar decomposition (\cref{trigonometry:thm:polar}) gives $\phi$, and the cosine subtraction
identity proves \cref{trigonometry:eq:amplitude}. The range, strict monotonicity and inverse cosine
on $[0,\pi]$ (\cref{trigonometry:prop:order}, \cref{trigonometry:thm:inverse}) give the count,
including both endpoints, where the two signs coincide modulo $2\pi$. For the explicit points use
the orthonormal coordinates $U=(AX+BY)/\rho$ and $W=(-BX+AY)/\rho$: the equations become
$U=D/\rho$ and $W^2=1-D^2/\rho^2=\mathcal D/\rho^2$. The angular derivative is
$-A\sin\theta+B\cos\theta=BX-AY=-\rho W$, of absolute value $\sqrt{\mathcal D}$, and the stated
Jacobian expresses the same transversality.
\end{proof}
Geometrically this is the complete intersection classification of a line $AX+BY=D$ with the unit
circle. Neither an infinite normal vector nor an infinitesimal distance from tangency creates an
exception. If $A=B=0$ the equation instead holds for every angle when $D=0$ and for none
otherwise. This also settles the side--side--angle ambiguity left open after
\cref{trigonometry:thm:sss}: the ray--circle problem is exactly this quadratic.

\subsection{Square-root splitting and its conditioning}
\begin{theorem}[Tangency, angular splitting, and loss of valuation]
\label{trigonometry:thm:tangency}
In \cref{trigonometry:thm:amplitude} suppose $D=\rho(1-\tau)$ with $0<\tau\in\mm_{\R}$. The two
solutions near $\phi$ are
\begin{equation}\label{trigonometry:eq:tangency}
 \theta_\pm=\phi\pm\sqrt{2\tau}\left(1+\frac\tau{12}+\frac{3\tau^2}{160}+O(\tau^3)\right),
\end{equation}
their separation has valuation $\tfrac12v(\tau)$, and the derivatives of the left-hand side at
these roots have absolute value $\rho\sqrt{\tau(2-\tau)}$. If $e\ne0$ and $e\prec\tau$, then with
$q(\tau)=\arccos(1-\tau)$ one has $\tau+e>0$ and
\begin{equation}\label{trigonometry:eq:contactstability}
 q(\tau+e)-q(\tau)=\frac e{\sqrt{\tau(2-\tau)}}\bigl(1+O(e/\tau)\bigr),
\end{equation}
the correction factor having standard part one; in particular
\begin{equation}\label{trigonometry:eq:contactval}
 v\bigl(q(\tau+e)-q(\tau)\bigr)=v(e)-\tfrac12v(\tau).
\end{equation}
\end{theorem}
\begin{proof}
Equation \cref{trigonometry:eq:tangency} is \cref{trigonometry:prop:acosendpoint}, and the
separation is $2q(\tau)$, whose valuation is unchanged by the ordinary factor two. Differentiating
\cref{trigonometry:eq:amplitude} gives magnitude $\rho\sin q(\tau)$ and
$\sin q(\tau)=\sqrt{1-(1-\tau)^2}=\sqrt{\tau(2-\tau)}$. For the perturbation formula write
$e=\tau u$ with $u$ infinitesimal; the exact expression $q(\tau)=\sqrt{2\tau}H(\tau)$ from
\cref{trigonometry:eq:acosseries}, with $H(0)=1$, may be expanded in the two infinitesimals
$\tau,u$. Its first derivative in $e$ is $1/\sqrt{\tau(2-\tau)}$, and every later term contains,
relative to that first term, a factor $u=e/\tau$ times a finite factor; the substitution is
support-controlled by \cref{trigonometry:lem:support}. The first term being nonzero and the
parenthesized multiplier a unit of standard part one gives \cref{trigonometry:eq:contactval}.
\end{proof}
The loss of one half of $v(\tau)$ is the inverse-Jacobian loss at a nearly double root. \emph{The
condition $e\prec\tau$ cannot be omitted from a uniform branch-matching theorem: $e=-\tau$ merges
the roots, and $e<-\tau$ destroys the two real intersections.}

\begin{example}[An infinite tangent coordinate]\label{trigonometry:ex:infinitetangent}
The line $X=-1+\tau$ meets the unit circle near $-1$ in the two points
$Y=\pm\sqrt{\tau(2-\tau)}$, whose half-angle coordinates from \cref{trigonometry:eq:cayley} are
$t=\pm\sqrt{(2-\tau)/\tau}$, and these are \emph{infinite} even though the points and their angles
are finite. Rational circle coordinates and finite angles therefore complement each other rather
than having identical size behaviour.
\end{example}

\begin{theorem}[The cosine fold]\label{trigonometry:thm:fold}
Let $\tau\in\mm_{\SC}$. If $\tau\ne0$, the solutions in the infinitesimal monad of zero of
$\cos\theta=1-\tau$ are exactly
\begin{equation}\label{trigonometry:eq:foldroots}
 \theta_\pm=\pm2\arcsin_{\mathrm{inf}}\sqrt{\tau/2}
 =\pm\sqrt{2\tau}\left(1+\frac\tau{12}+\frac{3\tau^2}{160}+\frac{5\tau^3}{896}+\cdots\right),
\end{equation}
where $\arcsin_{\mathrm{inf}}$ is the zero-centred formal inverse-sine germ and either square root
may be selected. Both are simple, with $v(\theta_\pm)=\tfrac12v(\tau)$ exactly. If $\tau=0$, zero
is the unique root there, of multiplicity two. For real $\tau>0$ both roots are real; for real
$\tau<0$ there is no real solution near zero, but there are exactly two conjugate purely imaginary
surcomplex solutions, each of multiplicity one.
\end{theorem}
\begin{proof}
On the infinitesimal monad, $1-\cos\theta=2\sin^2(\theta/2)$, and $\theta\mapsto\sin(\theta/2)$ is
a germ with nonzero linear coefficient, hence a bijection of that monad onto itself with inverse
$w\mapsto2\arcsin_{\mathrm{inf}}w$; all its formal inverse identities evaluate by
\cref{trigonometry:lem:support}. The equation becomes $w^2=\tau/2$, with exactly the two stated
roots. At a nonzero root $\sin\theta=2\sin(\theta/2)\cos(\theta/2)\ne0$, so the derivative is
nonzero; at $\tau=0$, $\cos\theta-1=-\theta^2/2+\cdots$ gives the double root. The expansion is
\cref{trigonometry:eq:acosseries}. Its coefficients are real and its germ odd, so a real positive
square root produces real roots and a purely imaginary square root produces purely imaginary ones,
while a real square cannot equal $\tau/2<0$.
\end{proof}
\emph{The complex multiplicity is constant across this real bifurcation}, which is the one-variable
version of the finite-algebra viewpoint of \cite{SourceGeometry}. The explicit inverse germ
supplies the proof; no sequence of Newton approximations is used as a convergence argument. The
same phenomenon in the strip is visible in \cref{trigonometry:ex:doublesplit}.

\subsection{The rank-two algebra that survives the collision}
Put $d=\mathcal D/\rho^2$ in \cref{trigonometry:thm:amplitude}. Over $\SC$ the coordinate algebra
of the line--circle intersection is
\begin{equation}\label{trigonometry:eq:quadalgebra}
 \mathcal A_d=\SC[W]/(W^2-d),
\end{equation}
whose elements are the unique remainders $a+bW$ with $W^2=d$ defining the multiplication; this
interprets the notation as a class algebra without forming a collection of proper-class cosets. It
has basis $1,W$ and dimension two \emph{for every} $d$, including $d=0$ and negative real $d$.
When $d<0$ there are no real points but the two complexified points still exist; when $d=0$ there
is a single point with the nonreduced algebra of dual numbers. This is the precise sense in which
multiplicity survives the change in the number of real solutions.

\begin{theorem}[A collision-stable residue pairing]\label{trigonometry:thm:residuepairing}
Define $\Lambda_d:\mathcal A_d\to\SC$ by $\Lambda_d(a+bW)=b$, the coefficient of $W$ in the unique
remainder. The pairing $(f,g)\mapsto\Lambda_d(fg)$ is nondegenerate for every $d$, with matrix
\begin{equation}\label{trigonometry:eq:pairmatrix}
 \begin{pmatrix}0&1\\1&0\end{pmatrix}
\end{equation}
in the basis $1,W$, of determinant $-1$. If $d\ne0$ and $s^2=d$ then for every polynomial $g$,
\begin{equation}\label{trigonometry:eq:residuesum}
 \Lambda_d(g)=\frac{g(s)}{2s}+\frac{g(-s)}{-2s}=\frac{g(s)-g(-s)}{2s},
\end{equation}
which is exactly the sum of the two simple residues of $g(W)/(W^2-d)$; at $d=0$ the corresponding
formula is $\Lambda_0(g)=g'(0)$. The trace of multiplication by $g$ is $\Lambda_d(2Wg)$.
\end{theorem}
\begin{proof}
The relation $W^2=d$ gives $\Lambda_d(1)=0$, $\Lambda_d(W)=1$ and $\Lambda_d(W^2)=\Lambda_d(d)=0$,
which is \cref{trigonometry:eq:pairmatrix}. Reducing $g$ modulo $W^2-d$ gives $g=a+bW$, whose
values at $\pm s$ are $a\pm bs$, so the right-hand side of \cref{trigonometry:eq:residuesum} is
$b$. At $d=0$ the reduction is $g(0)+g'(0)W$, and the coefficient of $W^{-1}$ in $g(W)/W^2$ is
$g'(0)$. Finally multiplication by $a+bW$ has matrix
$\left(\begin{smallmatrix}a&bd\\ b&a\end{smallmatrix}\right)$, of trace
$2a=\Lambda_d(2W(a+bW))$.
\end{proof}
For infinitesimal positive $d$ the two residues in \cref{trigonometry:eq:residuesum} may
individually be infinite while their sum remains well behaved: for $g=1$ they are
$\pm(2\sqrt d)^{-1}$ and cancel exactly, while for $g=W$ they are both $1/2$ and sum to $1$. This
elementary model exhibits the structural phenomenon emphasized in
\cite{SourceGeometry,SourceIntersections}: \emph{the quotient and its pairing are more stable than
the list of separated simple points.}

\subsection{A coupled four-point angular collision}\label{trigonometry:sec:coupled}
Consider two infinitesimal angles $\alpha,\beta\in\mm_{\C}$ and parameters $s,u\in\mm_{\C}$ with
\begin{equation}\label{trigonometry:eq:coupled}
 \sin^2\alpha+\sin^2\beta=s,\qquad \sin\alpha\sin\beta=u .
\end{equation}
(Here, and only here, the letters $s,u$ denote these parameters rather than a semiperimeter or a
unit direction.) Set $x=\sin\alpha$, $y=\sin\beta$ and then $X=x+y$, $Y=x-y$; both coordinate
changes are invertible analytic germs, and the system becomes
\begin{equation}\label{trigonometry:eq:coupleddiagonal}
 X^2=s+2u,\qquad Y^2=s-2u .
\end{equation}
The associated finite intersection algebra is
\[
 \mathcal B_{s,u}=\SC[X,Y]/(X^2-s-2u,\ Y^2-s+2u).
\]
Division by the two monic quadratics gives the basis $1,X,Y,XY$, so
$\dim_{\SC}\mathcal B_{s,u}=4$ for \emph{every} infinitesimal parameter value. Every support point
lies in the infinitesimal monad, since its coordinate squares are infinitesimal; the total complex
multiplicity there is therefore always four. \emph{This is the algebra of the entire finite zero
set, not the local algebra at one selected root}: it decomposes into the local algebras of the
distinct roots, and the invertible analytic sine-coordinate changes preserve each local
multiplicity. The angle solutions are
\begin{equation}\label{trigonometry:eq:coupledroots}
 \alpha=\arcsin\frac{\epsilon_1\sqrt{s+2u}+\epsilon_2\sqrt{s-2u}}2,
 \qquad
 \beta=\arcsin\frac{\epsilon_1\sqrt{s+2u}-\epsilon_2\sqrt{s-2u}}2,
\end{equation}
with $\epsilon_1,\epsilon_2\in\{1,-1\}$ and repetitions counted by local multiplicity; the
square-root choices may be fixed arbitrarily and the sign choices exhaust them.

For real infinitesimal parameters there are four distinct real solutions exactly when $s+2u>0$ and
$s-2u>0$; if precisely one of those quantities vanishes and the other is positive there are two
distinct real solutions of multiplicity two; at $s=u=0$ there is one solution of multiplicity four;
if either quantity is negative no solution has both angles real, because real $X,Y$ have
nonnegative squares; and if neither vanishes there are always four distinct complex solutions. The
angular Jacobian is
\[
 J=2\cos\alpha\cos\beta\,(x^2-y^2)=2\cos\alpha\cos\beta\,XY,
\]
whose vanishing locus in parameter space is $s^2-4u^2=0$. The cosine factors are units in the
infinitesimal monad, so they change neither the collision locus nor any multiplicity. This is the
only genuinely multivariable instance in this report, and the dimension four is a \emph{total}
multiplicity, not a claim about any individual root.

\subsection{Three stability certificates, and what each is for}
\label{trigonometry:sec:threecertificates}
The remainder of this section proves two sharp perturbation theorems; a third, of a different
shape, is \cref{trigonometry:thm:cluster} below. They are genuinely different statements and none
subsumes another, so all three are kept.
\begin{itemize}
\item \Cref{trigonometry:thm:stability} bounds the \emph{displacement of one simple root} of a
 trigonometric polynomial with finite coefficients under a perturbation of coefficient valuation
 $\sigma$, under the threshold $\sigma>2\kappa$ with $\kappa=v(f'(a))$, and
 \cref{trigonometry:prop:sharp} proves both that the threshold cannot be weakened to $\sigma\ge
 2\kappa$ and that both displayed exponents are attained. It is the sharpest statement available
 when one tracks a single named root of a polynomial with finite coefficients.
\item \Cref{trigonometry:thm:conditioned} instead inverts \emph{one scalar function}, the inverse
 cosine, at a prescribed interior angle, and supplies the second-order term $-c\eps^2/(2s^3)$
 together with a remainder bound $v(\mathcal R)\ge3\sigma-5\kappa$. It is sharper than
 \cref{trigonometry:thm:stability} whenever the second-order behaviour of the inverse is what one
 needs, and it applies at an angle rather than to a polynomial; its own sharpness construction is
 different and explicit.
\item \Cref{trigonometry:thm:cluster} certifies instead an \emph{exact multiplicity} for a whole
 cluster of roots, together with a coefficient-support bound, and makes no smallness demand beyond
 infinitesimality. It is the right tool when roots collide and individual displacements are not
 the object of interest, and it is the only one of the three that survives a multiple unperturbed
 root.
\end{itemize}

\begin{lemma}[A simple residue root of a polynomial]\label{trigonometry:lem:hensel}
Let $H(Y)$ be a real polynomial with finite surreal coefficients whose coefficientwise standard
part is $Y+b_0$, $b_0\in\R$. Then $H$ has exactly one root with standard part $-b_0$, and that
root is simple.
\end{lemma}
\begin{proof}
Put $y_0=-b_0$ and $\eta=H(y_0)\in\mm_{\R}$. If $\eta=0$ then $y_0$ is a root. Otherwise choose a
positive infinitesimal $\varrho$ with $|\eta|\prec\varrho$, for instance $\varrho=\sqrt{|\eta|}$.
Polynomial expansion gives $H(y_0\pm\varrho)=\eta\pm\varrho(1+\vartheta_\pm)$ with
$\vartheta_\pm\in\mm_{\R}$, so the two values have opposite signs, and the polynomial
intermediate-value property of the real closed field $\No$ supplies a root between them. For any
finite $y,z$ with the required standard part the divided difference is $1$ plus an infinitesimal,
because the polynomial is finite and its reduced polynomial is linear; hence it is nonzero, which
proves uniqueness, and the corresponding derivative calculation proves simplicity.
\end{proof}
This replaces a Newton-iteration limit by a single intermediate-value step with an infinitesimal
test radius --- the only available substitute, since by \cref{trigonometry:prop:topology} a
Newton sequence has no fine-topological limit to converge to.

\begin{theorem}[Sharp angular root-stability bound]\label{trigonometry:thm:stability}
Let $f$ be a real-valued trigonometric polynomial with \emph{finite coefficients} --- i.e.\ each
Laurent coefficient lies in $\OO_{\C}$, finite in magnitude, \emph{not} finite in number --- and
let $a\in\OO_{\R}$ be a simple zero. Put $A=f'(a)$ and $\kappa=v(A)\ge0$. Let $g$ be a real-valued
trigonometric polynomial all of whose Laurent coefficients have valuation at least $\sigma$, where
\begin{equation}\label{trigonometry:eq:stabilitythreshold}
 \sigma>2\kappa .
\end{equation}
Then $f+g$ has exactly one zero $a+h$ in the angular neighbourhood $v(h)>\kappa$; it is simple, and
\begin{align}
 v(h)&\ge\sigma-\kappa,\label{trigonometry:eq:stabilityfirst}\\
 v\left(h+\frac{g(a)}{f'(a)}\right)&\ge2\sigma-3\kappa.\label{trigonometry:eq:stabilitysecond}
\end{align}
The assertion holds for value groups of arbitrary rank; $\sigma$ and $\kappa$ need not be real
numbers.
\end{theorem}
\begin{proof}
All Taylor coefficients of $f(a+H)$ are finite, and those of $g(a+H)$ have valuation at least
$\sigma$, since only finitely many frequencies occur and $\cis a$ has valuation zero. Write
\[
 f(a+H)=AH+Q(H),\qquad g(a+H)=g(a)+BH+Q_g(H),
\]
where $Q,Q_g$ have order at least two with coefficient valuations at least $0$ and $\sigma$
respectively, and $v(B)\ge\sigma$. Set $r=\sigma-\kappa>\kappa$ and $\lambda=t^r$. Formally
substituting $H=\lambda Y$ and dividing by $A\lambda$ gives $Y+b+P(Y)$ with
$b=g(a)/(A\lambda)\in\OO_{\R}$, every nonconstant error coefficient of $P$ having strictly positive
valuation: the normalized degree-$j$ term with $j\ge2$ has valuation at least $(j-1)r-\kappa>0$.

To turn this into an existence proof without an iteration limit, use the exact rational chart of
\cref{trigonometry:thm:cayley},
\[
 H=2\arctan(\lambda Y/2),\qquad \cis(a+H)=\cis a\cdot C(\lambda Y/2).
\]
Since $f$ and $g$ have finite trigonometric degree, their pullbacks are rational functions of $Y$;
clearing the common denominator, whose standard part is one for finite $Y$, and dividing by
$A\lambda$, yields a finite polynomial whose coefficientwise reduction is $Y+\st(b)$. The chart's
difference from $H=\lambda Y$ begins in degree three and also has positive normalized valuation.
\Cref{trigonometry:lem:hensel} now gives a finite root, whence $h$ with $v(h)\ge r$: this is
existence and \cref{trigonometry:eq:stabilityfirst}.

For uniqueness on the whole stated neighbourhood take $v(h_1),v(h_2)>\kappa$. Strong Taylor divided
differences give $(f+g)(a+h_1)-(f+g)(a+h_2)=(A+D)(h_1-h_2)$ with
$v(D)\ge\min\{v(h_1),v(h_2),\sigma\}>\kappa$, so $A+D\ne0$; the same estimate for the derivative at
the new root proves simplicity. These divided-difference series are legitimate strong sums: each
nonlinear term introduces an infinitesimal factor, and the finitely many original coefficient
supports give a common well-ordered support by \cref{trigonometry:lem:support}.

Finally the root equation gives $h+A^{-1}g(a)=-A^{-1}(Q(h)+Bh+Q_g(h))$, and $v(h)\ge r$ bounds its
valuation below by $\min\{2r-\kappa,\ \sigma+r-\kappa,\ \sigma+2r-\kappa\}=2\sigma-3\kappa$, which
is \cref{trigonometry:eq:stabilitysecond}.
\end{proof}
This is the scalar angular counterpart of the inverse-Jacobian theorem of
\cite[\S9]{SourceIntersections}, with a direct algebraic existence argument available because a
trigonometric polynomial becomes rational in a tangent chart.

\begin{proposition}[The threshold and both exponents are sharp]\label{trigonometry:prop:sharp}
For uniform simple-root persistence, $\sigma>2\kappa$ cannot be weakened to $\sigma\ge2\kappa$, and
the exponents in \cref{trigonometry:eq:stabilityfirst} and
\cref{trigonometry:eq:stabilitysecond} are attained.
\end{proposition}
\begin{proof}
Let $0<\tau\in\mm_{\R}$ and $f(\theta)=\sin^2\theta-\tau^2$. At $a=\arcsin\tau$,
$f'(a)=2\tau\sqrt{1-\tau^2}$, so $\kappa=v(\tau)$. Adding the constant $g=\tau^2$, of valuation
$2\kappa$, merges the two roots near zero into the double root of $\sin^2\theta$: at
$\sigma=2\kappa$ there is no uniform simple-root persistence.

For $g=d$ with $v(d)=\sigma>2\kappa$, the root near $a$ is $a+h=\arcsin\sqrt{\tau^2-d}$, and
expansion in the infinitesimal ratio $d/\tau^2$ gives
\[
 h=-\frac d{2\tau\sqrt{1-\tau^2}}
 +\frac{(2\tau^2-1)d^2}{8\tau^3(1-\tau^2)^{3/2}}+\text{terms of higher valuation}.
\]
The first coefficient has valuation $-\kappa$ and the second $-3\kappa$, while successive terms
gain at least $\sigma-2\kappa>0$ in valuation. Thus $v(h)=\sigma-\kappa$ and
$v(h+g(a)/f'(a))=2\sigma-3\kappa$ for this family: both estimates are equalities. The expansion is
justified by strong binomial and inverse-sine composition, not by any rank-one convergence
assertion.
\end{proof}

\begin{theorem}[Conditioned inversion of cosine]\label{trigonometry:thm:conditioned}
Let $\theta\in(0,\pi)$, $c=\cos\theta$, $s=\sin\theta>0$, and let $0\ne\eps\in\No$ satisfy
\begin{equation}\label{trigonometry:eq:condition}
 \sigma:=v(\eps)>2\kappa,\qquad \kappa:=v(s)\ge0 .
\end{equation}
Then $c+\eps\in(-1,1)$, and with $h=\arccos(c+\eps)-\theta$,
\begin{equation}\label{trigonometry:eq:inversecos}
 h=-\frac\eps s-\frac{c\eps^2}{2s^3}+\mathcal R,
 \qquad v(\mathcal R)\ge3\sigma-5\kappa,
 \qquad v(h)=\sigma-\kappa,
 \qquad v\!\left(h+\frac\eps s\right)\ge2\sigma-3\kappa .
\end{equation}
\emph{The strict threshold in \cref{trigonometry:eq:condition} cannot be weakened to equality
uniformly near an extremum.}
\end{theorem}
\begin{proof}
The distance from $c$ to the nearer endpoint of $[-1,1]$ is $1-|c|=s^2/(1+|c|)$. Condition
\cref{trigonometry:eq:condition} says that $\eps/s^2$ is infinitesimal, so $|\eps|$ is smaller than
that distance; hence $c+\eps\in(-1,1)$.

Set $u=\eps/s^2$ and write $h=sq$. The addition formula gives
\begin{equation}\label{trigonometry:eq:normalizedinverse}
 \frac{\cos(\theta+sq)-c}{s^2}
 =c\,\frac{\cos(sq)-1}{s^2}-\frac{\sin(sq)}s
 =-q-\frac c2q^2+\frac{s^2}6q^3+\cdots,
\end{equation}
every coefficient on the right being a polynomial with rational coefficients in the finite
parameters $c$ and $s^2$, and the linear coefficient being $-1$. Formal inversion therefore yields
$q=-u-\frac c2u^2+u^3H(u)$ with every coefficient of $H$ again such a polynomial. Evaluation at the
infinitesimal $u$ is legitimate: write the finitely many parameters as ordinary standard parts plus
infinitesimals; all expanded terms are supported in the positive monoid generated by those
infinitesimals and $u$, and positive powers of $u$ ensure a finite contribution at each exponent
(\cref{trigonometry:lem:support}, and \cref{trigonometry:rem:notlimit}(ii) for why finitely many
parameters suffice). Thus $H(u)$ is finite, $v(H(u))\ge0$.

The constructed $h=sq$ has $h/s\in\mm_{\R}$ and therefore stays inside the same interior-angle
branch, the distance to either endpoint being at least $s$ because
$s=\sin\theta\le\min(\theta,\pi-\theta)$ by \cref{trigonometry:prop:order}. It solves the desired
equation, so strict monotonicity of cosine identifies it with the stated difference of inverse
values. Multiplying the inverse series by $s$ gives $\mathcal R=su^3H(u)$ with
$v(\mathcal R)\ge\kappa+3(\sigma-2\kappa)=3\sigma-5\kappa$, while the linear term has valuation
$\sigma-\kappa$ and every later term strictly larger valuation.

For sharpness let $\eta>0$ be infinitesimal and put $c=1-\eta$, so $s^2=2\eta-\eta^2$ has valuation
$v(\eta)$. Taking $\eps=\eta$ gives $\sigma=2\kappa$ and moves the target to $1$, merging the two
real roots, so the perturbed root is no longer simple; taking $\eps=2\eta$, of the same valuation,
moves the target above $1$ and the real solutions disappear altogether. A uniform simple-root
stability theorem therefore cannot replace $>$ by $\ge$.
\end{proof}
The roles of the two powers of the small derivative are now explicit: one factor of $s^{-1}$
amplifies the displacement, the other measures how close the target is to a critical value. The
leading displacement has order $\sigma-\kappa$ and the nonlinear correction order at least
$2\sigma-3\kappa$, matching the inverse-Jacobian pattern of \cite{SourceIntersections}. Near a
noncritical ordinary angle, $\kappa=0$ and the condition is simply that the target perturbation be
infinitesimal.

\section{Canonical complex trigonometry on a class-sized strip}\label{trigonometry:sec:strip}

\subsection{Real exponential and hyperbolic functions}
We use the established Gonshor exponential
$\exp_{\No}:(\No,+)\to(\No_{>0},\cdot)$ and its inverse $\log_{\No}$ \cite{Gonshor,vdDE}. They are
order-preserving group isomorphisms, agree with the ordinary real functions, and satisfy
\begin{equation}\label{trigonometry:eq:expreallocal}
 \exp_{\No}(x+h)=\exp_{\No}(x)\Ex(h)\qquad(h\in\mm_{\R}).
\end{equation}
For every $y\in\No$ put
$\sinh y=\frac12(\exp_{\No}y-\exp_{\No}(-y))$ and $\cosh y=\frac12(\exp_{\No}y+\exp_{\No}(-y))$;
the hyperbolic addition identities, $\cosh^2y-\sinh^2y=1$ and the usual derivatives follow from
the group law and \cref{trigonometry:eq:expreallocal}, and $\sinh y=0$ exactly at $y=0$. For
$x\ge1$ set $\arcosh x=\log_{\No}(x+\sqrt{x^2-1})$ and, for $|x|<1$,
$\artanh x=\frac12\log_{\No}\frac{1+x}{1-x}$; the real exponential identities show these are the
inverses of $\cosh$ on $\No_{\ge0}$ and of $\tanh$.

\subsection{Two complementary strips and an exponential quotient}
Define the additive subclasses
\[
 \mathscr H=\{x+iy:x\in\No,\ y\in\OO_{\R}\},\qquad
 \mathscr V=\{x+iy:x\in\OO_{\R},\ y\in\No\},
\]
so that $i\mathscr V=\mathscr H$. These are \emph{not} fixed-width ordinary strips. On the first
define
\begin{equation}\label{trigonometry:eq:stripexp}
 E(x+iy)=\exp_{\No}(x)\cis y .
\end{equation}

\begin{theorem}[Canonical strip exponential]\label{trigonometry:thm:stripexp}
The map $E:\mathscr H\to\SC^\times$ is a surjective group homomorphism with kernel $2\pi i\Z$. It
is fine-locally analytic and satisfies $E(z+h)=E(z)\Ex(h)$ for infinitesimal $h\in\SC$ and
$E'(z)=E(z)$. Consequently
\begin{equation}\label{trigonometry:eq:stripquotient}
 \mathscr V/(2\pi\Z)\longrightarrow\SC^\times,\qquad [z]\longmapsto E(iz),
\end{equation}
is a group isomorphism.
\end{theorem}
\begin{proof}
The group law follows from the real exponential and finite-angle addition. If $E(x+iy)=1$ its norm
forces $x=0$ and \cref{trigonometry:thm:polar} forces $y\in2\pi\Z$. For $w\ne0$ choose a finite
angle $\theta$ of $w$; then $E(\log_{\No}|w|+i\theta)=w$. The local formula follows by splitting an
infinitesimal increment into real and imaginary parts and using
\cref{trigonometry:eq:expreallocal}; this is an exact strongly summable local expansion, valid even
when $E(z)$ is infinite, since multiplication by a fixed surcomplex coefficient preserves strong
summability and one may choose a sufficiently small fine neighbourhood. Multiplication by $i$
proves \cref{trigonometry:eq:stripquotient}.
\end{proof}
The quotient is an algebraic statement, not a claim that this class strip is a classical
topological universal covering space.

\begin{definition}[Canonical strip sine and cosine]\label{trigonometry:def:stripsincos}
For $z\in\mathscr V$ set
\begin{equation}\label{trigonometry:eq:stripsincos}
 \Sin z=\frac{E(iz)-E(-iz)}{2i},\qquad \Cos z=\frac{E(iz)+E(-iz)}2 .
\end{equation}
When $z$ is a finite real surreal these are the lowercase functions already defined; when $z$ is
finite surcomplex they are the Taylor lifts of the ordinary entire complex sine and cosine.
\end{definition}

\begin{theorem}[Complex-variable identities and zero sets]\label{trigonometry:thm:striptrig}
For $z=x+iy\in\mathscr V$,
\begin{align}
 \Sin(x+iy)&=\sin x\cosh y+i\cos x\sinh y,&
 \Cos(x+iy)&=\cos x\cosh y-i\sin x\sinh y,\\
 |\Sin(x+iy)|^2&=\sin^2x+\sinh^2y,&
 |\Cos(x+iy)|^2&=\cos^2x+\sinh^2y.
 \label{trigonometry:eq:stripmoduli}
\end{align}
On the entire strip the addition identities, $\Sin^2z+\Cos^2z=1$, $\Sin'=\Cos$ and $\Cos'=-\Sin$
hold exactly, and $\Sin,\Cos$ are fine-locally analytic. The zeros of $\Sin$ are precisely $\pi\Z$
and those of $\Cos$ precisely $\pi/2+\pi\Z$; all are simple, and $\Tan=\Sin/\Cos$ has simple poles
at the latter points.
\end{theorem}
\begin{proof}
Insert \cref{trigonometry:eq:stripexp} into \cref{trigonometry:eq:stripsincos} and use the real
hyperbolic definitions for the first two formulas; squaring real and imaginary parts and using
$\cosh^2y=1+\sinh^2y$ gives the norm identities. All other identities follow algebraically from
the group law for $E$ and its local derivative. A zero in either norm identity forces
$\sinh y=0$, hence $y=0$ since the real exponential is injective; a sum of nonnegative squares in
$\No$ vanishes only if each term does. The kernel of $\cis$ then gives the listed zeros, whose
derivatives are $\pm1$, giving simplicity and the pole statement.
\end{proof}
In particular $\Sin(i\omega)=i\sinh\omega$ is canonically defined, even though this construction
prescribes no value for $\sin\omega$: \emph{infinite growth along an imaginary direction and
infinite oscillation along a real direction are different problems.}

\begin{theorem}[Surjectivity and complete fibres of strip sine]
\label{trigonometry:thm:sinefibers}
Every $w\in\SC$ is $\Sin z$ for some $z\in\mathscr V$. Modulo $2\pi$ it has exactly two preimages
counted with multiplicity, and the two coincide exactly when $w=1$ or $w=-1$. For
$z_1,z_2\in\mathscr V$,
\begin{equation}\label{trigonometry:eq:sinefibers}
 \Sin z_1=\Sin z_2
 \iff
 z_1-z_2\in2\pi\Z\ \text{ or }\ z_1+z_2\in\pi+2\pi\Z .
\end{equation}
Cosine is likewise surjective, since $\Cos z=\Sin(z+\pi/2)$.
\end{theorem}
\begin{proof}
Put $u=E(iz)$; the equation $\Sin z=w$ becomes the quadratic
\begin{equation}\label{trigonometry:eq:sinequadratic}
 u^2-2iwu-1=0 .
\end{equation}
The algebraically closed field $\SC$ supplies two roots counted with multiplicity, both nonzero
since their product is $-1$, and each has exactly one preimage modulo $2\pi$ by
\cref{trigonometry:thm:stripexp}. The discriminant is $4(1-w^2)$, which gives the multiplicity
assertion. Alternatively, subtracting the two sine values and using
$\Sin z_1-\Sin z_2=2\Cos\frac{z_1+z_2}2\Sin\frac{z_1-z_2}2$ together with the zero sets of
\cref{trigonometry:thm:striptrig} gives \cref{trigonometry:eq:sinefibers}.
\end{proof}

\begin{corollary}[Explicit inverse branches]\label{trigonometry:cor:sineinverse}
For $w\in\SC$ solve \cref{trigonometry:eq:sinequadratic}, $u=iw\pm\sqrt{1-w^2}$; neither root is
zero. For either root take a finite argument $x=\Arg(u)$ and put
\begin{equation}\label{trigonometry:eq:inversecomplexsine}
 z=x-i\log_{\No}|u| .
\end{equation}
Then $\Sin z=w$, and all solutions in $\mathscr V$ are $z+2\pi k$ and $\pi-z+2\pi k$, $k\in\Z$. For
example, for real $w>1$, including infinite $w$,
\[
 z=\frac\pi2\pm i\log_{\No}\bigl(w+\sqrt{w^2-1}\bigr)
\]
are the two solution classes.
\end{corollary}
\begin{proof}
Equation \cref{trigonometry:eq:inversecomplexsine} gives $E(iz)=u$ by
\cref{trigonometry:thm:stripexp}, and the second root of the quadratic is $-u^{-1}$, while
$E(i(\pi-z))=-E(iz)^{-1}$; \cref{trigonometry:thm:sinefibers} does the rest.
\end{proof}
Thus canonical complex trigonometry is confined neither to finite surcomplex values nor to finite
imaginary arguments, and no phase at an infinite \emph{real} argument is needed even to solve
$\Sin z=w$ for an arbitrary surcomplex target. This is a genuine extension of inverse complex
trigonometry to infinitely large values, with the choice of real exponential stated explicitly.
The unresolved phase choice is specifically at infinite real arguments, and
\cref{trigonometry:sec:global} is devoted to it.

\section{Trigonometric polynomials and sharp algebraic root counts}
\label{trigonometry:sec:trigpoly}

Let $n\in\N$ be ordinary and $c_{-n},\dots,c_n\in\SC$ arbitrary. Define on $\mathscr V$
\begin{equation}\label{trigonometry:eq:trigpoly}
 T(z)=\sum_{k=-n}^nc_kE(ikz).
\end{equation}
This is a \emph{finite} sum, so arbitrary surreal sizes of the coefficients cause no summability
problem; on finite real angles it is the trigonometric polynomial
$T(\theta)=\sum_kc_k\cis(k\theta)$, equivalently the Laurent polynomial $T(u)=\sum_kc_ku^k$ on
$\TT(\No)$. It is real-valued on finite real angles exactly when $c_{-k}=\bar c_k$ for all $k$.

\begin{theorem}[Algebraization and the $2n$ root bound]\label{trigonometry:thm:polyroots}
Suppose $T\ne0$ and set
\begin{equation}\label{trigonometry:eq:algebraization}
 P(U)=\sum_{k=-n}^nc_kU^{k+n}.
\end{equation}
Under $U=E(iz)$, the zeros of $T$ modulo $2\pi$ correspond bijectively, \emph{and with
multiplicity}, to the nonzero roots of $P$. Hence $T$ has at most $2n$ zeros modulo $2\pi$, counted
with multiplicity, on all of $\mathscr V$, and exactly $2n$ if $c_{-n}c_n\ne0$; in particular the
number of finite real angular zeros is at most $2n$, a bound attained by $\sin(n\theta)$, which has
the $2n$ distinct simple classes $\theta=k\pi/n$.
\end{theorem}
\begin{proof}
The identity $T(z)=U^{-n}P(U)$ and the bijection \cref{trigonometry:eq:stripquotient} give the
set-theoretic correspondence. Near a point with value $U_0\ne0$,
\[
 E(i(z+h))=U_0\Ex(ih)=U_0+iU_0h+O(h^2)
\]
has nonzero linear coefficient, and substitution by such a formal local coordinate preserves the
order of a zero; multiplication by the unit $U^{-n}$ also preserves it. Polynomial factorization in
$\SC$ gives the total number of nonzero roots as $\deg P-\operatorname{ord}_0P\le2n$, with equality
when both endpoint coefficients are nonzero.
\end{proof}
\emph{The degree bound is an ordinary finite integer even when the coefficients have arbitrary
infinite normal forms.} This is the definition of angular multiplicity used throughout: in a local
angular coordinate the derivative is a unit, so the angular multiplicity of a zero of $T$
\emph{is}, and equals, the polynomial multiplicity of $P$. There cannot be a nonzero finite
trigonometric polynomial vanishing at every ordinary real angle, since its associated polynomial
would have infinitely many distinct ordinary unit-circle roots; this also proves uniqueness of the
finite Fourier coefficients over $\SC$. Infinitesimal coefficients can split a multiple zero into a
cluster, as in \cref{trigonometry:thm:fold} and \cref{trigonometry:thm:cluster}, but cannot violate
the total-degree bound.

\begin{corollary}[Chebyshev and extremum equations]\label{trigonometry:cor:chebyshev}
For an ordinary $n\ge1$ and $w\in\SC$, the equation $\Cos(nz)=w$ has exactly $2n$ solutions modulo
$2\pi$ counted with multiplicity. For a nonconstant real trigonometric polynomial of degree at most
$n$, its derivative has at most $2n$ finite real zeros modulo $2\pi$.
\end{corollary}
\begin{proof}
The associated equation for the first assertion is $U^{2n}-2wU^n+1=0$, with nonzero constant and
leading coefficients. For the second, differentiation multiplies $c_k$ by $ik$, preserves the
frequency band, and does not annihilate a nonconstant polynomial in characteristic zero; apply
\cref{trigonometry:thm:polyroots}.
\end{proof}
\emph{The derivative statement bounds stationary angles; it does not invoke a general compactness
or extreme-value theorem on a surreal interval.} On all of $\No$, after a global extension has been
chosen in \cref{trigonometry:sec:global}, trigonometric polynomials repeat their zeros with the
larger period class of \cref{trigonometry:cor:globalzeros}, so a global count without a phase
quotient would be a different and generally infinite assertion.

\section{Support-controlled lifting of trigonometric zero clusters}
\label{trigonometry:sec:clusters}

This section specializes the positive-support philosophy of
\cite{SourceGeometry,SourceIntersections}. The proof is polynomial, so that the geometric theory
does not come to depend on their more substantial analytic division results. Let
\[
 T_0(z)=\sum_{k=-n}^na_kE(ikz),\qquad a_k\in\C,\quad a_{-n}a_n\ne0,
\]
be an ordinary trigonometric polynomial, perturbed to
$T(z)=\sum_k(a_k+e_k)E(ikz)$ with $e_k\in\mm_{\C}$, possibly zero, and let
$S\subset\Gamma_{>0}$ be the union of the supports of the nonzero $e_k$.

\begin{theorem}[Exact cluster multiplicity and coefficient support]\label{trigonometry:thm:cluster}
Suppose $\theta_0\in\C$ is an ordinary zero of $T_0$ of multiplicity $m$, and put
$u_0=e^{i\theta_0}$. Then $T$ has \emph{exactly} $m$ zeros, counted with multiplicity, in the
complex infinitesimal monad $\theta_0+\mm_{\C}$. Normalize the algebraized polynomial
\cref{trigonometry:eq:algebraization} to be monic; its unique monic factor $A$ collecting the roots
in $u_0+\mm_{\C}$ has degree $m$ and reduces to $(U-u_0)^m$, and the coefficients of
$A-(U-u_0)^m$, and of the complementary factor minus its ordinary reduction, are supported in
$S^*\setminus\{0\}$.

If $\theta_0$ is real, both $T_0$ and $T$ are real on finite real arguments, and $m=1$, then the
unique lifted zero is real. \emph{No corresponding reality conclusion is asserted for an arbitrary
multiple zero.}
\end{theorem}
\begin{proof}
Divide $P(U)=\sum(a_k+e_k)U^{k+n}$ by $a_n+e_n$; the inverse leading coefficient is an ordinary
nonzero scalar times a geometric series in an infinitesimal, so the normalized perturbation has
support in $S^*$ by \cref{trigonometry:lem:support}. Its ordinary reduction has the coprime monic
factorization $P_0=A_0B_0$ with $A_0=(U-u_0)^m$ and $B_0(u_0)\ne0$. Seek monic factors
$A=A_0+\Delta A$, $B=B_0+\Delta B$ of the same degrees with positive-support corrections. The
finite-dimensional ordinary linear map
\begin{equation}\label{trigonometry:eq:sylvester}
 (X,Y)\longmapsto XB_0+A_0Y,\qquad \deg X<m,\quad\deg Y<2n-m,
\end{equation}
is an isomorphism onto the polynomials of degree below $2n$: its kernel is zero by coprimality and
the degree restrictions, and the dimensions agree.

At a positive exponent $\nu\in S^*$ the coefficient of $AB=P$ is a linear equation with map
\cref{trigonometry:eq:sylvester}, whose right-hand side is the prescribed coefficient of $P-P_0$
minus the already determined coefficient of $\Delta A\,\Delta B$; in that product the two exponents
are positive and strictly below $\nu$. Transfinite recursion on the well-ordered monoid $S^*$
therefore constructs the factors, each equation involving only finitely many terms by
\cref{trigonometry:lem:support}, and the same recursion proves uniqueness. \emph{This is a
coefficient recursion, not a convergence claim for a Newton sequence}; by
\cref{trigonometry:rem:notlimit}(i) no valuation-improvement argument would be available.

Translate $U=u_0+H$. Then $A(u_0+H)$ is monic of degree $m$ with all lower coefficients
infinitesimal, and every one of its roots is infinitesimal: if a root $h$ were appreciable or
infinite, dividing its equation by $h^m$ would give $1$ plus a finite sum of infinitesimals equal
to zero. Its $m$ roots exist in an algebraically closed divisible Hahn workspace, and the
complementary factor cannot vanish at $u_0+h$ for infinitesimal $h$, its standard part there being
$B_0(u_0)\ne0$. The local bijection $U\mapsto\theta_0-i\Logm(U/u_0)$ transports these roots to the
specified angle monad, its nonzero linear term preserving multiplicity. Finally conjugation
preserves a real trigonometric equation and the monad of a real $\theta_0$; when there is exactly
one root, that root must be fixed by conjugation.
\end{proof}
\emph{The support bound belongs to the factor coefficients, not to individual roots.} A root may
require a fractional scale such as $t^{\lambda/2}$ even when the coefficient perturbation uses only
$t^\lambda$ --- which is exactly what the tangency of \cref{trigonometry:thm:tangency} exhibits.

\begin{example}[A double zero splitting off the real angle axis]\label{trigonometry:ex:doublesplit}
The equation $1-\Cos z=-\eps$, with $\eps>0$ infinitesimal, has two zeros in the monad of zero,
\[
 z=\pm i\arcosh(1+\eps),\qquad
 \arcosh(1+\eps)=\log_{\No}\bigl(1+\eps+\sqrt{2\eps+\eps^2}\bigr),
\]
since $\Cos(iy)=\cosh y$. They have valuation $\tfrac12v(\eps)$ and are \emph{not} real. This is
the complex continuation of the disappearance of real intersections after tangency in
\cref{trigonometry:thm:tangency}, and the strip counterpart of the negative-$\tau$ case of
\cref{trigonometry:thm:fold}.
\end{example}

\section{Finite Fourier calculus and positivity on the surreal circle}
\label{trigonometry:sec:fourier}

Nothing in this section concerns an uncontrolled infinite Fourier series. The number of frequencies
and of sample points is an ordinary finite integer, while the coefficients themselves may have
arbitrarily complicated set-sized surreal supports.

\subsection{Exact quadrature by roots of unity}
Let $M\in\N_{>0}$ be ordinary and $\zeta=\cis(2\pi/M)$, an ordinary primitive $M$th root of unity
by \cref{trigonometry:cor:representatives}. For every integer $k$,
\begin{equation}\label{trigonometry:eq:orthogonalityfinite}
 \frac1M\sum_{j=0}^{M-1}\zeta^{jk}=
 \begin{cases}1,&M\mid k,\\0,&M\nmid k,\end{cases}
\end{equation}
a finite geometric-sum identity, therefore unchanged over $\SC$.

\begin{theorem}[Finite Fourier inversion, sampling, and Parseval]\label{trigonometry:thm:fourier}
Let $T(u)=\sum_{k=-N}^Nc_ku^k$ with $c_k\in\SC$ and let $M>2N$. For \emph{any} base direction
$u_0\in\TT(\No)$,
\begin{align}
 c_k&=\frac1M\sum_{j=0}^{M-1}T(u_0\zeta^j)(u_0\zeta^j)^{-k},\qquad|k|\le N,
 \label{trigonometry:eq:fourierinverse}\\
 \frac1M\sum_{j=0}^{M-1}\bigl|T(u_0\zeta^j)\bigr|^2&=\sum_{k=-N}^N|c_k|^2 .
 \label{trigonometry:eq:parsevalsample}
\end{align}
More generally, for arbitrary $z_0,\dots,z_{M-1}\in\SC$ and
$\widehat z_k=\frac1M\sum_jz_j\zeta^{-jk}$,
\begin{equation}\label{trigonometry:eq:dft}
 z_j=\sum_{k=0}^{M-1}\widehat z_k\zeta^{jk},\qquad
 \sum_{j=0}^{M-1}|z_j|^2=M\sum_{k=0}^{M-1}|\widehat z_k|^2 .
\end{equation}
Thus $M$ exact samples determine every coefficient, even when the coefficients or the base
direction have infinitesimal parts.
\end{theorem}
\begin{proof}
Insert the finite Laurent sum into \cref{trigonometry:eq:fourierinverse}: distinct frequency
differences have absolute value at most $2N<M$, so \cref{trigonometry:eq:orthogonalityfinite} kills
every off-diagonal term. Multiplying $T$ by its conjugate and applying the same identity gives
\cref{trigonometry:eq:parsevalsample}. For \cref{trigonometry:eq:dft}, substitute the definition of
$\widehat z_k$, use the geometric-sum identity for inversion, and substitute the inversion formula
and its conjugate into $\sum_jz_j\bar z_j$. Every sum is finite and $\bar\zeta=\zeta^{-1}$, so no
support or convergence condition arises.
\end{proof}
All nonzero terms in the norm identity are positive surreals; \emph{Fourier cancellation cannot
destroy energy at a hidden infinitesimal order.} Exact coefficient recovery is quite different from
testing only the signs of the sample values: in \cref{trigonometry:ex:invisible} below, every
ordinary-angle sample is positive, yet the exact sample values still contain the infinitesimal
information, and the reconstructed polynomial exposes its negative windows.

\subsection{Coefficientwise integration and Parseval}
For an ordinary analytic $f$ with primitive $F$ on a real interval and finite endpoints
$a=a_0+\eps$, $b=b_0+\eta$, define the lifted integral by
\begin{align}
 \HInt_a^bf\sh(\theta)\,d\theta
 &=F\sh(b)-F\sh(a)\label{trigonometry:eq:endpointintegral}\\
 &=\int_{a_0}^{b_0}f(t)\,dt
 +\sum_{m\ge1}\frac{F^{(m)}(b_0)}{m!}\eta^m
 -\sum_{m\ge1}\frac{F^{(m)}(a_0)}{m!}\eps^m,\notag
\end{align}
the constant of the primitive cancelling. More generally a \emph{uniformly supported} family
$F(\theta)=\sum_\gamma t^\gamma f_\gamma(\theta)$ along an ordinary parameter interval --- that is,
an element of the ring (R1) of \cref{trigonometry:sec:rings}, with one well-ordered support and all
coefficient functions analytic on that one ordinary interval --- is integrated coefficientwise,
$\HInt_a^bF=\sum_\gamma t^\gamma\int_{a_0}^{b_0}f_\gamma$, with the same endpoint corrections. This
convention is not a definition of an unrestricted Riemann integral on all surreal functions, and
other surreal integration frameworks have substantially different scope and hypotheses \cite{CE}.
In particular
\begin{equation}\label{trigonometry:eq:orthogonality}
 \frac1{2\pi}\HInt_0^{2\pi}\cis(k\theta)\,d\theta=
 \begin{cases}1,&k=0,\\0,&k\ne0,\end{cases}\qquad k\in\Z,
\end{equation}
where the ordinary coefficient integral is the usual real or complex integral; its Hahn assembly is
not an ordinary Riemann limit over the full class interval.

\begin{theorem}[Finite Fourier identities over $\SC$]\label{trigonometry:thm:parseval}
If $F(\theta)=\sum_{k=-n}^nf_k\cis(k\theta)$ and $G(\theta)=\sum_{k=-m}^mg_k\cis(k\theta)$,
extending coefficients by zero where needed, then
\begin{align}
 f_j&=\frac1{2\pi}\HInt_0^{2\pi}F(\theta)\cis(-j\theta)\,d\theta,\\
 \frac1{2\pi}\HInt_0^{2\pi}F(\theta)\overline{G(\theta)}\,d\theta&=\sum_kf_k\bar g_k,
 \qquad
 \frac1{2\pi}\HInt_0^{2\pi}|F(\theta)|^2\,d\theta=\sum_k|f_k|^2,
 \label{trigonometry:eq:parseval}
\end{align}
the last expression being positive unless $F=0$. For a real trigonometric polynomial
$P(\theta)=a_0+\sum_{k=1}^n(a_k\cos k\theta+b_k\sin k\theta)$ the same convention gives
\begin{equation}\label{trigonometry:eq:continuousparseval}
 \frac1{2\pi}\HInt_0^{2\pi}P(\theta)^2\,d\theta=a_0^2+\frac12\sum_{k=1}^n(a_k^2+b_k^2).
\end{equation}
\end{theorem}
\begin{proof}
Expand the finite sums, use $\overline{\cis(k\theta)}=\cis(-k\theta)$, and apply
\cref{trigonometry:eq:orthogonality}; every operation on the frequency indices is finite.
Positivity follows from the ordering of $\No$ and positivity of each squared norm. Multiplication
by fixed surreal constants is legitimate because Hahn products have finite convolution at each
exponent.
\end{proof}
These identities remain valid for infinite and for infinitesimal coefficients. \emph{They do not
establish convergence or completeness for an unrestricted infinite Fourier series.} In particular,
infinitely many ordinary Fourier summands at the same Hahn exponent do not form a strongly summable
family merely because the corresponding real Fourier series converges.

\subsection{Dirichlet and Fej\'er kernels, and why they are not an approximation theorem}
Finite geometric sums and \cref{trigonometry:thm:identities} give
\begin{align}
 D_N(\theta)&=\sum_{k=-N}^N\cis(k\theta)=\frac{\sin((N+\tfrac12)\theta)}{\sin(\theta/2)},
 \label{trigonometry:eq:dirichlet}\\
 K_N(\theta)&=\frac1{N+1}\left|\sum_{j=0}^N\cis(j\theta)\right|^2
 =\frac1{N+1}\left(\frac{\sin((N+1)\theta/2)}{\sin(\theta/2)}\right)^2,
 \label{trigonometry:eq:fejerkernel}
\end{align}
the quotient formulas being used away from $\theta\in2\pi\Z$, where the finite sums give the
removable values $D_N=2N+1$ and $K_N=N+1$. In particular $K_N\ge0$ at every finite surreal angle,
with normalized mean $\mathcal M(K_N)=1$, where $\mathcal M(T)=c_0$ is the rotation-invariant
normalized mean, equal to the quadrature average whenever $M>N$ and, after expanding the finitely
many coefficients into normal form, to $\frac1{2\pi}\HInt_0^{2\pi}T$. These are exact algebraic
positivity and normalization statements. \emph{They are not, by themselves, a fine-topological
approximation theorem as $N$ grows.} For a fixed nonzero infinitesimal $\theta$, every ordinary
finite $N$ still has $(N+1)\theta$ infinitesimal, so $K_N(\theta)\sim N+1$, not the ordinary
fixed-nonzero-real-angle behaviour. One must specify a coefficient topology or another asymptotic
framework before formulating any infinite Fourier convergence theorem.

\subsection{Fej\'er--Riesz factorization with surcomplex coefficients}
The classical Fej\'er--Riesz theorem factors a nonnegative trigonometric polynomial as a squared
modulus; see \cite{DR} for its established context. \emph{The theorem below is a real-closed-field
extension of that classical result, not a claim of a new spectral-factorization theorem.} Its
relevance here is that the positivity hypothesis must range over surreal, not merely ordinary,
directions, and that coefficients and zeros may carry arbitrarily fine surreal scales while the
factorization stays finite and exact.

\begin{lemma}[Nonnegative polynomials over a real closed field]\label{trigonometry:lem:twosquares}
Let $F$ be real closed, $K=F[i]$, and $p\in F[t]$ nonnegative at every $t\in F$. Then there is
$q\in K[t]$ with $p(t)=q(t)\overline{q(t)}$ for all $t\in F$, and $\deg q=\tfrac12\deg p$ when
$p\ne0$.
\end{lemma}
\begin{proof}
For $p=0$ take $q=0$. Otherwise factor $p$ over $F$ into real linear factors and irreducible
quadratics. Every real root has even multiplicity: an odd multiplicity would change the sign across
that root, as is seen by taking a sufficiently small increment avoiding the finitely many other
roots. The leading constant is positive, and each monic irreducible quadratic is $(t-a)^2+b^2$ with
$b>0$. Take one factor $t-a-ib$ from each conjugate pair, half the multiplicity of each real linear
factor, and the positive square root of the leading constant; their product is $q$.
\end{proof}

\begin{theorem}[Surcomplex Fej\'er--Riesz theorem]\label{trigonometry:thm:fejer}
Let $T(\theta)=\sum_{k=-n}^nc_k\cis(k\theta)$ with $c_{-k}=\bar c_k$ and $c_k\in\SC$. The
following are equivalent.
\begin{enumerate}
\item $T(\theta)\ge0$ for \emph{every} finite surreal angle $\theta$, equivalently $T(u)\ge0$ for
 every $u\in\TT(\No)$.
\item There is $Q(U)=\sum_{j=0}^nq_jU^j\in\SC[U]$ with
 \begin{equation}\label{trigonometry:eq:fejer}
 T(\theta)=\bigl|Q(\cis\theta)\bigr|^2\qquad\text{for every finite surreal }\theta .
 \end{equation}
\end{enumerate}
If $T\ne0$, then $Q$ may be chosen with no zero in $|U|<1$; with that choice and the normalization
$Q(0)>0$ it is \emph{unique}.
\end{theorem}
\begin{proof}
(2)$\Rightarrow$(1) is immediate. Conversely, use the rational parametrization
\cref{trigonometry:eq:cayley} and form
\begin{equation}\label{trigonometry:eq:cayleypositive}
 p(t)=(1+t^2)^n\sum_{k=-n}^nc_kC(t)^k .
\end{equation}
Since $C(t)=(1+it)^2/(1+t^2)$ and $C(t)^{-1}=(1-it)^2/(1+t^2)$, the denominators cancel and $p$ is
a polynomial of degree at most $2n$, with real values and real coefficients by the Hermitian
symmetry of the $c_k$. The hypothesis and the surjectivity of the rational parametrization give
$p(t)\ge0$ for every $t\in\No$, so \cref{trigonometry:lem:twosquares} gives $p(t)=|q(t)|^2$ with
$q\in\SC[t]$ of degree at most $n$. Define
\begin{equation}\label{trigonometry:eq:fejerconstruction}
 Q(U)=\left(\frac{1+U}2\right)^nq\!\left(\frac{U-1}{i(U+1)}\right),
\end{equation}
a polynomial of degree at most $n$ by the degree bound on $q$. For $U=C(t)$ one has
$(1+U)/2=1/(1-it)$, hence $|Q(C(t))|^2=|q(t)|^2/(1+t^2)^n=T(C(t))$. This proves equality at every
unit point except possibly $U=-1$; on the circle, replacing conjugate powers by negative powers
makes the difference a Laurent polynomial which, after multiplication by $U^n$, is an
ordinary-degree polynomial with infinitely many roots, hence identically zero.

For the normalized form, let $d$ be the actual degree, so $c_d,c_{-d}\ne0$, and put $P(U)=U^dT(U)$,
of degree $2d$ with nonzero constant term and satisfying the self-reciprocal identity
\begin{equation}\label{trigonometry:eq:selfreciprocal}
 P(U)=U^{2d}\overline{P(1/\bar U)},
\end{equation}
where the right side means coefficient conjugation and reciprocal substitution. Algebraic
closedness pairs every root $\varsigma\ne0$ with $1/\bar\varsigma$, preserving multiplicity. A root
on the unit circle has even multiplicity: at a unit root $u_0$ use the nonsingular rational local
parameter $U=u_0C(t)$, under which the pullback of $T$ is a real rational function with nonvanishing
denominator; if its zero order were odd, its signs at sufficiently small positive and negative
surreal $t$ would be opposite --- factor out the first nonzero power of $t$ and bound the remaining
finite polynomial terms by the leading coefficient --- and one of those signs would be negative,
contradicting the hypothesis. \emph{The argument explicitly permits the test values $t$ to be
infinitesimal}, which is what makes it work over $\No$.

Now select, for each non-unit reciprocal pair, the root outside the unit disk with its multiplicity,
and half the multiplicity of every unit root: $d$ roots in all. Let $Q_0$ be their monic product.
The degree-$2d$ polynomials $P(U)$ and $U^dQ_0(U)\overline{Q_0(1/\bar U)}$ have the same roots with
the same multiplicities and nonzero constant terms, so they differ by a nonzero constant $\varkappa$;
on the circle this says $T=\varkappa|Q_0|^2$, and evaluating at a unit point away from the finite
root set, where $T>0$ and $|Q_0|^2>0$, gives $\varkappa\in\No_{>0}$. Set $Q=\sqrt\varkappa\,Q_0$
and multiply by a unit constant to make $Q(0)>0$. Any factor with no zeros inside the disk must
select exactly the same outside roots and half of each unit-root multiplicity, so two such factors
differ by a unit scalar, and their positive values at zero force that scalar to be one.
\end{proof}
The proof is finite algebra over a real closed field: no real compactness theorem, no Hilbert-space
completion, and no limit of spectral factors.

\begin{corollary}[Fourier coefficient bounds from positivity]\label{trigonometry:cor:coeffbounds}
Under the nonnegativity hypothesis of \cref{trigonometry:thm:fejer},
\[
 c_0=\sum_{j=0}^n|q_j|^2,\qquad |c_k|\le c_0\quad(-n\le k\le n),
\]
and $c_0>0$ when $T\ne0$. If $n\ge1$ the highest coefficient satisfies the sharper bound
$2|c_n|\le c_0$.
\end{corollary}
\begin{proof}
Expanding \cref{trigonometry:eq:fejer} gives $c_k=\sum_{j=0}^{n-k}q_{j+k}\bar q_j$ for $k\ge0$, and
finite-dimensional Cauchy--Schwarz --- proved over $\No$ by expanding a sum of squared norms ---
gives $|c_k|^2\le(\sum_j|q_{j+k}|^2)(\sum_j|q_j|^2)\le c_0^2$. For $k=n$, $c_n=q_n\bar q_0$ and
$2|q_n||q_0|\le|q_n|^2+|q_0|^2\le c_0$.
\end{proof}

\begin{example}[Ordinary-angle testing misses an infinitesimal negative window]
\label{trigonometry:ex:invisible}
Fix $0<\eps\in\mm_{\R}$ and put
\begin{equation}\label{trigonometry:eq:invisible}
 T(\theta)=(\sin\theta-\eps)^2-\eps^4,
\end{equation}
a real-valued trigonometric polynomial of degree two. It is strictly positive at every
\emph{ordinary} real $\theta$: if $\sin\theta\ne0$ is an ordinary real, the square is a positive
real leading term; if $\sin\theta=0$, the value is $\eps^2-\eps^4>0$. Nevertheless at the finite
surreal angle $\theta_*=\arcsin\eps$,
\[
 T(\theta_*)=-\eps^4<0,
\]
so no factorization \cref{trigonometry:eq:fejer} exists for this $T$. The negative window near zero
is exactly
\[
 \arcsin(\eps-\eps^2)<\theta<\arcsin(\eps+\eps^2),
\]
of width $2\eps^2(1+\vartheta)$ with $\vartheta\in\mm_{\R}$, centred to first order at an angle of
scale $\eps$; there is a second such window near $\pi-\arcsin\eps$. The four boundary angles are
simple zeros and attain the bound of \cref{trigonometry:thm:polyroots}. The same phenomenon occurs
with cosine: $(\cos\theta-\eps)^2-\eps^4$ is positive at every ordinary real angle and equals
$-\eps^4$ at $\theta=\arccos\eps$.

\emph{This is a concrete reason not to replace positivity on $\TT(\No)$ by positivity of a
real-angle trace, even at a fixed very small trigonometric degree. It does not contradict
\cref{trigonometry:prop:lift}:} the coefficients in \cref{trigonometry:eq:invisible} depend on a
non-real infinitesimal parameter, so that polynomial is not a single ordinary real-coefficient
analytic function. For contrast,
\[
 |u-(1+\eps)|^2=\eps^2+2(1+\eps)(1-\cos\theta)\qquad(u=\cis\theta)
\]
\emph{is} positive at every surreal direction, with exact minimum $\eps^2$ at $u=1$: a positive
trigonometric polynomial whose positivity margin is genuinely infinitesimal.
\end{example}

\section{Spherical trigonometry over \texorpdfstring{$\No$}{No}}\label{trigonometry:sec:spherical}

The distinction between algebraic geometry and angle evaluation extends beyond the plane. This
section and the next record two substantial analogues, proved from Gram identities, using only
three-dimensional vectors and finite matrix operations.

Let $P,Q,R\in\No^3$ be linearly independent unit vectors and define the minor central side angles
\[
 a=\arccos(Q\cdot R),\qquad b=\arccos(R\cdot P),\qquad c=\arccos(P\cdot Q),
\]
all in $(0,\pi)$ by \cref{trigonometry:thm:inverse}. At $P$ the unit tangent vectors toward $Q$ and
$R$ are
\[
 T_{PQ}=\frac{Q-\cos c\,P}{\sin c},\qquad T_{PR}=\frac{R-\cos b\,P}{\sin b},
\]
and $A\in(0,\pi)$ is their angle; define $B,C$ cyclically.

\begin{theorem}[Spherical cosine and sine laws]\label{trigonometry:thm:spherical}
With the above notation,
\begin{align}
 \cos a&=\cos b\cos c+\sin b\sin c\cos A,\label{trigonometry:eq:sphericalcos}\\
 \frac{\sin A}{\sin a}&=\frac{\sin B}{\sin b}=\frac{\sin C}{\sin c}
 =\frac{|\det(P,Q,R)|}{\sin a\sin b\sin c},\label{trigonometry:eq:sphericalsine}
\end{align}
together with the cyclic cosine laws. On a sphere of arbitrary positive surreal radius $\varrho$
the corresponding coherent minor arc lengths are $\varrho a$, $\varrho b$, $\varrho c$.
\end{theorem}
\begin{proof}
The tangent vectors are unit vectors perpendicular to $P$, and their dot product is
\[
 T_{PQ}\cdot T_{PR}=\frac{\cos a-\cos b\cos c}{\sin b\sin c},
\]
which is \cref{trigonometry:eq:sphericalcos}. In the tangent plane the magnitude of the determinant
of two unit tangent vectors is the sine of their angle; since $P$ is a unit normal,
\[
 \sin A=|\det(P,T_{PQ},T_{PR})|=\frac{|\det(P,Q,R)|}{\sin b\sin c},
\]
the cross-product identity used here being a polynomial consequence of the Euclidean Gram
determinant and therefore valid over $\No$. Dividing by $\sin a$ and relabelling gives
\cref{trigonometry:eq:sphericalsine}. Scaling the circle containing each great-circle arc gives the
arc-length statement from \cref{trigonometry:thm:arcs}.
\end{proof}
Linear independence excludes great-circle degeneracy and makes every sine in a denominator nonzero.
The theorem allows all three points to be infinitesimally close, so it includes spherical triangles
on arbitrarily small angular scales, and equally spherical triangles on a sphere of infinite
radius. \emph{It is a family of metric identities under explicit nondegeneracy hypotheses, not a
complete global geometric or measure-theoretic foundation for spherical geometry over $\No$.}


{\raggedright Spatial spherical trigonometry over $\No^3$ --- polar duality, reconstruction,
the exceptional side-side-angle family, branch-correct excess, area and
holonomy --- is developed in \path{docs/surreal/euclidean-three-space/}
(\texttt{e3:sec:spherical-laws} to \texttt{e3:sec:differential}); its
\texttt{e3:thm:phase} is the Hahn-workspace form of
\cref{trigonometry:thm:polar}.\par}
\section{Hyperbolic trigonometry in the full surreal disk}
\label{trigonometry:sec:hyperbolic}

\subsection{A disk metric including infinitesimal boundary distances}
Let $\Dd_{\No}=\{z\in\SC:|z|<1\}$. Unlike the ordinary-disk thickening used in some coherent
analytic statements, this disk includes points infinitesimally inside the unit circle whose
standard parts lie on its boundary. For $a\in\Dd_{\No}$ set
\begin{equation}\label{trigonometry:eq:diskmap}
 \varphi_a(z)=\frac{z-a}{1-\bar az},
\end{equation}
whose denominator is nonzero because $|\bar az|<1$. Direct expansion gives
\begin{equation}\label{trigonometry:eq:diskidentity}
 1-|\varphi_a(z)|^2=\frac{(1-|a|^2)(1-|z|^2)}{|1-\bar az|^2}>0,
\end{equation}
so $\varphi_a$ is a disk automorphism with inverse $(w+a)/(1+\bar aw)$; rotations are disk
automorphisms as well. Define
\begin{equation}\label{trigonometry:eq:diskdistance}
 \varrho(z,w)=\left|\frac{z-w}{1-\bar wz}\right|,\qquad
 d_H(z,w)=\log_{\No}\frac{1+\varrho(z,w)}{1-\varrho(z,w)}=2\artanh\varrho(z,w).
\end{equation}

\begin{proposition}[Disk invariance and metric properties]\label{trigonometry:prop:diskmetric}
The quantities $\varrho$ and $d_H$ are invariant under rotations and under the maps $\varphi_a$,
and $d_H$ is an $\No$-valued metric on $\Dd_{\No}$.
\end{proposition}
\begin{proof}
Invariance of $\varrho$ follows by substituting \cref{trigonometry:eq:diskmap}, putting differences
over common denominators and cancelling absolute values; invariance of $d_H$ follows. Symmetry,
positivity and separation of points follow from the definition and
\cref{trigonometry:eq:diskidentity}. For the triangle inequality move the middle point to zero: if
$|z|=r$, $|w|=s$ and the angle between them is $\theta$, then
\[
 \varrho(z,w)^2=\frac{r^2+s^2-2rs\cos\theta}{1+r^2s^2-2rs\cos\theta}
 \le\left(\frac{r+s}{1+rs}\right)^2,
\]
since for fixed $r,s$ the left-hand ratio, written as
$1-(1-r^2)(1-s^2)/(1+r^2s^2-2rs\cos\theta)$, is largest when $\cos\theta=-1$. Applying the
increasing logarithm and using
\[
 \frac{1+(r+s)/(1+rs)}{1-(r+s)/(1+rs)}=\frac{1+r}{1-r}\cdot\frac{1+s}{1-s}
\]
gives $d_H(z,w)\le d_H(z,0)+d_H(0,w)$.
\end{proof}
At a vertex $A$ the hyperbolic angle between the sides towards $B$ and $C$ is the Euclidean angle
between $\varphi_A(B)$ and $\varphi_A(C)$. This is invariant under disk automorphisms: after moving
the vertex to zero, the remaining automorphism fixing zero is a rotation, as direct algebra with
the fractional-linear formulas shows; equivalently the nonzero derivative
$\varphi_a'(z)=(1-|a|^2)/(1-\bar az)^2$ preserves angles. \emph{All these angles are finite, even
when their side distances are infinite.}

\subsection{Hyperbolic cosine and sine laws}
\begin{theorem}[Hyperbolic triangle laws over $\No$]\label{trigonometry:thm:hyperbolic}
Let $A,B,C$ be distinct points of $\Dd_{\No}$ not on one hyperbolic line, with
$a=d_H(B,C)$, $b=d_H(C,A)$, $c=d_H(A,B)$ and hyperbolic interior angles
$\alpha,\beta,\gamma\in(0,\pi)$. Then
\begin{align}
 \cosh a&=\cosh b\cosh c-\sinh b\sinh c\cos\alpha,\label{trigonometry:eq:hypcos}\\
 \frac{\sin\alpha}{\sinh a}&=\frac{\sin\beta}{\sinh b}=\frac{\sin\gamma}{\sinh c}.
 \label{trigonometry:eq:hypsine}
\end{align}
The distances $a,b,c$ may be infinite.
\end{theorem}
\begin{proof}
Move $A$ to zero and rotate $B$ to $r>0$, writing $C=s\cis\alpha$ after a possible orientation
reversal. Here $r=\tanh(c/2)$ and $s=\tanh(b/2)$, so $\cosh c=(1+r^2)/(1-r^2)$ and
$\sinh c=2r/(1-r^2)$, similarly for $b$. With $q=\varrho(B,C)$,
\[
 q^2=\frac{r^2+s^2-2rs\cos\alpha}{1+r^2s^2-2rs\cos\alpha},\qquad
 \cosh a=\frac{1+q^2}{1-q^2},
\]
and substitution yields
$\cosh a=\bigl((1+r^2)(1+s^2)-4rs\cos\alpha\bigr)/\bigl((1-r^2)(1-s^2)\bigr)$, which is
\cref{trigonometry:eq:hypcos}; by invariance the same calculation applies at every vertex. For the
sine law set $X=\cosh a$, $Y=\cosh b$, $Z=\cosh c$; the cosine law gives
\[
 \sin^2\alpha=\frac{1+2XYZ-X^2-Y^2-Z^2}{(Y^2-1)(Z^2-1)} .
\]
Dividing by $\sinh^2a=X^2-1$ produces an expression symmetric in $a,b,c$, so the squares of all
three ratios in \cref{trigonometry:eq:hypsine} agree; the ratios being strictly positive, the
unsquared equality follows.
\end{proof}

\begin{example}[An infinitely distant point inside the disk]\label{trigonometry:ex:infinitedistance}
For $0<\eps\in\mm_{\R}$ and $z=1-\eps$,
\[
 d_H(0,z)=\log_{\No}\frac{2-\eps}\eps=-\log_{\No}\eps+\log2+\log_{\No}(1-\eps/2),
\]
which is infinite, the last logarithm being a strongly summable infinitesimal series. \emph{No
value of sine or cosine at an infinite real circular angle is required to obtain either this
distance or the triangle laws.}
\end{example}

\begin{remark}[A second model, not restated]\label{trigonometry:rem:hyperboloid}
The same two laws can be proved in the Lorentz hyperboloid model
$\Hy^2_{\No}=\{P\in\No^3:L(P,P)=1,\ P_0>0\}$ with $L(P,Q)=P_0Q_0-P_1Q_1-P_2Q_2$ and
$a=\arcosh L(Q,R)$ etc., by the Lorentz Gram-determinant computation that proves
\cref{trigonometry:thm:spherical}: the tangent vectors $T_{PQ}=(Q-\cosh c\,P)/\sinh c$ have
$L$-norm $-1$, their $-L$ inner product is
$(\cosh b\cosh c-\cosh a)/(\sinh b\sinh c)$, and the Lorentz Gram matrix of $P,T_{PQ},T_{PR}$ has
determinant $\sin^2\alpha$, giving
$\sin\alpha=|\det(P,Q,R)|/(\sinh b\sinh c)$. We do not restate the theorem in that model: it is the
same pair of identities, and one model suffices. What the hyperboloid form makes visible is that
$L(P,Q)\ge1$ with equality only for $P=Q$ --- from
$L(P,Q)\ge\sqrt{1+p^2}\sqrt{1+q^2}-pq\ge1$ and
$(1+p^2)(1+q^2)-(1+pq)^2=(p-q)^2$ --- so the distances are genuinely positive. Unlike the circular
functions, the hyperbolic functions use the ordered real surreal exponential \emph{without}
discarding a purely infinite part, so infinite distances represent genuine large hyperbolic
separation while the angles between tangent directions remain finite.
\end{remark}

\section{Coherent arc length, sector area, and the limit pitfall}
\label{trigonometry:sec:arcs}

\begin{theorem}[Circle length and sector area at every radius]\label{trigonometry:thm:arcs}
For $O\in\SC$, $R\in\No_{>0}$ and finite $a\le b$, consider the unwrapped parametrized arc
$z(\theta)=O+R\cis\theta$, $a\le\theta\le b$. Its coherent arc length, computed from its speed with
\cref{trigonometry:eq:endpointintegral}, is $R(b-a)$, and its oriented sector area about $O$ is
$R^2(b-a)/2$. These count turns according to the chosen unwrapped interval; in particular one turn
has coherent circumference $2\pi R$ and the disk has area $\pi R^2$.
\end{theorem}
\begin{proof}
One has $z'(\theta)=iR\cis\theta$, so $|z'|=R$, and the integral of this constant is $R(b-a)$,
including the infinitesimal endpoint corrections. The usual parametrized sector functional has
integrand $\tfrac12\im(\overline{z-O}\,z')=R^2/2$, whose integral is $R^2(b-a)/2$. All expressions
are finite algebraic combinations of a fixed surreal scalar with scalar-lifted trigonometric
functions, so their coefficientwise integrals have valid common support in the ring (R1); the
chosen primitives are $R\theta$ and $R^2\theta/2$, and a change of their one global additive
constant cancels at the endpoints.
\end{proof}
As a related contour identity, direct parametrization gives
$\HInt_{|z-O|=R}\frac{dz}{z-O}=2\pi i$. \emph{This is a coherent contour functional, not an
assertion about ordinary paths being connected in the full fine topology}, and the word
``coherent'' is part of the definition: no Riemann-integral theory for every fine-continuous class
function is being claimed, nor any ordinary real-parameter fine-continuous traversal of a surreal
circle, which \cref{trigonometry:prop:topology} forbids. For an infinitesimal positive central
angle $\delta$, arc and chord are relatively equivalent, $2R\sin(\delta/2)\sim R\delta$, with
quadratic relative discrepancy, independently of whether $R$ is tiny or enormous; and for a simple
arc of angular span $\delta\in[0,\pi]$ the chord comparison reads
$\frac2\pi L\le\text{chord}\le L$.

\begin{proposition}[Failure of fine convergence of inscribed perimeters]
\label{trigonometry:prop:perimeters}
Let $R>0$ be any surreal radius and let $P_n=2nR\sin(\pi/n)$, $n=3,4,\dots$, be the perimeters of
the ordinary finite regular inscribed polygons of \cref{trigonometry:cor:polygon}. The sequence
does \emph{not} converge to $2\pi R$ in the fine topology, and the set $\{P_n:n\ge3\}$ has \emph{no
least upper bound} in $\No$.
\end{proposition}
\begin{proof}
The ordinary real numbers $q_n=P_n/R=2n\sin(\pi/n)$ satisfy $q_n<2\pi$ and converge to $2\pi$ in
$\R$. For any positive surreal infinitesimal $\eta$ the positive ordinary gap $2\pi-q_n$ exceeds
$\eta$ for every $n$, so $2\pi R-P_n>R\eta$ for every $n$, ruling out fine convergence. An infinite
upper bound for the $q_n$ cannot be least, because $2\pi$ is a smaller upper bound; a finite
candidate with standard part below $2\pi$ is not an upper bound by ordinary real convergence; a
candidate with standard part above $2\pi$ can be decreased; and every finite candidate with
standard part exactly $2\pi$ can be decreased by a positive infinitesimal and remain an upper
bound, each ordinary gap still exceeding any such correction. Multiplication by $R$ is an order
isomorphism and preserves the conclusion.
\end{proof}
Thus the circumference formula of \cref{trigonometry:thm:arcs} is a genuine coherent-integration
theorem. \emph{It must not be justified by importing a real supremum or a real sequence limit into
$\No$ without specifying the interpretation.} This example identifies exactly which classical
limiting step may not be transplanted.

\section{Infinite real phases: a canonical choice, its classification, and two obstructions}
\label{trigonometry:sec:global}

\subsection{Why the Maclaurin series is not the definition}
For a nonzero infinitely large $x\in\No$ one has $v(x)<0$, and the terms of the formal sine series
have valuations $v\bigl((-1)^nx^{2n+1}/(2n+1)!\bigr)=(2n+1)v(x)$, a strictly \emph{decreasing}
sequence. Their supports cannot have a well-ordered union, so that family is not strongly summable
and merely writing its formal expression defines nothing. At a finite argument with nonzero real
standard part the direct Maclaurin family also fails to be strongly summable, its ordinary terms
repeatedly contributing at exponent zero; this is exactly why
\cref{trigonometry:eq:sinfinite}--\cref{trigonometry:eq:cosfinite} evaluate the real centre first
in ordinary analysis and sum only the infinitesimal corrections. Failure of direct strong
summability at a real argument is not a failure of the canonical finite lift.

\subsection{The finite-part projection and the canonical normalization}
Every real surreal decomposes uniquely as
\begin{equation}\label{trigonometry:eq:split}
 x=p+r+\eps,\qquad p\in\Pi,\quad r\in\R,\quad\eps\in\mm_{\R},
\end{equation}
where $\Pi$ consists of zero and the normal forms supported entirely at negative $t$-exponents,
whose nonzero elements are purely infinite. Put $\fin(x)=r+\eps$, an \emph{additive} projection
$\No\twoheadrightarrow\OO_{\R}$ with kernel $\Pi$, so that $\No=\Pi\oplus\OO_{\R}$; the finite part
retains both the ordinary constant and all infinitesimal terms, and is not merely the standard
part. It is \emph{not} multiplicative:
$\fin(\omega\cdot\omega^{-1})=1$ while $\fin(\omega)\fin(\omega^{-1})=0$. Reducing phases is thus
compatible with adding angles, not with reducing the factors of their products.

The omnific integers have the normal-form description $\Oz=\Pi\oplus\Z$. They form a discrete
subring: sums and products of purely infinite normal forms have no positive exponents and no
constant term, and every positive element of $\Oz$ is at least $1$. They form an integer part of
$\No$: writing $x=p+f$ with $f$ finite and choosing the ordinary integer $n$ with $n\le f<n+1$
(taking the appropriate side when $f$ is infinitesimally adjacent to an integer) gives
$p+n\le x<p+n+1$. Since $\Pi$ is closed under multiplication and division by nonzero ordinary real
constants,
\begin{equation}\label{trigonometry:eq:periodclass}
 2\pi\Oz=\Pi+2\pi\Z,\qquad \pi\Oz=\Pi+\pi\Z .
\end{equation}
Ehrlich and Kaplan's additional characterization is that $\Oz$ is the unique integer part of $\No$
which is initial in its simplicity tree \cite[Lemma 11.1]{EK}. \emph{We use that published
characterization to explain the word canonical; we do not need to reproduce its tree-theoretic
proof.}

\begin{definition}[Canonical real-surreal sine and cosine]\label{trigonometry:def:globaltrig}
For every $x\in\No$ set
\[
 \Sin x=\sin(\fin x),\qquad \Cos x=\cos(\fin x),\qquad \Phi(x)=\Cos x+i\Sin x .
\]
\end{definition}
This is exactly the integer-part normalization of \cite[\S11]{EK}: reducing $x$ modulo $2\pi\Oz$
leaves the same finite representative modulo $2\pi\Z$ as reducing $\fin x$. \emph{The infinitesimal
part is retained}, so for $p\in\Pi$ and infinitesimal $\eps$ one has $\Sin p=0$, $\Cos p=1$ and
$\Sin(p+\eps)=\eps-\eps^3/6+\cdots$; the operation is not standard part. The notation is consistent
with \cref{trigonometry:def:stripsincos}, with which it agrees wherever both apply.

\begin{theorem}[Canonical global exponential and its period class]
\label{trigonometry:thm:globalexp}
Define $\Exp(x+iy)=\exp_{\No}(x)\Phi(y)$. Then $\Exp:(\SC,+)\to(\SC^\times,\cdot)$ is onto, extends
the ordinary complex exponential, and satisfies
\begin{equation}\label{trigonometry:eq:expkernel}
\begin{gathered}
 \ker\Exp=2\pi i\,\Oz,\qquad |\Exp z|=\exp_{\No}(\re z),\qquad
 \Exp'=\Exp,\\
 \Exp(z+h)=\Exp(z)\Ex(h)\quad(h\in\mm_{\SC}).
\end{gathered}
\end{equation}
Putting $\Sin z=(\Exp(iz)-\Exp(-iz))/(2i)$ and $\Cos z=(\Exp(iz)+\Exp(-iz))/2$ extends
\cref{trigonometry:def:globaltrig} to all of $\SC$, satisfies Euler's formula
$\Exp(iz)=\Cos z+i\Sin z$, the addition identities, $\Sin^2z+\Cos^2z=1$, $\Sin'=\Cos$ and
$\Cos'=-\Sin$; and for real surreal $x,y$,
\begin{equation}\label{trigonometry:eq:trighyper}
 \Sin(x+iy)=\Sin x\cosh y+i\Cos x\sinh y,\qquad
 \Cos(x+iy)=\Cos x\cosh y-i\Sin x\sinh y .
\end{equation}
Moreover $(\Cos z+i\Sin z)^n=\Cos(nz)+i\Sin(nz)$ for every ordinary integer $n$, whence
$\Cos(nz)=T_n(\Cos z)$ and $\Sin((n+1)z)=\Sin z\,U_n(\Cos z)$ as polynomial identities, with no
division problem at a sine zero.
\end{theorem}
\begin{proof}
Both $\fin$ and the finite phase map are additive-to-multiplicative in the appropriate sense, so
$\Exp$ is a group homomorphism, and its modulus is immediate. For $w\ne0$ write $w=\rho\cis\theta$
with $\theta$ finite (\cref{trigonometry:thm:polar}); then $w=\Exp(\log_{\No}\rho+i\theta)$. If
$\Exp(x+iy)=1$ the modulus forces $x=0$ and \cref{trigonometry:thm:polar} forces
$\fin(y)\in2\pi\Z$, so \cref{trigonometry:eq:split} and \cref{trigonometry:eq:periodclass} give the
kernel. For infinitesimal $h=a+ib$ one has $\fin(y+b)=\fin(y)+b$, so the real exponential Taylor
law \cref{trigonometry:eq:expreallocal} and the finite phase law give
$\Exp(z+h)=\Exp(z)\Ex(a)\Ex(ib)=\Exp(z)\Ex(h)$; this is a local surcomplex power series, proving
fine-local analyticity and the derivative. Substitution gives the remaining identities, and
expanding products by the exponential group law gives the addition and Pythagorean identities.
\end{proof}
This calculation identifies the map with the canonical surcomplex exponential of
\cite[Theorem 11.1]{EK}. \emph{The local power-series argument establishes the fine-local analytic
property needed here without asserting uniqueness from a differential equation.}

\begin{corollary}[Zero classes and periods]\label{trigonometry:cor:globalzeros}
On the full surcomplex plane,
\[
 \{z:\Sin z=0\}=\pi\Oz,\qquad \{z:\Cos z=0\}=\pi/2+\pi\Oz .
\]
All these zeros are real surreals and are simple as fine-analytic zeros. The period class of
$\Sin$, of $\Cos$ and of the pair is $2\pi\Oz$. Hence $2\pi$ is the least positive period, but
\emph{the period class is not the cyclic group $2\pi\Z$}: it contains every purely infinite real
surreal.
\end{corollary}
\begin{proof}
$\Sin z=0$ says $\Exp(2iz)=1$, so $2iz\in2\pi i\Oz$; for cosine use $\Cos z=\Sin(z+\pi/2)$. At a
sine zero $\Cos^2z=1$, so the derivative is nonzero. A common period is exactly an element $T$ with
$\Exp(iT)=1$, hence $T\in2\pi\Oz$; if $T$ is a period of sine alone, evaluating at $0$ and $\pi/2$
gives $\Sin T=0$ and $\Cos T=1$, so it is a common period, and similarly for cosine.
\end{proof}
This is compatible with every finite-angle geometric theorem above, whose angle ambiguity is still
only $2\pi\Z$. It is also the first appearance of the phenomenon that
\cref{trigonometry:thm:infiniteperiods} shows to be unavoidable.

\subsection{Classification of all circular group-law extensions}
\begin{theorem}[All extensions of the finite phase]\label{trigonometry:thm:characters}
For every group homomorphism $\chi:(\Pi,+)\to(\TT(\No),\cdot)$ the formula
\begin{equation}\label{trigonometry:eq:globalcis}
 \Psi_\chi(p+f)=\chi(p)\cis f,\qquad p\in\Pi,\quad f\in\OO_{\R},
\end{equation}
defines a group homomorphism $\No\to\TT(\No)$ extending canonical finite-angle $\cis$, and every
such extension is uniquely of this form. Writing $\Psi_\chi=C_\chi+iS_\chi$ one obtains real-valued
global sine and cosine satisfying the addition formulas, $C_\chi^2+S_\chi^2=1$ and, in fine-local
calculus, $S_\chi'=C_\chi$, $C_\chi'=-S_\chi$; for every $h\in\mm_{\R}$,
\begin{equation}\label{trigonometry:eq:globallocal}
 \Psi_\chi(x+h)=\Psi_\chi(x)\Ex(ih).
\end{equation}
Each $E_\chi(x+iy)=\exp_{\No}(x)\Psi_\chi(y)$ is then an onto exponential homomorphism with
$E_\chi'=E_\chi$ and $|E_\chi(z)|=\exp_{\No}(\re z)$, and the corresponding global sine and cosine
agree with the canonical strip functions of \cref{trigonometry:sec:strip}. The canonical choice
$\chi=1$ is exactly $\Psi_1=\Phi$ of \cref{trigonometry:def:globaltrig}.
\end{theorem}
\begin{proof}
The direct-sum decomposition $\No=\Pi\oplus\OO_{\R}$ and the two group laws prove that
\cref{trigonometry:eq:globalcis} is a homomorphism with the desired restriction. Conversely the
restriction of any extension to $\Pi$ determines $\chi$, and its value at $p+f$ must be the product
of its two restrictions. An infinitesimal increment belongs to the finite summand, which gives
\cref{trigonometry:eq:globallocal}; expanding its right-hand side proves fine-local analyticity and
the derivative identities. The proof of \cref{trigonometry:thm:stripexp} applies verbatim to
$E_\chi$.
\end{proof}

\begin{example}[Any prescribed phase at $\omega$]\label{trigonometry:ex:phases}
Let $\ell(p)\in\R$ be the coefficient of $\omega=t^{-1}$ in a purely infinite normal form. For any
finite angle $\tau\in\OO_{\R}$ the rule $\chi_\tau(p)=\cis(\tau\ell(p))$ is a character, because
$\ell$ is additive and $\tau\ell(p)$ is finite, and its extension satisfies
$\Psi_{\chi_\tau}(\omega)=\cis\tau$. Thus $S(\omega)$ and $C(\omega)$ can jointly be assigned
\emph{any} point of the surreal unit circle without changing the finite-angle theory, the addition
laws, or the fine-local differential equations: phases $0$, $\pi/2$, $\pi$ give respectively
$(C(\omega),S(\omega))=(1,0),(0,1),(-1,0)$. The corresponding complex family
$E_{\chi_\tau}$, with $E_{\chi_0}(i\omega)=1$ and $E_{\chi_\pi}(i\omega)=-1$, is the family already
constructed in \cite{SourceAnalysis}; the character formulation exposes the entire freedom.
\end{example}
The trivial character gives the explicit normal-form extension $\Psi_1(x)=\cis(\fin x)$. \emph{It
is a possible convention, not a consequence of the local trigonometric axioms alone}: the
classification above concerns group homomorphisms with the specified finite restriction, and is not
a classification of all solutions of a differential equation. In particular it does not by itself
answer the separate robustness questions formulated in \cite[\S11]{EK} for more general initial
fields and their integer parts. Sources 20 and 21 use it to answer one specific possibility in the
first of them negatively --- initiality of the kernel-multiplier ring does not force the canonical
exponential (\cref{trigonometry:per:thm:robustness}) --- and those questions remain open otherwise
(\cref{trigonometry:per:sec:ek}). Conversely, the existence of these alternatives does not negate the canonical
status of the initial-integer-part construction: canonicity there refers to that additional
structure, not to uniqueness under the weaker analytic hypotheses. The two statements are
complementary, and both are kept.

\begin{theorem}[Infinite periods are unavoidable]\label{trigonometry:thm:infiniteperiods}
Every homomorphic extension in \cref{trigonometry:thm:characters} has a kernel element in every
coset $p+\OO_{\R}$ of the finite subgroup; in particular, for every nonzero purely infinite $p$ it
has an infinite period. \emph{No extension agreeing with canonical finite-angle $\cis$ can have
kernel exactly $2\pi\Z$}, and the full surreal additive quotient $\No/2\pi\Z$ cannot be identified
with $\TT(\No)$ by any extension of the finite phase.
\end{theorem}
\begin{proof}
Finite-angle $\cis$ is already surjective onto the whole unit circle
(\cref{trigonometry:thm:polar}), so choose a finite $a$ with $\cis a=\chi(p)^{-1}$; then
$\Psi_\chi(p+a)=1$. For nonzero purely infinite $p$, the element $p+a$ remains infinitely large in
modulus. Being in the kernel of a group homomorphism makes it a period of the phase and of both
coordinate functions.
\end{proof}
This is stronger than a bare example of nonuniqueness, and stronger than the observation that one
particular convention happens to create extra zeros: the ordinary statement ``all periods are
ordinary multiples of $2\pi$'' is \emph{incompatible} with extending the circular group law to all
surreal real inputs while retaining its finite-angle meaning. The canonical choice of
\cref{trigonometry:cor:globalzeros}, whose period class is $2\pi\Oz$, is the instance of this
phenomenon that the integer-part normalization produces; the theorem says no other choice escapes
it. Each $E_\chi$ agrees with the canonical strip functions of \cref{trigonometry:sec:strip}, so
all disagreement occurs strictly beyond that strip, and the real-period and zero sets outside it
are correspondingly larger and choice-dependent. \Cref{trigonometry:per:thm:cofinal} strengthens the
theorem: every such extension has cofinally large periods all of whose rational multiples are periods.

\subsection{Two coherence obstructions, both over the fixed-domain ring}
Recall from \cref{trigonometry:sec:rings} the ring (R1): a \emph{common-domain Hahn-analytic}
datum on an ordinary neighbourhood $J$ of a point is a family $H=\sum_{\gamma\in S}h_\gamma
t^\gamma$ with \emph{one} well-ordered support $S$ and ordinary real-analytic coefficients
$h_\gamma$ all defined on the \emph{same} $J$; evaluation at finite thickened points uses Taylor
substitution, and ordinary holomorphic coefficients on a complex neighbourhood may be used instead.
This is the real restriction of the coherence convention of \cite{SourceAnalysis}, and it is
strictly stronger than independently prescribing local power series on different monads. Both
theorems in this subsection are statements about (R1)--(R2) \emph{and about no other coefficient
ring}; neither is asserted over the radius-free ring (R3) or the formal ring (R4).

\begin{theorem}[No coherent infinite-frequency bounded sine]
\label{trigonometry:thm:infinitefrequency}
Let $S_\chi$ be any global extension from \cref{trigonometry:thm:characters} and let $L>0$ be
infinite. There is \emph{no} common-domain Hahn-analytic $H$ over the ring (R1) on an ordinary
neighbourhood of zero whose evaluation satisfies $H\sh(x)=S_\chi(Lx)$ throughout its finite surreal
thickening.
\end{theorem}
\begin{proof}
For every ordinary real $r$ in the neighbourhood the right-hand side lies in $[-1,1]$, so every
negative-exponent coefficient $h_\gamma(r)$ of $H(r)$ is zero by uniqueness of normal form and
finiteness. This holds for all such $r$, so every negative-exponent coefficient function vanishes
identically, and $H'(0)=\sum_\gamma h_\gamma'(0)t^\gamma$ has no negative exponents and is finite.
But equality on the thickening implies equality of the fine derivatives at zero, and by
\cref{trigonometry:eq:globallocal}, or by the chain rule with the constant scalar $L$,
\[
 \left.\frac d{dx}S_\chi(Lx)\right|_{x=0}=L,
\]
which is infinite. This contradiction proves the theorem.
\end{proof}
\emph{The requirement that the coefficient functions share an ordinary neighbourhood is essential},
which is to say: the theorem is about (R1)--(R2) and would be a different statement over any other
ring. On a sufficiently small surreal scale every $S_\chi(Lx)$ still has its usual local power
series; what fails is the compatibility of that infinite-frequency behaviour with a bounded,
common-domain ordinary analytic coefficient description.

\begin{theorem}[Why all-scale coherence forces a polynomial]\label{trigonometry:thm:allscale}
Suppose a global function $f$ had, at \emph{every} positive surreal radius $\varrho$, a
Hahn-coherent chart over the ring (R1) for $w\mapsto f(\varrho w)$ on the infinitesimal thickening
of one fixed ordinary neighbourhood of zero. Then $f$ has only finitely many nonzero Taylor
coefficients at zero. In particular neither the global exponential nor any global sine $S_\chi$
satisfies this condition.
\end{theorem}
\begin{proof}
Write $a_n=f^{(n)}(0)/n!$. The support of every $a_n\varrho^n$ is contained in the chart's single
well-ordered support, because ordinary differentiation introduces no new Hahn exponents. If
infinitely many $a_n$ are nonzero, choose a positive surreal $\mathit\Lambda$ larger than four
times every $|v(a_n)|$ and take $\varrho=t^{-\mathit\Lambda}$. For successive nonzero indices
$m>n$,
\[
 v(a_m\varrho^m)-v(a_n\varrho^n)=v(a_m)-v(a_n)-(m-n)\mathit\Lambda<0,
\]
producing an infinite decreasing sequence inside the alleged chart support, a contradiction.
\end{proof}
The stronger polynomial classification, including global propagation of that finite germ, is proved
in \cite{SourceAnalysis}; the displayed obstruction already suffices here, and \emph{this is a
restriction on uniform analytic descriptions over (R1), not a failure of the algebraic
trigonometric identities or of the arbitrary-radius circle geometry.} A circle with enormous radius
uses the map $R\cis\theta$ with a \emph{finite} angle; it does not require $w\mapsto\Sin(Rw)$ to
possess one Hahn support over an entire ordinary angular domain. It is also why the all-scale
coherence category, in which entire functions are polynomials, is not the category being asserted
for these global sine functions.

\section{The rotation group \texorpdfstring{$\SO(2,\No)$}{SO(2,No)}}\label{trigonometry:rot:sec:main}

This section is the contribution of source~19, \emph{Rotations of the Surreal Plane: The Norm-One
Torus, Infinitesimal Angles, and Topology at Every Scale} (September 22, 2026). It organizes the
finite-angle theory of \cref{trigonometry:sec:circle} around the group
\begin{equation}\label{trigonometry:rot:eq:defG}
 \SO(2,\No)=\{A\in M_2(\No):A^{\mathsf T}A=I,\ \det A=1\},
\end{equation}
whose structure theorem
\begin{equation}\label{trigonometry:rot:eq:main}
 \SO(2,\No)\ \cong\ \TT(\No)\ \cong\ \OO_{\R}/2\pi\Z\ \cong\ \TT(\R)\times(\mm_{\R},+)
\end{equation}
is exactly \cref{trigonometry:prop:rotation} together with \cref{trigonometry:thm:polar}, and is
printed only there. In this section the word \emph{rotation} always means the algebraic isometry;
continuity is an additional property in a specified topology, never part of an informal picture of
turning a rigid object through ordinary real time. The entries of a rotation matrix lie in
$[-1,1]$, but there is a proper class of infinitesimal corrections to them: bounded coordinates do
not make the group set-sized, and they do not make it compact in any ordinary sense.

\subsection{What source 19 re-derives, and the conventions of this section}
\label{trigonometry:rot:sec:overlap}
Source~19 was written against repository commit \texttt{dcf86662b574}, at which the present
article already stood unchanged, and it inspected only the opening of this article. It states
that the finite-angle quotient, the split real/infinitesimal phase, the projective half-angle
coordinates and the distinction between global normalization and arbitrary extensions were already
here, and presents itself as a group-centred synthesis rather than a discovery of these. The overlap
is larger than it records: the results in the following table were also already here, and each is
printed once, at the place indicated, with both sources credited. The only genuinely different proof
among them is kept as a marked second route (\cref{trigonometry:rot:rem:rootsroute}). For the
evaluation of formal series at infinitesimals, source 19 cites the surreal analytic use of Neumann's
lemma in \cite[\S3]{BMgerms}; here that role is played by \cref{trigonometry:lem:support}.
\begin{center}\small
\begin{longtable}{L{0.47\textwidth}L{0.45\textwidth}}
\toprule
Source 19 statement & Printed as\\
\midrule\endhead
$\TT(F)\cong\SO(2,F)$ by $a+ib\mapsto\left(\begin{smallmatrix}a&-b\\b&a\end{smallmatrix}\right)$; the group is abelian & \cref{trigonometry:prop:rotation}\\
Structure theorem; logarithmic linearization $\TT(\No)\cong\TT(\R)\times\mm_{\R}$; finite-angle uniformization with kernel $2\pi\Z$; polar decomposition & \cref{trigonometry:thm:polar}, \cref{trigonometry:eq:anglegroup,trigonometry:eq:circlesplit}; same argument (conjugation makes the logarithm purely imaginary), not a second route\\
Infinitesimal exponential and logarithm, inverse and addition laws, formal evaluation on infinitesimals, leading-term control & \cref{trigonometry:eq:explog,trigonometry:eq:logleading}, \cref{trigonometry:lem:support,trigonometry:lem:leading}\\
$\cis(r+\eps)=e^{ir}\Ex(i\eps)$ and its addition law & \cref{trigonometry:eq:finiteEuler}, \cref{trigonometry:thm:identities}\\
Principal angle in $(-\pi,\pi]$ with endpoint correction; a branch of $\Arg$ is not additive & \cref{trigonometry:cor:representatives}, \cref{trigonometry:rem:cut}\\
Rational chart $C$, its inverse and its homogeneous group law & \cref{trigonometry:thm:cayley}, in this article's homogeneous convention (below)\\
Angular distance; chord $=2\sin(\text{angle}/2)$ and its series; equal valuations of angular and chordal distance & \cref{trigonometry:thm:metric}\\
All $n$th roots; no infinitesimal torsion & \cref{trigonometry:cor:representatives}; second route in \cref{trigonometry:rot:rem:rootsroute}\\
$v(\sin\eps)=v(\eps)$, $v(1-\cos\eps)=2v(\eps)$; $v(u-v)=v(\ellm(u)-\ellm(v))$ in one monad & \cref{trigonometry:prop:leading}; \cref{trigonometry:eq:localangle}\\
Displacement $|uz-vz|=|z||u-v|$ at any radius & \cref{trigonometry:eq:rotdisplacement}; its valuation form is \cref{trigonometry:rot:eq:dispval}\\
Set-sized subsets closed and discrete; convergent set-indexed nets eventually constant; no nonconstant ordinary real paths & \cref{trigonometry:prop:topology}; the Cauchy-net clause is added in \cref{trigonometry:rot:sec:topology}\\
$\fin$, $\Oz=\Pi\oplus\Z$, the normalized global phase, kernel $2\pi\Oz$, the Ehrlich--Kaplan exponential & \cref{trigonometry:def:globaltrig}, \cref{trigonometry:thm:globalexp}, \cref{trigonometry:cor:globalzeros}\\
All extensions of the finite phase; explicit extensions with any value at $\omega$; unavoidable infinite periods & \cref{trigonometry:thm:characters}, \cref{trigonometry:ex:phases} (a finite $\tau$ there covers source 19's real parameter), \cref{trigonometry:thm:infiniteperiods} (its coset form is source 19's ``for every infinite $x$ a finite $\theta$ with $x-\theta$ in the kernel'')\\
$(\cis)'=i\cis$; every global phase obeys the same local law; Taylor data on each monad select $\cis$ & \cref{trigonometry:thm:identities}, \cref{trigonometry:eq:globallocal}, \cref{trigonometry:prop:normalization}\\
A derivation on $\No$ is not a fine derivative & the closing paragraph of \cref{trigonometry:sec:foundations}\\
\bottomrule
\end{longtable}
\end{center}

\emph{Notation.} Source 19's symbols are converted once, here, and not used again. Its $\OO,\mm$ are
$\OO_{\R},\mm_{\R}$. Its $\Exp,\operatorname{Log}$ on infinitesimals are $\Ex,\Logm$, because $\Exp$ in this report
is the global exponential of \cref{trigonometry:thm:globalexp}. Its $H_0$, $H_\chi$ and
purely infinite summand $\mathcal J$ are $\Phi$, $\Psi_\chi$ and $\Pi$; its angular distance
$d_{\mathrm{ang}}$ is $\angle$; its coefficient $a_\omega(p)$ is the $\ell(p)$ of
\cref{trigonometry:ex:phases}. Its infinitesimal angle $\ell(u)$ is written $\ellm(u)$ below, since
the bare letter $\ell$ is taken by that coefficient; its identity monad $U$ is $\Tm$, since $U$ is
used for Chebyshev polynomials and for $E(iz)$; its rotation matrix $R(a,b)$ is $\Rot(a,b)$, since
$R$ is a circumradius; its generator $J$ is $\mathsf J$, since $J$ is an interval in (R1); its
conjugation $\kappa$ is $\mathsf c$, since $\kappa=v(f'(a))$ in \cref{trigonometry:thm:stability};
its dual-number variable $\delta$, double-angle map $p_2$, Clifford bivector $E$, invariant
differential $\alpha$, rescaling automorphism $F_\lambda$ and circles $S_r$ are $\mathsf d$,
$\mathrm{sq}$, $e_{12}$, $\varpi$, $\Xi_a$ and $\mathcal S_\rho$. Source 19 writes the homogeneous
chart as $[s:t]\mapsto C(t/s)$ with law $[s:t]*[p:q]=[sp-tq:sq+tp]$; in this article's convention
$[p:q]\mapsto C(p/q)$ this is exactly \cref{trigonometry:eq:projectiveaddition}. Throughout this
section $F$ denotes a real closed field and $L=F[i]$; in the topological subsections $F$ is a
\emph{set} field, and in the algebraic ones the field $\No$ itself may be taken, all arguments being
finite algebra.

For $u\in\TT(\No)$ with $u_0=\st(u)$ put
\begin{equation}\label{trigonometry:rot:eq:ellm}
 \ellm(u)=-i\Logm(u/u_0)\in\mm_{\R},\qquad u=u_0\Ex\bigl(i\ellm(u)\bigr),\qquad
 \Tm=\{u\in\TT(\No):\st(u)=1\}.
\end{equation}
By \cref{trigonometry:thm:polar}, $\ellm$ is additive, $\Tm\to(\mm_{\R},+)$ is an isomorphism, and
\begin{equation}\label{trigonometry:rot:eq:exact}
 1\longrightarrow\Tm\longrightarrow\TT(\No)\xrightarrow{\ \st\ }\TT(\R)\longrightarrow1
\end{equation}
is split by the literal inclusion of the ordinary unit circle. The splitting uses the embedded real
field, coefficientwise standard part and formal Taylor evaluation; once these are fixed no basis or
arbitrary section is chosen, and the infinitesimal angle of a unit is recovered exactly, not up to
an infinitesimal error.

\emph{Size convention.} $\No$ is a proper class, so statements about $\SO(2,\No)$ are class-group
statements, readable in a class theory supporting the surreal construction or relative to a
universe. No proper-class Hahn sum is formed, no quantifier ranges over a set of all class
functions, and the group is never treated as a set-sized manifold. Topological statements about the
full class are made with ball neighbourhoods and explicit quantifiers, not with a set of all open
subclasses.

\begin{remark}[Infinitesimal angles are not a surreal vector space]\label{trigonometry:rot:rem:notvector}
$\mm_{\R}$ is an $\OO_{\R}$-module and an $\R$-vector space, but not an $\No$-vector space:
$\omega^{-1}\in\mm_{\R}$ while $\omega\cdot\omega^{-1}=1\notin\mm_{\R}$. Accordingly $\Tm$ is
torsion-free, uniquely divisible, and an $\R$-vector group through $r\odot\Ex(i\eps)=\Ex(ir\eps)$,
$r\in\R$; the same formula does not extend to all surreal scalars while staying in $\Tm$. That the
algebraic Lie algebra of \cref{trigonometry:rot:sec:lie} is one-dimensional over $\No$ is a different
statement, about tangent vectors.
\end{remark}

\subsection{Isometries fixing the origin}\label{trigonometry:rot:sec:isometry}
\Cref{trigonometry:prop:rotation} identifies the orientation-preserving \emph{linear} isometries.
Linearity need not be assumed.

\begin{theorem}[Isometry rigidity]\label{trigonometry:rot:thm:isometry}
Let $F$ be a real closed field, $\|x\|=\sqrt{x_1^2+x_2^2}$ on $F^2$, and let $f:F^2\to F^2$ satisfy
$f(0)=0$ and $\|f(x)-f(y)\|=\|x-y\|$ for all $x,y$. Then $f$ is $F$-linear and orthogonal. Neither
continuity nor surjectivity is assumed. The maps of this kind with positive orientation are exactly
$\SO(2,F)$.
\end{theorem}
\begin{proof}
Polarization $2\langle x,y\rangle=\|x\|^2+\|y\|^2-\|x-y\|^2$ shows that $f$ preserves inner
products. Put $u=f(e_1)$, $w=f(e_2)$; their Gram matrix is the identity, so they form an orthonormal
basis. The inner products of $f(x)$ with $u$ and $w$ are $x_1$ and $x_2$, hence $f(x)=x_1u+x_2w$.
This is linear and orthogonal, and an orthogonal matrix has determinant $\pm1$.
\end{proof}
Preserving only $\|f(x)\|=\|x\|$ would not suffice: circles of different radii can be rotated by
different units, destroying pairwise distances. A field automorphism that changes surreal distances
is not an isometry in this exact surreal-valued sense.

Every origin-fixing isometry is therefore uniquely $z\mapsto uz$ or $z\mapsto u\bar z$ with
$u\in\TT(F)$. With $\mathsf c(z)=\bar z$ one has $\mathsf c\,u\,\mathsf c=u^{-1}$, so
\begin{equation}\label{trigonometry:rot:eq:orthogonal}
 \Orth(2,F)\cong\TT(F)\rtimes\{1,\mathsf c\},
\end{equation}
the nontrivial element acting by inversion, and subtracting the image of the origin gives the
Euclidean group $F^2\rtimes\Orth(2,F)$. \emph{These are algebraic semidirect products, not
assertions about connected components in an unspecified topology.}

\subsection{The rational chart: exceptional cases and the logarithmic bridge}
\label{trigonometry:rot:sec:chart}
In the affine coordinate of \cref{trigonometry:thm:cayley} the exceptional cases of the law
$t\oplus s=(t+s)/(1-ts)$ are
\[
 t\oplus s=\infty\ (ts=1),\qquad \infty\oplus t=-1/t\ (t\ne0),\qquad
 \infty\oplus0=\infty,\qquad \infty\oplus\infty=0;
\]
the identity is $0$ and the inverse is $t\mapsto-t$, with $\infty$ fixed. For instance $C(1)=i$,
and $1\oplus1=\infty$ records $i^2=-1$. In homogeneous coordinates the inverse of $[p:q]$ is
$[-p:q]$, and $[p:q]=[r:s]$ exactly when $ps-qr=0$. Near the missing half-turn,
\begin{equation}\label{trigonometry:rot:eq:halfturn}
 C(t)+1=\frac{2(1+it)}{1+t^2}:
\end{equation}
an infinite chart parameter is a finite angle infinitesimally close to $\pi$, not an infinite
angle, which is one reason the projective formula is more robust than an unqualified use of
$\tan(\theta/2)$.

\begin{remark}[A parametrization of points, not an isomorphism of varieties]
\label{trigonometry:rot:rem:projective}
\Cref{trigonometry:thm:cayley} concerns $F$-points and semialgebraic maps. It is \emph{not} an
isomorphism of the complete variety $\PP^1_F$ with the affine algebraic torus: over $F[i]$ the
denominator $p^2+q^2$ vanishes at nonzero pairs, and identifying the two statements would erase
exactly the complex points at which the rational formula has poles.
\end{remark}

For infinitesimal $t\in\mm_{\R}$ the chart value lies in $\Tm$, and
\begin{equation}\label{trigonometry:rot:eq:arctanbridge}
 \ellm\bigl(C(t)\bigr)=-i\Logm C(t)=2\arctan t=2\sum_{n\ge0}\frac{(-1)^nt^{2n+1}}{2n+1},
\end{equation}
because $\Logm C(t)=\Logm(1+it)-\Logm(1-it)$ and the even powers cancel. This is a direct bridge
between exact rational arithmetic and the infinitesimal angle, consistent with $t=\tan(\theta/2)$
in \cref{trigonometry:thm:cayley} and with \cref{trigonometry:eq:smallseries}. On infinitesimal
parameters the law $t\oplus s$ stays infinitesimal and is linearized by $2\arctan$; no branch issue
arises there.

\subsection{Powers, roots, torsion and finite subgroups}\label{trigonometry:rot:sec:torsion}

\begin{remark}[Second route to all roots, via the splitting]\label{trigonometry:rot:rem:rootsroute}
Write $u=u_0\Ex(i\eps)$ with $u_0=e^{ir}\in\TT(\R)$ and $\eps=\ellm(u)$. For every ordinary integer
$n$ the direct-product law gives $u^n=u_0^n\Ex(in\eps)$. For ordinary $n\ge1$ the solutions of
$v^n=u$ are exactly
\begin{equation}\label{trigonometry:rot:eq:roots}
 v_k=e^{i(r+2\pi k)/n}\,\Ex(i\eps/n),\qquad 0\le k<n .
\end{equation}
Indeed, writing $v=d\,\Ex(i\eta)$, uniqueness of the split factors in
\cref{trigonometry:eq:circlesplit} gives $d^n=u_0$ and $n\eta=\eps$; the first equation has the
displayed $n$ ordinary solutions and the second exactly one, and the $v_k$ are distinct because
their standard parts are. This is source 19's proof. The count in
\cref{trigonometry:cor:representatives} uses instead the kernel of $\cis$ and the degree bound for
$X^n-u$; the route here uses neither, only the splitting and ordinary complex roots.
\end{remark}

\begin{corollary}[No new finite-order rotations]\label{trigonometry:rot:cor:torsion}
The torsion subgroup of $\TT(\No)$ consists exactly of the ordinary complex roots of unity and is
isomorphic to $\Q/\Z$. Every finite subgroup of $\TT(\No)\cong\SO(2,\No)$ is the literal standard
group $\mu_n=\{e^{2\pi ik/n}:0\le k<n\}$ for one ordinary $n\ge1$.
\end{corollary}
\begin{proof}
If $u^n=1$ then $n\ellm(u)=0$, so $\ellm(u)=0$ and $u$ is an ordinary root of unity
(\cref{trigonometry:cor:representatives}). A finite subgroup has a finite exponent $N$, hence lies
in the cyclic group $\mu_N$, and is $\mu_n$ for a divisor $n$ of $N$.
\end{proof}
Thus infinitesimal rotations produce new elements of \emph{infinite} order, not new regular polygons
with a finite nonstandard symmetry order; ordinary finite order is the only meaning of torsion used.

\begin{corollary}[Divisibility]\label{trigonometry:rot:cor:divisible}
$\SO(2,\No)$ is divisible, has no nontrivial homomorphism to a finite group, and has no proper
subgroup of finite index.
\end{corollary}
\begin{proof}
Divisibility is \cref{trigonometry:rot:eq:roots}. If $f$ maps onto a finite group of order $N$,
write $u=v^N$; then $f(u)=f(v)^N=1$. A finite-index subgroup of an abelian group gives a finite
quotient.
\end{proof}

\emph{Square roots without transcendental evaluation.} For $u=a+ib\ne-1$ in $\TT(F)$ the unique
square root with positive real part is
\begin{equation}\label{trigonometry:rot:eq:sqrt}
 w=\sqrt{\frac{1+a}2}+i\,\frac{b}{2\sqrt{(1+a)/2}},
\end{equation}
since $b^2=(1-a)(1+a)$ gives $w^2=u$; the other root is $-w$, and the roots of $-1$ are $\pm i$. Only
ordered-field square roots are used. This branch suits a half-angle near the identity and is
\emph{not} a globally multiplicative choice: no multiplicative section of $u\mapsto u^2$ exists on
the whole circle, since it would send $-1$, of order two, to an element whose square is $-1$ but
whose order divides two. The obstruction is elementary and persists over the surreals.

\emph{The double-angle map and $\operatorname{Spin}(2)$.} The map
$\mathrm{sq}:\TT(F)\to\TT(F)$, $a+ib\mapsto(a^2-b^2)+2iab$, is onto with kernel $\{\pm1\}$. In the
Clifford algebra with $e_1^2=e_2^2=-1$ put $e_{12}=e_1e_2$, so $e_{12}^2=-1$; then
\begin{equation}\label{trigonometry:rot:eq:spin}
 (a+be_{12})\,e_1\,(a-be_{12})=(a^2-b^2)e_1+2ab\,e_2 ,
\end{equation}
with the corresponding expression for $e_2$, so conjugation by the spinor $a+be_{12}$, $a^2+b^2=1$,
induces the rotation $\mathrm{sq}(a+ib)$. This is a finite algebraic calculation and does not depend
on ordinary path lifting in the fine topology.

\begin{proposition}[Finite orthogonal subgroups]\label{trigonometry:rot:prop:dihedral}
A finite subgroup of $\Orth(2,\No)$ is either some $\mu_n$ or conjugate, by a rotation, to the
standard dihedral group $\mu_n\rtimes\{1,\mathsf c\}$.
\end{proposition}
\begin{proof}
Its rotation subgroup is some $\mu_n$ by \cref{trigonometry:rot:cor:torsion}. If it contains a
reflection $z\mapsto u\bar z$, choose $v$ with $v^2=u$. Conjugating by multiplication by $v$ sends
$z\mapsto u\bar z$ to $z\mapsto v^{-1}u\bar v\bar z=\bar z$ and leaves all rotations unchanged, so the
group becomes $\mu_n\cup\mu_n\mathsf c$.
\end{proof}
The conjugacy is stronger than saying the abstract finite groups have familiar names, and weaker
than saying all their matrices are literally real: before conjugation a reflection axis may have a
genuinely nonstandard direction.

\subsection{Valuation scales and precision}\label{trigonometry:rot:sec:scales}
For $\gamma>0$ in the value group put
$\Tm^{\ge\gamma}=\{u\in\Tm:v(u-1)\ge\gamma\}$ and $\Tm^{>\gamma}=\{u\in\Tm:v(u-1)>\gamma\}$. Since
$v(\Ex(i\eps)-1)=v(\eps)$ (\cref{trigonometry:eq:logleading}), these are subgroups, identified by
$\ellm$ with the valuation subgroups of $\mm_{\R}$.

\begin{theorem}[One real angular coordinate at each scale]\label{trigonometry:rot:thm:graded}
For every surreal $\gamma>0$ there is a canonical isomorphism
\begin{equation}\label{trigonometry:rot:eq:graded}
 \Tm^{\ge\gamma}/\Tm^{>\gamma}\ \cong\ (\R,+),
\end{equation}
sending the class of $u$ to the coefficient of $t^\gamma$ in $\ellm(u)$.
\end{theorem}
\begin{proof}
On angles of valuation at least $\gamma$, extraction of the coefficient at $t^\gamma$ is additive,
onto (by the angles $rt^\gamma$), and has kernel the angles of valuation greater than $\gamma$.
Transport through $\ellm$.
\end{proof}
This is a scale \emph{filtration}, not a decomposition of $\Tm$ as a direct sum of one copy of $\R$
per surreal exponent: individual angles can have infinite well-ordered supports, and neither an
unrestricted product nor the finite-support vectors captures the normal-form condition. Measuring the
cosine defect alone loses first-order information, since $v(1-\cos\eps)=2v(\eps)$
(\cref{trigonometry:prop:leading}): a cosine error of scale $t^{2\gamma}$ can encode an angular change
of scale $t^\gamma$.

\emph{Displacement at any radius.} For $u,v\in\TT(\No)$ and $z\in\SC$ one has exactly
$|uz-vz|=|z||u-v|$, the same amplification in every direction, with no supremum over vectors. Hence
for $z\ne0$ and a nonzero infinitesimal angle $\eps$, by \cref{trigonometry:eq:rotdisplacement},
\begin{equation}\label{trigonometry:rot:eq:dispval}
 v\bigl(\Ex(i\eps)z-z\bigr)=v(z)+v(\eps),
\end{equation}
so an infinitesimal angle at an infinite radius can produce an infinitesimal, appreciable or infinite
displacement according to the sign of $v(z)+v(\eps)$. Here and below $x\sim y$ is the algebraic
relation of \cref{trigonometry:sec:conventions}, not a sequence limit.

\begin{example}[A finite displacement from an infinitesimal angle]\label{trigonometry:rot:ex:cayleydisp}
Take $t=\omega^{-1}$ and $u=C(t)$. The rational formula gives exactly
\[
 u\,\omega=\omega-\frac{2\omega^{-1}}{1+\omega^{-2}}+i\,\frac2{1+\omega^{-2}} .
\]
The angle $\ellm(u)=2\arctan\omega^{-1}$ is infinitesimal, yet $u\omega-\omega$ is finite with
standard part $2i$: the longitudinal change is infinitesimal and the transverse displacement
appreciable. Likewise the single angle $\eps=\omega^{-2}$ displaces the points $\omega$, $\omega^2$,
$\omega^3$ by amounts $\sim\omega^{-1}$, $\sim1$, $\sim\omega$, while the rotation matrix, the same in
all three cases, is infinitesimally close to the identity.
\end{example}

\emph{Precision under products, powers and roots.} If $u_j=c_j\Ex(i\eps_j)$ and approximations
$\widehat u_j=c_j\Ex(i\widehat\eps_j)$ have the same standard phases, then for an ordinary finite
number of factors
\begin{equation}\label{trigonometry:rot:eq:producterror}
 \ellm\Bigl(\prod_ju_j\Big/\prod_j\widehat u_j\Bigr)=\sum_j(\eps_j-\widehat\eps_j),
\end{equation}
whose valuation is at least the least individual error valuation and may be larger through
cancellation. Multiplication by a fixed ordinary nonzero integer, or division by it when a fixed root
branch is selected, preserves a nonzero infinitesimal valuation, so ordinary-degree powers and
matched root branches neither lose nor gain valuation precision; a mismatched standard root branch
causes an appreciable error.

\emph{Hahn products.} If a set-indexed family $(\eps_j)_{j\in I}$ of infinitesimal angles is
strongly summable, its sum is infinitesimal and
\begin{equation}\label{trigonometry:rot:eq:hahnproduct}
 \prod_{j\in I}^{\mathrm{Hahn}}\Ex(i\eps_j)=\Ex\Bigl(i\sum_{j\in I}\eps_j\Bigr).
\end{equation}
To justify the notation, expand the left side over finite-support choices of powers of the factors.
Each exponent of a resulting term is a finite word in the well-ordered positive union of the
supports, which has only finitely many representations by \cref{trigonometry:lem:support}, and each
letter comes from only finitely many members of the family by strong summability; so each exponent
receives finitely many terms, each coefficient identity reduces to a finite exponential identity, and
the value is independent of ordering. For example, with $t=\omega^{-1}$,
$\prod_{n\ge1}^{\mathrm{Hahn}}\Ex(it^n)=\Ex\bigl(it/(1-t)\bigr)$. \emph{The finite partial products
do not converge to this value in the fine topology} (\cref{trigonometry:rot:sec:topology}), and
repeating one fixed nonzero infinitesimal angle infinitely often is not a strongly summable family,
so \cref{trigonometry:rot:eq:hahnproduct} assigns no canonical value to an ``infinite number of
identical turns.''

\emph{Surreal exponents of a rotation.} For an ordinary integer $n$, $u^n$ is a group operation. For
a nonordinary surreal $s$ the group axioms supply nothing analogous. On $\Tm$ the unique logarithm
gives the partial definition
\begin{equation}\label{trigonometry:rot:eq:surrealpower}
 \bigl(\Ex(i\eps)\bigr)^s:=\cis(s\eps)\qquad\text{whenever }s\eps\in\OO_{\R},
\end{equation}
an analytic scaling of a uniquely chosen infinitesimal angle, \emph{not} a transfinite iteration of
the group operation. Beyond that domain a global phase extension (\cref{trigonometry:thm:characters})
supplies additional choices; for general $u$ even the choice of a finite argument is a branch, and
changing it by $2\pi n$ can change the result after multiplication by $s$.

\subsection{The action on the plane and its polynomial invariants}\label{trigonometry:rot:sec:action}
For $\rho>0$ let $\mathcal S_\rho=\{z\in\SC:|z|=\rho\}$. The action of $\SO(2,\No)$ on
$\mathcal S_\rho$ is free and transitive: the unique rotation taking $z$ to $w$ is multiplication by
$w/z$, and $uz=z$ with $z\ne0$ forces $u=1$. The orbits are the origin and one circle for each
positive surreal radius. The absence of nonzero fixed points of a nonidentity rotation is an
algebraic fact and needs no continuous motion from the identity. The products $\langle z,w\rangle$
and $[z,w]$ of \cref{trigonometry:eq:dotcross} are invariant under simultaneous rotation, and the
cyclic order of directions is defined by ordering the two semicircles with determinant signs, or by
finite arguments modulo $2\pi$ --- not by ordinary real paths in the fine topology.

\begin{proposition}[Polynomial invariants]\label{trigonometry:rot:prop:invariants}
For a real closed field $F$, the polynomial invariants of the algebraic rotation group acting on one
plane vector are
\begin{equation}\label{trigonometry:rot:eq:invariants}
 F[X,Y]^{\SO(2)}=F[X^2+Y^2],
\end{equation}
and the same holds for invariance under all $F$-valued rotations, in particular with surreal
coefficients.
\end{proposition}
\begin{proof}
Over $L=F[i]$ put $Z=X+iY$, $W=X-iY$; a unit $u$ acts by $Z\mapsto uZ$, $W\mapsto u^{-1}W$, so
$Z^aW^b$ has weight $a-b$. Independence of distinct Laurent powers of $u$ forces an invariant to
contain only terms with $a=b$, so it lies in $L[ZW]=L[X^2+Y^2]$, and conjugation fixes exactly
$F[X^2+Y^2]$. For invariance under $F$-points, the rational chart supplies infinitely many $u$, and a
Laurent polynomial in one variable vanishing at infinitely many points is zero.
\end{proof}
\emph{This is a statement about polynomial functions}, not a classification of all class functions
on orbits; arbitrary functions of the radius are nonpolynomial invariants.

\subsection{The norm-one torus and its Lie algebra}\label{trigonometry:rot:sec:lie}
The coordinate ring of the rotation group over $F$ is
$A_F=F[\mathsf a,\mathsf b]/(\mathsf a^2+\mathsf b^2-1)$. The gradient $(2\mathsf a,2\mathsf b)$
vanishes nowhere on the circle, so the group is smooth of dimension one. Over $L$ put
$\mathsf z=\mathsf a+i\mathsf b$, with inverse $\mathsf a-i\mathsf b$; then
$\mathsf a=(\mathsf z+\mathsf z^{-1})/2$, $\mathsf b=(\mathsf z-\mathsf z^{-1})/(2i)$ and
\begin{equation}\label{trigonometry:rot:eq:torussplit}
 A_F\otimes_FL\cong L[\mathsf z,\mathsf z^{-1}].
\end{equation}
So the rotation group is the norm-one torus $\ker(\operatorname{Res}_{L/F}\Gm\to\Gm)$, split over
$L$ and anisotropic over $F$: conjugation acts on its character lattice $\Z$ by $n\mapsto-n$, and
there is no nontrivial $F$-defined character to $\Gm$. These are the standard algebraic-group
interpretations \cite[Chapters 10 and 12]{Milne}. For $\No$ this fixed algebraic group is applied at
surreal points, rather than treating the class $\No$ as a set-sized coefficient field of scheme
theory; each finite construction lives in a real closed Hahn workspace.

\emph{Tangent vectors.} With dual numbers $F[\mathsf d]/(\mathsf d^2)$, a tangent vector $I+\mathsf dX$
satisfies $(I+\mathsf dX)^{\mathsf T}(I+\mathsf dX)=I+\mathsf d(X^{\mathsf T}+X)$, so $X$ is
skew-symmetric, hence traceless, and the determinant adds no condition:
\begin{equation}\label{trigonometry:rot:eq:lie}
 \mathfrak{so}(2,F)=F\mathsf J,\qquad
 \mathsf J=\begin{pmatrix}0&-1\\1&0\end{pmatrix},\qquad [a\mathsf J,b\mathsf J]=0 .
\end{equation}
This is the precise sense in which the surreal tangent space is one-dimensional over $\No$. From
$\mathsf J^2=-I$,
\begin{equation}\label{trigonometry:rot:eq:formalexp}
 \exp(T\mathsf J)=I\sum_{n\ge0}\frac{(-1)^nT^{2n}}{(2n)!}+\mathsf J\sum_{n\ge0}\frac{(-1)^nT^{2n+1}}{(2n+1)!},
\end{equation}
Hahn-evaluable at infinitesimal $T$ and equal to $\Rot(\cos T,\sin T)$ at finite $T$ by the
real-plus-infinitesimal prescription, where
$\Rot(a,b)=\left(\begin{smallmatrix}a&-b\\b&a\end{smallmatrix}\right)$. At an arbitrary infinite
$T$ the displayed series is not strongly summable, its $\mathsf J$-component being the sine family
of \cref{trigonometry:sec:global}: the Lie algebra has all
surreal scalars, but the canonical finite-angle exponential is defined only on $\OO_{\R}\mathsf J$.

\begin{proposition}[No algebraic additive exponential]\label{trigonometry:rot:prop:noalgexp}
Every morphism of algebraic groups $\Ga\to\SO(2)$ over a real closed field is trivial.
\end{proposition}
\begin{proof}
Over $L$ the target is $\Gm$, and a regular map $\mathbb A^1_L\to\Gm$ is a unit of $L[T]$, a
constant, which must be $1$. Triviality after scalar extension gives triviality over $F$.
\end{proof}
The analytic phase is therefore not a missing algebraic morphism: even over $\R$ the exponential is
transcendental, and over the surreals its domain and summation convention must also be specified.

\emph{The invariant differential.} On the circle $\varpi=\mathsf a\,d\mathsf b-\mathsf b\,d\mathsf a$
is invariant, $\mathsf a\,d\mathsf a+\mathsf b\,d\mathsf b=0$ by differentiating the equation, and
the rational chart pulls it back to
\begin{equation}\label{trigonometry:rot:eq:cayleyform}
 C^*\varpi=\frac{2\,dt}{1+t^2},
\end{equation}
locally $d(2\arctan t)=d\theta$. This is an identity of algebraic differential forms whose local
primitive is evaluated by the infinitesimal rule; \emph{it does not by itself define integration over
arbitrary full-class contours.} The tangent space, the algebraic dimension and this form survive
even though the fine topology has no nonconstant ordinary real paths; a connected-component argument
cannot replace the algebraic tangent calculation.

\subsection{Algebraic representations and nonalgebraic symmetries}
\label{trigonometry:rot:sec:representations}

\begin{theorem}[Algebraic endomorphisms]\label{trigonometry:rot:thm:endomorphisms}
For the norm-one torus over a real closed field $F$,
\begin{equation}\label{trigonometry:rot:eq:endomorphisms}
 \End_{F,\mathrm{alg}}(\SO(2))\cong\Z,\qquad n\longleftrightarrow(u\mapsto u^n),
\end{equation}
and the algebraic automorphisms are exactly the identity and inversion.
\end{theorem}
\begin{proof}
Over $L$ a regular homomorphism $\Gm\to\Gm$ is an invertible Laurent polynomial $c\mathsf z^n$, and
the identity value forces $c=1$. Each $\mathsf z^n$ commutes with the descent involution
$\mathsf z\mapsto\mathsf z^{-1}$, so it is defined over $F$. An inverse would have exponent $m$ with
$mn=1$.
\end{proof}
For $n\ne0$ the power maps are isogenies of degree $|n|$, whose surreal fibres
\cref{trigonometry:rot:eq:roots} describes; they multiply the infinitesimal angle exactly by $n$.

Over $L$ a finite-dimensional algebraic representation of $\Gm$ is a finite direct sum of
integer-weight spaces: write its coaction $v\mapsto\sum_nv_n\otimes\mathsf z^n$ (finitely many terms
for each $v$); coassociativity makes each $v_n$ a weight-$n$ vector, the counit axiom gives
$v=\sum_nv_n$, and independence of Laurent monomials makes the sum direct. Over $F$ conjugation pairs
weight $n$ with weight $-n$, so the irreducible algebraic representations are the trivial one and,
for each ordinary $n\ge1$, the two-dimensional
\begin{equation}\label{trigonometry:rot:eq:weightrep}
 u\longmapsto\Rot\bigl(\re(u^n),\im(u^n)\bigr),
\end{equation}
and every finite-dimensional algebraic representation is a direct sum of these. This is the rank-one
case of torus representation theory \cite[\S12h]{Milne}; the weight calculation shows why \emph{no
new weight labels occur: the weights are ordinary integers}, not omnific integers or surreal
frequencies. Finite linear combinations of these matrix coefficients are algebraic trigonometric
polynomials (\cref{trigonometry:thm:polyroots}). \emph{Nothing here supplies a Fourier convergence
theorem, a Haar measure on the full class, or a classification of arbitrary abstract
representations.}

For each nonzero ordinary real $a$ define
\begin{equation}\label{trigonometry:rot:eq:nonalgauto}
 \Xi_a\bigl(u_0\Ex(i\eps)\bigr)=u_0\Ex(ia\eps),\qquad u_0\in\TT(\R),\ \eps\in\mm_{\R}.
\end{equation}
It is a group automorphism with inverse $\Xi_{1/a}$, fixes every ordinary rotation, rescales every
infinitesimal angle, and is continuous in the fine topology by \cref{trigonometry:rot:thm:finesplit}.
For $a\ne1$ it is not algebraic: an algebraic endomorphism fixing all ordinary units has exponent $1$
and is the identity. No basis of a large vector space is chosen. Likewise
\begin{equation}\label{trigonometry:rot:eq:unipotent}
 u\longmapsto\begin{pmatrix}1&\ellm(u)\\0&1\end{pmatrix}\in\GL_2(\No)
\end{equation}
is a homomorphism, nontrivial on infinitesimal rotations and trivial on ordinary ones. The line
spanned by $(1,0)$ is invariant and has no invariant complement, since a vector $(x,y)$ with $y\ne0$
is moved by $\ellm(u)y$ in its first coordinate; so this fine-continuous representation is not
semisimple, consistent with the semisimplicity of the algebraic representations above.

\subsection{The fine topology of the circle}\label{trigonometry:rot:sec:topology}
The fine topology is that of \cref{trigonometry:prop:topology}: balls $|z-a|<\rho$ with $\rho$ any
positive surreal, the circle carrying the induced topology. Coordinatewise order neighbourhoods give
the same topology, because $\max\{|x|,|y|\}\le|x+iy|\le|x|+|y|$; addition, multiplication, inversion
away from zero and conjugation are continuous by their elementary difference identities, so the
circle is a class topological group in this explicit neighbourhood sense. \emph{All positive surreal
radii are allowed}: restricting to real radii, or to radii in one set-sized Hahn subfield, gives a
different topology; in particular a set-sized Hahn field $K_\Gamma$ has its own intrinsic order
topology, while the topology it inherits from $\SC$ is discrete by
\cref{trigonometry:prop:topology}. Enlarging the admissible radii changes the topology but no
algebraic identity.

The existence of a positive surreal below any \emph{set} of positive surreals, used in the proof of
\cref{trigonometry:prop:topology}, also gives: \emph{every set-indexed Cauchy net is eventually
constant.} Indeed, take a positive lower bound for all its nonzero pairwise distances; on a tail all
pairwise distances are smaller than it, hence zero. \emph{This is special to the full surreal class}:
it is not a theorem about a non-Archimedean real closed set field with its intrinsic topology, it does
not forbid a properly class-indexed net of radii from approaching zero, and the restriction on index
size is essential.

The subgroup $\mm_{\R}\subset\No$ is open, since it contains $(-\omega^{-1},\omega^{-1})$ and its
translates by elements of $\mm_{\R}$; an open subgroup is closed, its complement being a union of open
cosets. The same holds for $\mm_{\C}$ and every nonzero multiple $\rho\,\mm_{\C}$. For $a\in\SC$ and
$\rho>0$ the set $a+\rho\,\mm_{\C}$ is therefore clopen and lies in the ball of radius $\rho$; and if
$a\ne b$, taking $\rho=|a-b|$ separates them, since $(b-a)/\rho$ has norm one and is not
infinitesimal.

\begin{proposition}[Total separation without isolated points]\label{trigonometry:rot:prop:totalseparated}
The fine circle $\TT(\No)$ is totally separated and has a basis of clopen neighbourhoods, but has no
isolated points. In particular topological connectedness is not the meaning of the algebraic
connectedness of the norm-one torus.
\end{proposition}
\begin{proof}
The clopen sets above separate any two points. Given $u$ and $\rho>0$, choose a positive
infinitesimal $\eta<\rho/4$; then $0<|\Ex(i\eta)-1|=2\sin(\eta/2)<\eta<\rho$, so $u\Ex(i\eta)$ is a
point of the ball other than $u$.
\end{proof}
So the circle is not discrete as a class, although every set-sized subset is; in the
universe-relative reading, small subsets are discrete while the whole larger-universe carrier is not.
No contradiction is involved.

\begin{theorem}[The splitting is topological only with a discrete real factor]
\label{trigonometry:rot:thm:finesplit}
With the fine topology on $\TT(\No)$ and on $\mm_{\R}$, the splitting \cref{trigonometry:eq:circlesplit}
is a homeomorphism
\begin{equation}\label{trigonometry:rot:eq:finesplit}
 \TT(\No)_{\mathrm{fine}}\ \cong\ \TT(\R)_{\mathrm{discrete}}\times(\mm_{\R})_{\mathrm{fine}},
\end{equation}
and the fine quotient topology on $\TT(\No)/\Tm$ is discrete, not the usual topology of the real
circle.
\end{theorem}
\begin{proof}
The monad $\Tm=\TT(\No)\cap(1+\mm_{\C})$ is open, and its cosets are the fibres of $\st$. Within one
fibre, for infinitesimal $\eps,\eta$,
$\tfrac12|\eps-\eta|\le|\Ex(i\eps)-\Ex(i\eta)|\le2|\eps-\eta|$, since the ratio is $1$ plus an
infinitesimal when $\eps\ne\eta$ (\cref{trigonometry:cor:phaseisometry}). So $\Ex(i\,\cdot)$ and
$\ellm$ are mutually inverse continuous maps on each monad. All monads being open, the circle is their
topological disjoint union, which is the product with a discrete real factor; every union of fibres
is open, which gives the quotient statement.
\end{proof}
The embedded ordinary circle is itself a closed discrete subspace; its usual topology is recovered
only by a coarser choice (\cref{trigonometry:rot:sec:othergeometry}), not as a subspace topology.

\begin{corollary}[Small subgroups]\label{trigonometry:rot:cor:subgroups}
Every subgroup of the fine circle generated by a set is closed and discrete, and none is dense.
\end{corollary}
\begin{proof}
Such a subgroup consists of finite words in the generators and their inverses, hence is a set, and
\cref{trigonometry:prop:topology} applies. It is not the whole circle, which is a proper class, the
rational chart embedding all of $\No$.
\end{proof}
In particular an ordinary irrational rotation has a countable orbit that is closed and discrete in
the fine topology, although its orbit is dense in the usual real circle: pigeonholing $N+1$ phases
into $N$ equal arcs gives a nonzero angular step below $2\pi/N$. Arbitrarily small \emph{real} steps
suffice for ordinary density but do not approach every surreal scale. There is consequently no
ordinary real Lie-group interpretation of the fine circle: a set-sized compact subspace would be
compact and discrete, hence finite, while no neighbourhood is finite. \emph{Classical compactness
arguments and Haar integration cannot be imported from the familiar picture of $\SO(2)$.}

Strong sums are not fine limits. With $t=\omega^{-1}$, $\Ex(it)$ is well defined, but the error of its
degree-$N$ Taylor polynomial has leading exponent $N+1$, so for every ordinary $N$ it exceeds the
fixed positive surreal $t^\omega=\omega^{-\omega}$ in absolute value; the partial sums do not converge
finely. Similarly $\Ex(it^n)$ does not converge finely to $1$, all its nonzero distances having finite
valuation $n<\omega$. In the intrinsic topology of the field $\R((t^{\Q}))$, by contrast, the integers
are cofinal in the value group, so $t^n\to0$ and $\Ex(it^n)\to1$ there; in a higher-rank fixed field
even this may fail if some positive value exceeds every $n\,v(t)$. These examples separate three
notions: exact Hahn summation, convergence in a fixed workspace, and convergence at all surreal
scales.

\subsection{The standard-part topology and semialgebraic paths}
\label{trigonometry:rot:sec:othergeometry}
The \emph{standard-part topology} on $\TT(\No)$ has open sets $\st^{-1}(V)$, $V$ open in the usual
topology of $\TT(\R)$. It deliberately forgets infinitesimal directions. It is a group topology, since
$\st$ is a homomorphism, but it is not $T_0$: distinct units with the same standard part have the same
neighbourhoods, and identifying topologically indistinguishable points yields the usual real circle.
The map $\st$ from the fine circle to the usual circle is continuous, preimages of real open sets
being unions of open monads, but it is not a quotient map: the quotient topology induced from the
fine source is discrete (\cref{trigonometry:rot:thm:finesplit}). \emph{These two uses of the same
set-theoretic map must not be conflated.}

Let $F$ now be a fixed real closed \emph{set} field. In the semialgebraic category a path is a
continuous semialgebraic map on $[0,1]_F$, not an ordinary real-parameter path. A unit $u\ne-1$ is
joined to $1$ by
\begin{equation}\label{trigonometry:rot:eq:semipath}
 \gamma_u(s)=\frac{(1-s)+su}{|(1-s)+su|},\qquad0\le s\le1,
\end{equation}
whose denominator never vanishes, since a zero would make $u$ a negative real unit, that is $-1$;
$1$ is joined to $-1$ by two such arcs through $i$, and multiplying by a starting unit joins any two
points. The nonnegative square root is semialgebraic, so this proves semialgebraic path
connectedness directly; the circle is closed and bounded, hence definably compact. The semialgebraic
fundamental group is the usual one,
\begin{equation}\label{trigonometry:rot:eq:sapi}
 \pi_1^{\mathrm{sa}}\bigl(\TT(F),1\bigr)\cong\Z;
\end{equation}
\emph{this is imported, not proved here}: it invokes the standard comparison and scalar-extension
theorems of semialgebraic homotopy \cite{DK,BaroOtero}, applicable because the circle is defined
without parameters. There is no conflict with \cref{trigonometry:prop:topology}: the topology is
intrinsic to $F$ and the parameter interval is $[0,1]_F$. The same formula on $[0,1]_{\No}$ with the
fine topology is a nonconstant continuous class map, but that domain is not connected in the ordinary
sense; and $[0,1]_{\R}$ with the topology induced from $\No$ is discrete, not the usual interval.

Four statements therefore hold simultaneously and answer different questions: the algebraic torus is
geometrically connected of dimension one; its semialgebraic $F$-points are definably path connected;
its full fine surreal points are totally separated; its standard-part quotient is the ordinary
connected circle. Correspondingly the exact sequence $2\pi\Z\to\OO_{\R}\to\TT(\No)$ is an algebraic
group statement with local analytic charts, \emph{not} a claim that a connected real line is a
universal covering space of the fine circle; semialgebraic and ordinary covering theories have their
own categories and hypotheses.

\subsection{Local angular velocities}\label{trigonometry:rot:sec:flows}
The fine derivative (\cref{trigonometry:sec:foundations}), which is local Taylor calculus and not
differentiation of a convergent sequence of partial sums \cite[\S3]{BMgerms}, gives $\frac{d}{d\theta}\cis\theta=i\cis\theta$
and $\frac{d}{d\theta}\Rot(\cos\theta,\sin\theta)=\mathsf J\,\Rot(\cos\theta,\sin\theta)$ on finite
angles. For a fixed $a\in\No$ the canonical local flow is
\begin{equation}\label{trigonometry:rot:eq:localflow}
 \Rot_a(x)=\Rot\bigl(\cos(ax),\sin(ax)\bigr),\qquad x\in D_a=\{x\in\No:ax\in\OO_{\R}\},
\end{equation}
with $D_a=\No$ if $a=0$ and otherwise $D_a=a^{-1}\OO_{\R}$, an open additive subgroup; on $D_a$,
$\Rot_a(x+y)=\Rot_a(x)\Rot_a(y)$, $\Rot_a'(x)=a\mathsf J\Rot_a(x)$ and $\Rot_a(0)=I$. Every surreal
angular velocity is allowed locally, but an infinite velocity requires a correspondingly small
canonical time domain.

Every extension $\Psi_\chi$ of \cref{trigonometry:thm:characters} obeys the same local law
\cref{trigonometry:eq:globallocal}, is fine-continuous, and has $\Psi_\chi'=i\Psi_\chi$ and
$\Psi_\chi^{(n)}=i^n\Psi_\chi$. The extensions of \cref{trigonometry:ex:phases} therefore have
identical finite behaviour and identical local differential laws but can disagree at $\omega$:
\emph{solving the local rotation equation and imposing $\Psi(0)=1$ does not choose a global phase
normalization.} A sharper warning concerns solutions not assumed to be homomorphisms.

\begin{example}[A locally constant factor]\label{trigonometry:rot:ex:locallyconstant}
Let $g(x)=1$ for $x\in\mm_{\R}$ and $g(x)=-1$ otherwise. Both regions are fine-clopen, so $g$ is
locally constant with zero fine derivative everywhere, and is not constant. On $\OO_{\R}$ the function
$f(x)=g(x)\cis x$ satisfies $f'=if$ and $f(0)=1$, yet differs from $\cis x$ off the identity monad.
\end{example}
So even on a finite surreal interval the ordinary mean-value and initial-value uniqueness arguments
cannot be used without further hypotheses. A coherent Taylor specification on every real-centred
monad together with the ordinary real values does select $\cis$ (\cref{trigonometry:prop:normalization});
that is additional global data, not a consequence of fine differentiability, and the classification
of \cref{trigonometry:thm:characters} is not a classification of solutions of the fine differential
equation. A derivation of the scalar field, such as a Berarducci--Mantova derivation \cite{BM}, is yet
another operation; no compatibility between it and a globally normalized value such as
$\Sin\omega=0$ is assumed.

\subsection{Exact computation with rotations}\label{trigonometry:rot:sec:computation}
Substantial rotation algebra can be implemented before arbitrary surreals are.
\begin{center}\small
\begin{tabular}{L{0.21\textwidth}L{0.33\textwidth}L{0.36\textwidth}}
\toprule
Representation & Best operations & Necessary caution\\
\midrule
$(a,b)$, $a^2+b^2=1$ & Action on coordinates, dot products, exact matrix operations & Normalizing an arbitrary pair needs a positive square root\\
Projective $[p:q]$ & Composition \cref{trigonometry:eq:projectiveaddition} and inverse $[-p:q]$ by field arithmetic alone & Keep the point at infinity; equality is $ps-qr=0$, not coordinatewise\\
$(u_0,\eps)$, $u_0\in\TT(\R)$, $\eps\in\mm_{\R}$ & Products, matched roots, valuations, infinitesimal error accounting & Exact real phase data and Hahn support information must be represented honestly\\
\bottomrule
\end{tabular}
\end{center}
In split form $(u_0,\eps)(w_0,\eta)=(u_0w_0,\eps+\eta)$ and $(u_0,\eps)^{-1}=(u_0^{-1},-\eps)$, exactly;
a chosen $N$th-root branch divides $\eps$ by $N$ and takes the specified ordinary root of $u_0$, and
\cref{trigonometry:rot:eq:producterror} gives the valuation of propagated error in a finite-support
approximation. \emph{No representation of arbitrary surreal class elements by finite strings is
asserted, and equality of effective Hahn-series names is not automatically decidable}; these concern
representations and algorithms, not the definition of the group.

Source 19 proposes a layered formal development: finite algebra (matrix/unit-circle isomorphism,
isometry rigidity, projective chart, action, invariants) over an ordered field with square roots; a
valuation layer ($\OO_{\R}$, $\mm_{\R}$, standard part, conjugation); a strong-series layer (exp/log
inversion, purely imaginary logarithm of a near-one unit), which yields the splitting; ordinary
complex phase for the finite-angle uniformization and kernel; a separate topology layer resting on
the lower-bound cut; a semialgebraic layer; and a global layer resting on $\No=\Pi\oplus\OO_{\R}$
rather than on a Taylor series at an infinite argument. It names the exact class-group splitting,
with its fine-topological version, as a next formal target, and says this order is compatible with the
repository's division between generic finite algebra, Hahn evaluation and concrete surreal
foundations. The collection's formalization ledger \texttt{docs/FORMALIZATION.md} is not among the
files source 19 lists as inspected. Already at its pin, and still now, that ledger records Lean proofs, on
the actual surcomplex field, of several of these layers as stated in the analysis report --- the
infinitesimal exponential/logarithm group isomorphism, the purely imaginary logarithm of a unit, and
surjectivity of the finite phase onto the unit circle with kernel exactly $2\pi\Z$ --- and of the
affine and projective circle charts over any ordered field (mapped, as prerequisites, to
\cref{trigonometry:thm:cayley}). The splitting \cref{trigonometry:eq:circlesplit} itself, its
fine-topological version \cref{trigonometry:rot:thm:finesplit}, and every label of this section have
no Lean mapping; nothing in this section is formally verified.

\subsection{Comparison with the classical circle}\label{trigonometry:rot:sec:comparison}
\begin{center}\small
\begin{longtable}{L{0.25\textwidth}L{0.27\textwidth}L{0.38\textwidth}}
\toprule
Question & Ordinary real plane & Surreal plane\\
\midrule\endhead
Origin-fixing rotations & $\SO(2,\R)$ & $\SO(2,\No)$, a class group\\
Matrix form & $\Rot(a,b)$, $a^2+b^2=1$ & The same equations\\
Geometric phase domain & $\R/2\pi\Z$ & $\OO_{\R}/2\pi\Z$\\
Infinitesimal rotations & Trivial & $\Tm\cong(\mm_{\R},+)$\\
Canonical factorization & Ordinary phase & Ordinary phase times exact infinitesimal phase\\
Torsion; finite subgroups & Roots of unity; cyclic & The same literal roots of unity and $\mu_n$\\
Roots of degree $n\ge1$ & Exactly $n$ & Exactly $n$, one infinitesimal lift per standard root\\
Algebraic Lie algebra & $\R\mathsf J$ & $\No\mathsf J$, in the universe-relative reading\\
Algebraic endomorphisms; weights & Integer power maps; integers & The same; still ordinary integers\\
Fine topology & The usual topology & Totally separated, no isolated points, set-sized subsets discrete\\
Ordinary real paths in the fine target & Nonconstant ones exist & All constant\\
Standard-part quotient & The circle itself & The usual circle, only for the coarser standard-part topology\\
Semialgebraic $\pi_1$ & $\Z$ & $\Z$ over a fixed real closed set field (imported)\\
Global canonical phase & Periodic exponential & $\Phi(x)=\cis(\fin x)$, kernel $2\pi\Oz$\\
Infinitely many repeated turns & Not a finite group operation & Still requires an explicit summation or phase convention\\
\bottomrule
\end{longtable}
\end{center}
The resulting picture is neither ``just the usual circle'' nor ``trigonometry breaks down'': the
group has the familiar finite algebra, a very large exact infinitesimal angular factor, and several
useful but inequivalent geometries, and which one is relevant depends on whether a question concerns
finite equations, asymptotic scales, ordinary observations or definable paths. Higher-dimensional
orthogonal groups are a natural extension, where infinitesimal logarithms are noncommutative and
Baker--Campbell--Hausdorff terms appear; source 19 leaves them aside, and the collection now contains
a separate report on the three-dimensional case, \texttt{docs/surreal/euclidean-three-space}.

\section{Period arithmetic and noncanonical surcomplex exponentials}\label{trigonometry:per:sec:main}

This section is the contribution of sources~20 and~21. It continues
\cref{trigonometry:sec:global}: \cref{trigonometry:thm:characters} says that every extension of the
finite phase to all real surreal arguments is $\Psi_\chi$ for a character $\chi$ of $\Pi$, and
\cref{trigonometry:thm:infiniteperiods} says that every such extension has infinite periods. Here we
study the \emph{arithmetic} of those periods: which of them remain periods after division by every
ordinary integer, which scalars preserve all of them under multiplication, and whether a choice of
$\chi$ can be made invariant under automorphisms of the underlying valued field. The main
application is a partial negative answer to the first open question of \cite[\S11.1]{EK}: a
surcomplex exponential with the canonical strip values, the modulus law and all local Taylor
identities can have kernel-multiplier ring exactly $\Z$, which is initial in the surreal
simplicity tree but is not an integer part (\cref{trigonometry:per:thm:robustness}). Only that
specific sufficiency possibility is refuted; the question itself stays open
(\cref{trigonometry:per:sec:ek}).

\subsection{What sources 20 and 21 contribute, and the conventions of this section}
\label{trigonometry:per:sec:overlap}
\emph{The two manuscripts.} Source~20 is \emph{Exponential Kernels and Arithmetic Rigidity in the
Surcomplex Field: Minimal Multiplier Rings, Nonsplit Periods, and the Limits of Canonical
Extension}; source~21 is \emph{Period Arithmetic and Noncanonical Surcomplex Exponentials: Profinite
Defects, Minimal Multiplier Rings, and a Partial Answer to an Ehrlich--Kaplan Question}. Both are
dated September 22, 2026, and both were written against repository commit
\texttt{465a54b479a1}, at which this article stood exactly as it did before they were merged here
(Sections~1--18 and the concluding section were unchanged between that commit and this merge).
Both take \cref{trigonometry:thm:globalexp,trigonometry:thm:characters,trigonometry:thm:infiniteperiods}
and the local law \eqref{trigonometry:eq:globallocal} as prerequisites, restate
\cref{trigonometry:thm:characters} in turn coordinates, and do not claim these results; they are
cited here and not proved again. Source~21 also records that no profinite-defect or
minimal-multiplier construction was found in this article at that commit, which is correct.

\emph{One spine, two bases of hypothesis.} The two manuscripts were written independently and prove
the same central fact: a normalized surcomplex exponential can have kernel-multiplier ring $\Z$. Source~21
is the base text of this section, because it proves more under weaker hypotheses: its set-sized theorem
(\cref{trigonometry:per:thm:minimalset}) is proved in ZFC on a fixed continuum-sized field such as
$\No(\eps_0)$ or $F_{\Q}=\R((t^{\Q}))$; it prescribes an arbitrary ordinary character on all of
$\Pi$ and an arbitrary profinite defect; and it keeps the standard part of the phase fixed at every
argument. Source~20's own set-sized examples live on arbitrarily large, unspecified fragments
$\No(\lambda)$ obtained by reflecting a construction on all of $\No$ made with global choice
(\cref{trigonometry:per:thm:reflection}). Source~20 contributes what source~21 lacks: a relative
construction with values in the ordinary circle, the reflection theorem, the external invariant
$\Hom(\cI,\Z)$, the first-order sentence $\mathsf{RingKer}$, definability of $\Z$ with an omitted
type, and the example of a non-initial multiplier ring. Source~21 contributes what source~20 lacks:
the profinite defect and its realization, the cardinal count on $F_{\Q}$, cofinal divisible periods
on $\No$, and the free action of valued-field translations with the shadow formula.

\emph{Printed once.} The following results are proved in both manuscripts; each is printed once, at
the place indicated, with both credited. The proofs agree in substance except where a second route is
marked.
\begin{center}\small
\begin{longtable}{L{0.30\textwidth}L{0.17\textwidth}L{0.17\textwidth}L{0.24\textwidth}}
\toprule
Result & Source 20 & Source 21 & Printed as\\
\midrule\endhead
Finite phase in turns, kernel $\Z$ & Prop.~3.1 & Prop.~2.1 & \cref{trigonometry:thm:polar} (radians); not reprinted\\
Character coordinates of a phase & Prop.~3.3 & Prop.~3.2 & \cref{trigonometry:thm:characters}; coordinates in \cref{trigonometry:per:prop:coordinates}\\
Phases $\leftrightarrow$ exponentials; modulus, conjugation, kernel & Prop.~3.4 & Prop.~8.1 & \cref{trigonometry:per:prop:exponentials}\\
Kernels are exactly the additive integer parts & Thm.~4.1 & Thm.~3.4 & \cref{trigonometry:per:thm:dictionary}\\
Local law and all fine derivatives & Prop.~10.1 & Props.~3.6, 8.3 & \cref{trigonometry:eq:globallocal}; \cref{trigonometry:per:prop:local}\\
Split period sequence $\Leftrightarrow$ additive angle lift $\Leftrightarrow$ zero defect & Thm.~7.1 & Thm.~4.5 & \cref{trigonometry:per:thm:split}\\
Explicit nonliftable profinite character & Lemma 7.2 (3-adic) & \S10 (2-adic) & \cref{trigonometry:per:ex:profinite}, both kept\\
Multiplier ring and the integer-part criterion & Lemma 5.1, Thm.~5.2 & Props.~7.1, 8.2, eq.~(8.7) & \cref{trigonometry:per:prop:ZE,trigonometry:per:thm:criterion}\\
Minimal multiplier ring $\Z$ by two-vector avoidance & Lemma 6.1, Thms.~6.2, 6.4 & Lemmas 7.2--7.6, Thms.~7.5, 7.7 & \cref{trigonometry:per:thm:minimalset,trigonometry:per:thm:minimalclass}; source 20's different route is \cref{trigonometry:per:thm:relative}\\
Initial multiplier ring that is not an integer part & Cor.~6.5 & Thm.~8.4 & \cref{trigonometry:per:thm:robustness}\\
Strong shift automorphisms & Thm.~11.1 & Thm.~9.2 & \cref{trigonometry:per:thm:shifts}, which is \cite[Theorem 5.2]{Saut}\\
No phase equivariant under valued-field automorphisms & Thm.~11.2 ($c=1/2$) & Thm.~9.3, Cor.~9.4 & \cref{trigonometry:per:thm:free}; source 20's case in \cref{trigonometry:per:rem:halfshift}\\
\bottomrule
\end{longtable}
\end{center}

\emph{Credits the sources owe.} The shift automorphisms of \cref{trigonometry:per:thm:shifts} are the
fixed-shift flows of the collection's report on field automorphisms of $\SC$ \cite[Theorem 5.2,
label \texttt{saut:thm:shiftflow}]{Saut}, with its parameters $\delta=1$ and $\lambda=[\,\cdot\,]_0$
and the sign of the flow parameter reversed; source~21 cites that theorem and \cite[\S4]{KKS},
source~20 cites neither the repository report nor the flows. The residue map, the maximal divisible
subgroup and the splitting and homomorphism facts behind
\cref{trigonometry:per:prop:divisible,trigonometry:per:thm:split,trigonometry:per:thm:homz} are
classical facts about $\Z$-groups \cite[\S\S2--4]{Jerabek}; source~21 says so, source~20 proves its
versions without this citation. Neither source claims them as new, and neither is credited with them
here.

\emph{Where the merge had to choose.} (i) Source~21's text is the base, and source~20's statements are
inserted where it has no counterpart. (ii) Where both prove one result, one proof is printed.
(iii) Source~20's minimal-multiplier construction is not the same theorem as source~21's: it changes
the ordinary character (outside a prescribed set-sized subspace) and keeps the infinitesimal
character zero, so that all phase values lie in the ordinary circle, which is what its reflection
theorem needs; source~21 fixes the ordinary character everywhere and changes only the infinitesimal
character. Both are printed (\cref{trigonometry:per:thm:relative} and
\cref{trigonometry:per:thm:minimalset,trigonometry:per:thm:minimalclass}). (iv) Source~20's
non-equivariance theorem is the special case $c=1/2$ of source~21's free action; source~20's shorter
proof is kept as a remark. (v) Two consequences are obtained by combining the sources and are stated by
neither: source~20's first-order sentence separates source~21's example on $\No(\eps_0)$ from the
canonical exponential, and source~21's defects give three pairwise nonisomorphic exponentials on
$\No(\eps_0)[i]$ in ZFC (\cref{trigonometry:per:cor:combined}). Both are marked where they are stated.

\emph{Set theory, kept attached to each statement.} Statements about a set-sized field are theorems
of ZFC. Statements about a class map on all of $\No$ are made in NBG, reading ``all phases'' as a
scheme with a class parameter rather than as a class whose elements are proper classes; every
normal form, coefficient family and recursion stage is a set, and a quotient of a proper class is
never treated as a set of classes. \emph{Global choice} (NBG with global choice, written GBC) is
used exactly where it is named: in \cref{trigonometry:per:thm:minimalclass}, in
\cref{trigonometry:per:thm:relative} and its variants, in the $\No[i]$ case of
\cref{trigonometry:per:thm:robustness}, and to build the class maps that
\cref{trigonometry:per:thm:reflection} reflects. The realization theorem
(\cref{trigonometry:per:thm:realize}) uses only a choice between set-sized vector spaces, even for a
defect defined on all of $\Pi$. The set-sized consequences of the global-choice constructions rest on
the conservativity of GBC over ZFC for statements about sets; no large cardinal, saturation hypothesis
or continuum hypothesis is used anywhere in this section.

\emph{Radians, turns, and the symbols of this section.} This article measures angles in radians: the
finite phase is $\cis$ with kernel $2\pi\Z$, and $\Phi$, $\Psi_\chi$, $\Exp$ and $E_\chi$ are exactly the
objects of \cref{trigonometry:sec:global}. Both sources measure phases in \emph{turns}, with finite phase
$a\mapsto e^{2\pi ia}$ and kernel $\Z$; a turn-phase at $x$ is here the radian phase at $2\pi x$. Because
the arithmetic concerns the ordinary integers, periods are divided by $2\pi$ once and for all: the
\emph{normalized period group} of a phase is $\cI(\Psi)=\{a:\Psi(2\pi a)=1\}$
(\cref{trigonometry:per:def:periodgroup}). The remaining conversions are:
\begin{itemize}
\item Purely infinite part: the sources' $\mathcal P$ is $\Pi$ (the letter $\PP$ is used for the
 projective line). Principal part $\pr(x)=x-\fin x\in\Pi$ (source 21's $\operatorname{pr}$; source 20
 writes $p(x)$).
\item Phases: source~21's admissible phase $\Phi$ and canonical phase $\Phi_0$, and source~20's
 normalized phase $c$ and canonical phase $c_0$, are here $\Psi$ and $\Phi$; the letter $\Phi$ keeps
 its meaning from \cref{trigonometry:def:globaltrig} and never denotes an arbitrary phase (other
 reports of the collection use $\Phi$ for unrelated objects, none of which occurs here). Source~20's circle $\mathcal U_F$ is $\TT(F)$, and its
 $\TT=\R/\Z$ is always written $\R/\Z$ here, since $\TT(F)$ is the unit circle.
\item Characters: the character $\chi:\Pi\to\TT(\No)$ of \cref{trigonometry:thm:characters} has
 coordinates $(\psi,\cL)$ (\cref{trigonometry:per:prop:coordinates}). Source~21's pair $(h,L)$ is
 $(\psi,\cL)$; source~20's ordinary-circle character $\chi:\mathcal P\to\R/\Z$ is $\psi$ with
 $\cL=0$, and its general character $\eta$ is our $\chi$. The letter $h$ keeps its meaning as a generic
 infinitesimal increment, and $L$ its meaning in \cref{trigonometry:thm:infinitefrequency}.
\item Exponentials: $E_\chi(x+iy)=\exp_{\No}(x)\Psi_\chi(y)$ as in \cref{trigonometry:thm:characters};
 source~20's $E_\chi$ with $y/(2\pi)$ and a turn-valued $\chi$ is our $E_\chi$ for the character with
 coordinates $(\chi,0)$. Source~20's finite phase $\mathrm e_0$ and circle map $\mathrm e(\theta)$ are
 $f\mapsto\cis(2\pi f)$ and $\theta\mapsto e^{2\pi i\theta}$; source~21's $\mathrm e_{\OO}$ and
 $\mathrm e_{\R}$ are the same two maps. (In their sources one macro name stands for $\exp_{\No}$ in
 source~20 and for $r\mapsto e^{2\pi ir}$ in source~21; the printed symbols differ, and neither is
 used here.) The strip exponential of
 \cref{trigonometry:thm:stripexp} keeps the letter $E$ without subscript.
\item Periods and multipliers: source~20's normalized kernel $D$ and source~21's $I_\Phi$ are $\cI$;
 source~21's maximal divisible subgroup $D_\Phi$ or $D(I)$ is $\cI^{\mathrm{div}}$ (so neither source's
 $D$ survives); source~20's $M(E)$ and source~21's $Z_E$ are $Z_E$; $\Mult$ is unchanged.
\item Profinite objects: source~21's defect $\delta_\Phi$ is $\pd_\Psi$ (the letter $\delta$ is the side gap
 $b+c-a$, and ``defect'' alone means the angle defect $d(\alpha)$ or the triangle-inequality defect,
 so this one is always called the \emph{profinite defect}); its residue map $\rho_I$ is $\res_{\cI}$;
 its $\mathcal V=\Zhat/\Z$ is written $\Zhat/\Z$; its $\Sigma$, $\partial$, $s$ are $\Qv$, $\cl$,
 $\varsigma$; its continuum $\mathfrak c$ is $2^{\aleph_0}$; its representative $\theta(p)$ is
 $\bar\psi(p)$; its retraction $r$ is $\mathfrak r$.
\item Fields: source~21's $K_\Gamma=\R((t^\Gamma))$ is $F_\Gamma$ of \cref{trigonometry:sec:rings}, since
 $K_\Gamma$ denotes $\C((t^\Gamma))$ in this article and in the collection's notation guide; its $K_0$ and
 $K_\eps$ are $F_{\Q}$ and $\No(\eps_0)$. Source~20's $F$ keeps its name here; birthdays are written
 out in words.
\item Automorphisms: source~21's $T_c$ and source~20's $\sigma_a$ are $\Sh_c$ ($T_n$ is a Chebyshev
 polynomial, $\sigma$ a valuation threshold, and $T_s$ in \cite{Saut} has the opposite sign); their
 exponent map $\ell(g)$ is $[g]_0$, the real coefficient of $\omega^0$ in the exponent $g$, since
 $\ell(p)$ is the coefficient of $\omega$ in \cref{trigonometry:ex:phases}. Source~21's $[X]p$ is that
 $\ell(p)$. Source~20's witness monomial $h_x$ is $\mu_x$.
\end{itemize}

\subsection{Phases on a field, periods, and exponentials}\label{trigonometry:per:sec:phases}
\emph{Fields.} In this section $F$ is one of the following: $\No$; a set-sized initial
trigonometric-exponential subfield $\No(\lambda)=\{x\in\No:\text{the birthday of $x$ is less than
$\lambda$}\}$ for an epsilon number $\lambda$, in particular $\No(\eps_0)$
\cite[\S11.2, Theorem 11.1]{EK}, \cite{vdDE}; or, where no real exponential is needed, a Hahn field
$F_\Gamma=\R((t^\Gamma))$, in particular $F_{\Q}$, or more generally a non-Archimedean
truncation-closed subfield of $\No$ containing $\R$ on which the finite phase is available. Initial
subfields are truncation closed. $\No(\eps_0)$ has cardinality $2^{\aleph_0}$, since its elements have
countable sign sequences, of countably many lengths, and it contains $\R$; it is \emph{not} asserted to
be a full Hahn field. Write $\OO_F,\mm_F,\Pi_F$ for the finite, infinitesimal and purely infinite
elements of $F$; truncation closure gives $F=\Pi_F\oplus\OO_F$ and $\OO_F=\R\oplus\mm_F$. The group
$\Pi_F$ is a $\Q$-vector space and a nonunital subring, and its nonzero elements are infinite. The
finite phase $\cis:(\OO_F,+)\to\TT(F)$ is onto with kernel $2\pi\Z$, and
$\TT(F)\cong\TT(\R)\times(\mm_F,+)$ through $\ellm$ of \cref{trigonometry:rot:eq:ellm}: this is
\cref{trigonometry:thm:polar} for $\No$, and the same proof for these $F$, as both sources note. The
occurrence of the ordinary $\Z$, not $\Oz$, in this kernel is essential.

\begin{definition}[Phases and normalized periods]\label{trigonometry:per:def:periodgroup}
A \emph{phase} on $F$ is a homomorphism $\Psi:(F,+)\to\TT(F)$ whose restriction to $\OO_F$ is
$\cis$. Its \emph{normalized period group} is $\cI(\Psi)=\{a\in F:\Psi(2\pi a)=1\}$, so that
$\ker\Psi=2\pi\cI(\Psi)$. An \emph{additive integer part} of $F$ is an additive subgroup
$\cI\subseteq F$ containing $1$, with least positive element $1$, such that every $x\in F$ has an
$a\in\cI$ with $a\le x<a+1$; no multiplicative closure is included. An \emph{integer part} is an
additive integer part that is a subring.
\end{definition}
The canonical phase of \cref{trigonometry:def:globaltrig} has $\cI(\Phi)=\Oz$; on an initial subfield
$F$ it has $\cI=\Oz\cap F$, the unique initial integer part of $F$ \cite[Lemma 11.1]{EK}, and on
$F_\Gamma$ the phase $\cis\circ\fin$ has $\cI=\Pi_F\oplus\Z$. \emph{Initial} means closed under
proper sign-sequence prefixes; it does not mean small, ordered or definable.

\begin{proposition}[Coordinates of a phase; restating \cref{trigonometry:thm:characters}]
\label{trigonometry:per:prop:coordinates}
Every phase on $F$ is $\Psi_\chi(p+f)=\chi(p)\cis f$ for a unique homomorphism
$\chi:\Pi_F\to\TT(F)$, and there are unique maps $\psi:\Pi_F\to\R/\Z$, additive, and
$\cL:\Pi_F\to\mm_F$, $\Q$-linear, with
\begin{equation}\label{trigonometry:per:eq:coordinates}
 \chi(2\pi p)=e^{2\pi i\psi(p)}\,\Ex\bigl(2\pi i\cL(p)\bigr)\qquad(p\in\Pi_F).
\end{equation}
Every such pair occurs; we write $\Psi_{\psi,\cL}$ for the phase. The canonical phase is
$\Psi_{0,0}=\Phi$.
\end{proposition}
\begin{proof}
The first assertion is \cref{trigonometry:thm:characters}, whose proof applies verbatim on $F$
(sources 20 and 21 state it so). Since $\Pi_F=2\pi\Pi_F$, compose $p\mapsto\chi(2\pi p)$ with the
splitting $\TT(F)\cong\TT(\R)\times\mm_F$: the first coordinate is an additive map to
$\TT(\R)\cong\R/\Z$, the second is additive between uniquely divisible groups, hence $\Q$-linear.
\end{proof}
Source~21 uses these coordinates to separate the ordinary phase from the infinitesimal one; source~20
works mainly with $\cL=0$, where every value of $\chi$ lies in the ordinary circle.

\begin{theorem}[Kernel dictionary]\label{trigonometry:per:thm:dictionary}
The map $\Psi\mapsto\cI(\Psi)$ is a bijection from phases on $F$ onto additive integer parts of $F$.
Every such $\cI$ has unique floors, satisfies $\cI\cap\OO_F=\Z$, and sits in an exact sequence
\begin{equation}\label{trigonometry:per:eq:periodsequence}
 0\longrightarrow\Z\longrightarrow\cI\xrightarrow{\ \pr\ }\Pi_F\longrightarrow0 .
\end{equation}
In particular a phase is determined by its kernel. If $\Pi_F\ne0$, then $\cI$ contains infinite
elements and is not cyclic.
\end{theorem}
\begin{proof}
Let $\Psi$ be a phase. The finite restriction gives $\cI\cap\OO_F=\Z$. For $x\in F$ choose finite $f$
with $\cis(2\pi f)=\Psi(2\pi x)$, so $x-f\in\cI$. Every finite $f=r+\eps$ has a unique ordinary integer
$n$ with $n\le f<n+1$: the ordinary floor of $r$, lowered by one when $r\in\Z$ and $\eps<0$. Then
$a=x-f+n\in\cI$ is a floor of $x$; two floors differ by an element of $\cI$ of absolute value less
than one, hence by $0$, and $1$ is least positive since $(0,1)$ is finite and meets $\cI$ in nothing.
Applying the same correction to $p\in\Pi_F$ proves that $\pr:\cI\to\Pi_F$ is onto; its kernel is
$\cI\cap\OO_F=\Z$.

Conversely let $\cI$ be an additive integer part. A finite element of $\cI$ minus its ordinary floor
lies in $\cI\cap[0,1)=\{0\}$, so $\cI\cap\OO_F=\Z$, and floors are unique. Put
$\Psi_{\cI}(y)=\cis\bigl(y-2\pi\lfloor y/2\pi\rfloor_{\cI}\bigr)$. For two arguments the floors of
the sum and of the summands differ by $0$ or $1$, which $\cis(2\pi\,\cdot\,)$ ignores, so
$\Psi_{\cI}$ is a homomorphism; it restricts to $\cis$ on $\OO_F$ and has normalized period group
$\cI$. Two phases with the same $\cI$ agree at $2\pi x$, because both take there the finite phase of
$x$ minus its common floor.

For $p\ne0$ in $\Pi_F$ every element of the fibre of $\pr$ over $p$ is infinite. If $\cI=\tau\Z$ were
cyclic, then $1=n\tau$ for an ordinary $n\ne0$, and every element of $\cI$ would be finite.
\end{proof}
The dictionary is source~20's Theorem 4.1 with Corollary 4.2 and source~21's Theorem 3.4; the proof
printed is source~21's. That $\cI$ meets every coset $p+\OO_F$ is
\cref{trigonometry:thm:infiniteperiods} in normalized form; the floor formulation isolates the
additional arithmetic used below. For $F=\No$ and $\Psi=\Phi$ the floor is Ehrlich and Kaplan's
integer-part floor \cite[\S11]{EK}.

\begin{proposition}[Division with remainder; source 20]\label{trigonometry:per:prop:division}
For every ordinary $n\ge1$ and $d\in\cI$ there are unique $e\in\cI$ and $r\in\{0,\dots,n-1\}$ with
$d=ne+r$. Hence $\cI/n\cI\cong\Z/n\Z$, induced by the inclusion of $\Z$.
\end{proposition}
\begin{proof}
Let $e$ be the $\cI$-floor of $d/n$; then $r=d-ne\in\cI\cap[0,n)$ is an ordinary integer. If
$ne+r=ne'+r'$, then $e-e'=(r'-r)/n\in\cI\cap\OO_F=\Z$ has absolute value less than one.
\end{proof}
This arithmetic is additive; it does not allow arbitrary elements of $\cI$ to be multiplied inside
$\cI$, and that is exactly where the counterexamples below enter. For set-sized $\cI$ the ordered
group $\cI$ is discretely ordered with divisible quotient $\cI/\Z\cong\Pi_F$, that is, a $\Z$-group, a
model of Presburger arithmetic in its ordered-group presentation \cite[\S2]{Jerabek}; this does not make
it a model of arithmetic with multiplication.

\emph{Related (batch 33).} For $\cI=\Oz$ the collection's report on exponential relations over omnific
integers (\nolinkurl{docs/foundations-and-computation/exponential-relations-over-omnific-integers})
re-derives this proposition independently (label \texttt{exr:prop:division}), with $x\equiv y\pmod n$
exactly when the constant terms of $x$ and $y$ are congruent, and reads the class-sized $\Oz$ formula by
formula as a $\Z$-group. It adds the zero predicates of finite exponential sums with algebraic
coefficients and real algebraic slopes. With rational slopes the result is a definitional expansion of
Presburger arithmetic, and $\Z$ is an elementary substructure (\texttt{exr:thm:rational}); with algebraic
slopes the theory is still decidable (\texttt{exr:thm:algdecision}). The boundary is the one stated
above: adjoining the graph of multiplication, as the relation $\exp(xy)=\exp(z)$, makes the theory
undecidable (\texttt{exr:thm:undecidable}).

\begin{proposition}[Phases and exponentials]\label{trigonometry:per:prop:exponentials}
Let $F$ carry a real exponential $\exp_F$, i.e.\ $F=\No$ or $F=\No(\lambda)$ with $\lambda$ an epsilon
number. For every phase $\Psi_\chi$ on $F$, $E_\chi(x+iy)=\exp_F(x)\Psi_\chi(y)$ is a surjective
homomorphism $(F[i],+)\to(F[i]^\times,\cdot)$ extending $\exp_F$ and the strip exponential of
\cref{trigonometry:thm:stripexp}, with
\begin{equation}\label{trigonometry:per:eq:explaws}
 |E_\chi(z)|=\exp_F(\re z),\qquad E_\chi(\bar z)=\overline{E_\chi(z)},\qquad
 \ker E_\chi=2\pi i\,\cI(\Psi_\chi).
\end{equation}
Conversely every additive-to-multiplicative homomorphism $E:F[i]\to F[i]^\times$ that agrees with the
strip exponential at finite imaginary arguments and satisfies $|E(z)|=\exp_F(\re z)$ is $E_\chi$ for
exactly one $\chi$, namely the restriction of $y\mapsto E(iy)$ to $\Pi_F$.
\end{proposition}
\begin{proof}
The forward assertions other than conjugation and the kernel are \cref{trigonometry:thm:characters}.
Conjugation follows from $\Psi(-y)=\Psi(y)^{-1}=\overline{\Psi(y)}$. A kernel element has modulus
$\exp_F(x)=1$, so $x=0$, and then $\Psi_\chi(y)=1$ means $y\in2\pi\cI$. Conversely $\Psi(y)=E(iy)$
has modulus one, is a homomorphism, and is $\cis$ on finite $y$; on the real axis $E=\exp_F$; the
group law gives $E(x+iy)=\exp_F(x)\Psi(y)$.
\end{proof}
Ehrlich and Kaplan establish the integer-part case \cite[Proposition 11.2]{EK}; here the formula is
used for additive integer parts that are not rings. The proposition is source~20's Proposition 3.4
and source~21's Proposition 8.1.

\emph{Related (batch 33).} For $F=\No$, on the additive domain $\overline{\mathbb Q}+\Pi$, whose imaginary
parts are finite algebraic numbers, every $E_\chi$ restricts to $a+p\mapsto e^a\exp_{\No}(p)$, the
exponential of the report on exponential relations over omnific integers (\texttt{exr:def:domain}).
There the values at finitely many distinct points are linearly independent over $\overline{\mathbb Q}$
(\texttt{exr:thm:independence}, from Gonshor's monomial theorem and the Lindemann--Weierstrass
theorem). So all the $E_\chi$ agree on that domain and are injective there; the choice of phase
classified here does not enter.

\begin{proposition}[The local law, for every phase]\label{trigonometry:per:prop:local}
For every phase $\Psi$ on $F$, every $x\in F$, $h\in\mm_F$, and, when $F$ has a real exponential,
every $z\in F[i]$ and $w\in\mm_F[i]$,
\[
 \Psi(x+h)=\Psi(x)\Ex(ih),\qquad E_\chi(z+w)=E_\chi(z)\Ex(w)=E_\chi(z)\sum_{n\ge0}\frac{w^n}{n!}.
\]
In the fine sense, with tolerances ranging over the positive elements of $F$,
$\Psi^{(n)}=i^n\Psi$ and $E_\chi^{(n)}=E_\chi$ for every ordinary $n$, and $\Psi$ is uniformly
continuous for the fine additive uniformity.
\end{proposition}
\begin{proof}
The identities are \cref{trigonometry:eq:globallocal} and its exponential form, on $F$. For the fine
derivative write $\Ex(w)=1+w+w^2R(w)$; for infinitesimal $w$ the factor $R(w)$ is finite with
standard part $1/2$, so $|R(w)|<1$. Given a positive tolerance $\epsilon\in F$, choose a positive
infinitesimal $\varrho\in F$; then
$0<|w|<\min\{\varrho,\epsilon/(1+|E_\chi(z)|)\}$ gives $|(E_\chi(z+w)-E_\chi(z))/w-E_\chi(z)|\le|E_\chi(z)||w|<\epsilon$, and induction
gives the higher derivatives; the phase is treated in the same way. Finally
$|\Psi(x+h)-\Psi(x)|=|\cis h-1|$ does not depend on $x$, and $(\cis h-1)/h$ is finite with standard
part $i$, so $|h|<\min\{\varrho,\epsilon/2\}$ forces $|\Psi(x+h)-\Psi(x)|<\epsilon$.
\end{proof}
This is source~20's Proposition 10.1 and source~21's Propositions 3.6 and 8.3; it is inherited, not a
new uniqueness theorem. The sums are strong Hahn sums, not limits of finite partial sums in the fine
topology. \emph{The local law does not remove the freedom}: for two phases $\Psi_1,\Psi_2$ the
quotient $\Psi_1\Psi_2^{-1}$ is a homomorphism trivial on $\OO_F$, hence constant on every coset of
$\OO_F$; those cosets are fine-open, so the quotient is locally constant and invisible to
differentiation (source~20). The derivative rules are the same for all extensions; they do not say
that the values of distinct exponentials agree at infinite points.

\subsection{Profinite residues and the profinite defect}\label{trigonometry:per:sec:defect}
Let $\Zhat=\varprojlim\Z/n\Z\cong\prod_p\Z_p$, with $\Z$ embedded diagonally. Throughout,
$\Zhat/\Z$ is an \emph{abstract} additive group, not a separated topological quotient ($\Z$ is dense
in $\Zhat$).

\begin{lemma}\label{trigonometry:per:lem:zhat}
$\Zhat/\Z$ is uniquely divisible, hence a $\Q$-vector space, of cardinality and $\Q$-dimension
$2^{\aleph_0}$.
\end{lemma}
\begin{proof}
Multiplication by $n>0$ is injective on $\Zhat$ with image the kernel of reduction modulo $n$; for
$\alpha\in\Zhat$ choose $r\in\Z$ with $\alpha\equiv r$ modulo $n$, so $(\alpha-r)/n\in\Zhat$. If
$n\alpha=m\in\Z$ then $n\mid m$ and $\alpha=m/n\in\Z$, so the quotient is torsion free. $\Zhat$ has
cardinality $2^{\aleph_0}$, a countable quotient leaves it unchanged, and an uncountable
$\Q$-vector space has cardinality equal to its dimension.
\end{proof}

Let $\cI=\cI(\Psi)$. For $a\in\cI$, $\Psi(2\pi a/n)$ is an $n$th root of unity, hence one of the
ordinary ones, $\Psi(2\pi a/n)=e^{2\pi ir_n(a)/n}$ with $r_n(a)\in\Z/n\Z$. These residues are compatible
under divisibility, so $\res_{\cI}(a)=(r_n(a))_n\in\Zhat$, and
\begin{equation}\label{trigonometry:per:eq:residue}
 a-r\in n\cI\iff r_n(a)\equiv r\pmod n .
\end{equation}
\begin{proposition}[The maximal divisible subgroup; classical]\label{trigonometry:per:prop:divisible}
$\res_{\cI}$ is additive and restricts to the diagonal embedding on $\Z$, and
\[
 \cI^{\mathrm{div}}:=\ker\res_{\cI}=\bigcap_{n\ge1}n\cI=\{a\in\cI:\Q a\subseteq\cI\}
\]
is the maximal divisible subgroup of $\cI$.
\end{proposition}
\begin{proof}
Additivity is immediate and \cref{trigonometry:per:eq:residue} gives the intersection formula. In the
uniquely divisible ambient group, $a\in n\cI$ for every $n$ says $a/n\in\cI$ for every $n$, i.e.\
$\Q a\subseteq\cI$; every divisible subgroup of $\cI$ has this property.
\end{proof}
This is the residue map and divisible kernel of the $\Z$-group $\cI$ \cite[Proposition 4.1 and
Remark 4.2]{Jerabek}, read off from a phase; source~21 presents it so.

\begin{theorem}[The profinite defect and its kernel; source 21]\label{trigonometry:per:thm:defect}
For $p\in\Pi_F$ choose $a\in\cI$ with $\pr(a)=p$ and put $\pd_\Psi(p)=\res_{\cI}(a)+\Z\in\Zhat/\Z$.
This is well defined and $\Q$-linear, and $\pr$ restricts to an isomorphism
$\cI^{\mathrm{div}}\cong\ker\pd_\Psi$. Writing $q:\Zhat\to\Zhat/\Z$,
\[
 \res_{\cI}(\cI)=q^{-1}\bigl(\pd_\Psi(\Pi_F)\bigr),\qquad
 \cI/\cI^{\mathrm{div}}\cong q^{-1}\bigl(\pd_\Psi(\Pi_F)\bigr)\subseteq\Zhat,
\]
so the nondivisible quotient is represented in a set of cardinality at most $2^{\aleph_0}$. For
set-sized $F$, $\Pi_F/\pr(\cI^{\mathrm{div}})\cong\pd_\Psi(\Pi_F)$ has $\Q$-dimension at most
$2^{\aleph_0}$; for $F=\No$ the same quotient is represented by the class map $\pd_\Psi$ with set-sized
range, without forming a set of proper classes.
\end{theorem}
\begin{proof}
Two lifts differ by an element of $\cI\cap\OO_F=\Z$, whose residue is diagonal; lifts of sums may be
added; source and target are uniquely divisible. A divisible subgroup meets $\Z$ trivially, so $\pr$
is injective on $\cI^{\mathrm{div}}$, and its image lies in $\ker\pd_\Psi$. If $\pd_\Psi(p)=0$, a lift
$a$ has $\res_{\cI}(a)=k\in\Z$ and $a-k\in\cI^{\mathrm{div}}$ has principal part $p$. If
$q(\alpha)=\pd_\Psi(p)$ and $a$ lifts $p$, then $\alpha=\res_{\cI}(a+k)$ for some $k\in\Z$; the
isomorphism is the first isomorphism theorem, read in the class case as the displayed residue map,
whose fibres are the cosets of $\cI^{\mathrm{div}}$.
\end{proof}

\begin{theorem}[When the period sequence splits]\label{trigonometry:per:thm:split}
For a phase $\Psi=\Psi_{\psi,\cL}$ with $\chi=\Psi|_{\Pi_F}$ the following are equivalent:
\begin{enumerate}[label=(\arabic*)]
\item $\pd_\Psi=0$;
\item the sequence \eqref{trigonometry:per:eq:periodsequence} splits as additive groups;
\item there is an additive retraction $\mathfrak r:F\to\OO_F$, the identity on $\OO_F$, with
 $\Psi(2\pi x)=\cis\bigl(2\pi\mathfrak r(x)\bigr)$ for all $x$;
\item there is an additive $H:\Pi_F\to\OO_F$ with $\chi(2\pi p)=\cis\bigl(2\pi H(p)\bigr)$;
\item the ordinary character $\psi$ has a $\Q$-linear lift $\Pi_F\to\R$ through $\R\to\R/\Z$.
\end{enumerate}
The retraction in (3) is unique when it exists, $\cI=\ker\mathfrak r\oplus\Z$, and
$\cI^{\mathrm{div}}=\ker\mathfrak r$; every additive $H$ as in (4) is $\Q$-linear.
\end{theorem}
\begin{proof}
(1)$\Rightarrow$(3): the isomorphism of \cref{trigonometry:per:thm:defect} gives a section
$\iota:\Pi_F\to\cI^{\mathrm{div}}$; $\mathfrak r(x)=x-\iota(\pr x)$ is finite, additive, fixes
$\OO_F$, and differs from $x$ by a period. (3)$\Rightarrow$(2): $p\mapsto p-\mathfrak r(p)$ is a section.
(2)$\Rightarrow$(1): an additive section has divisible image, hence lies in $\cI^{\mathrm{div}}$, and
it meets every principal part. If $\mathfrak r_1,\mathfrak r_2$ give the same phase, $\mathfrak
r_1-\mathfrak r_2$ is an additive map from the divisible group $F$ to $\Z$, hence zero; $\mathfrak
r(x)\in\Z$ characterizes $\cI$, and all rational divisions of $x$ lie in $\cI$ exactly when
$\mathfrak r(x)=0$.

(2)$\Leftrightarrow$(4) (source 20's route): given $H$, $p\mapsto p-H(p)$ is an additive section into
$\cI$; given a section $\iota$, $H(p)=p-\iota(p)$ is finite and additive with $\chi(2\pi p)=\cis(2\pi H(p))$.
Additive maps between torsion-free divisible groups are $\Q$-linear. (4)$\Leftrightarrow$(5): write
$H=A+\mathcal E$ with $A$ real and $\mathcal E$ infinitesimal; by the splitting of $\TT(F)$,
\cref{trigonometry:per:eq:coordinates} and $\chi(2\pi p)=\cis(2\pi H(p))$ are equivalent to
$A+\Z=\psi$ and $\mathcal E=\cL$. The infinitesimal coordinate always lifts; all obstruction lies in
$\psi$.
\end{proof}
This is source~21's Theorem 4.5 (conditions (1)--(3)) and source~20's Theorem 7.1 (conditions (2),
(4), (5)). A real argument for each single value of $\chi$ always exists; condition (5) asks for them
\emph{simultaneously and additively}. Replacing an arbitrary character by a real-valued linear
functional therefore discards every nonzero profinite defect, while a zero defect says nothing about
multiplicative closure (\cref{trigonometry:per:sec:multiplier}).

\begin{theorem}[Homomorphisms to $\Z$; source 20]\label{trigonometry:per:thm:homz}
The sequence \eqref{trigonometry:per:eq:periodsequence} splits if and only if there is a nonzero
homomorphism $\cI\to\Z$. For set-sized $F$, $\Hom(\cI,\Z)\cong\Z$ if it splits and $\Hom(\cI,\Z)=0$ if
it does not. For $F=\No$ the statement concerns existence and pointwise description of class
homomorphisms, not a set of all of them.
\end{theorem}
\begin{proof}
If it splits, $\cI\cong\Z\oplus\Pi_F$, and every homomorphism from the divisible group $\Pi_F$ to $\Z$
vanishes, so the homomorphisms are the multiples of the first projection. Conversely let
$\varphi:\cI\to\Z$ be nonzero and $m=\varphi(1)$; if $m=0$, $\varphi$ factors through
$\cI/\Z\cong\Pi_F$ and vanishes. By \cref{trigonometry:per:prop:division} with $n=|m|$, $d=ne+r$
gives $\varphi(d)=n\varphi(e)+rm\in m\Z$, so $\varphi/m$ is a retraction of $\Z\hookrightarrow\cI$, and
$\pr$ maps its kernel bijectively onto $\Pi_F$.
\end{proof}
The group facts in \cref{trigonometry:per:thm:split,trigonometry:per:thm:homz} are classical for
$\Z$-groups \cite[\S\S2--4]{Jerabek}; source~20 proves them without that citation.

\subsection{Realizing every profinite defect}\label{trigonometry:per:sec:realize}
Let $\Qv=\Hom(\Q,\R/\Z)$, an abstract group (the character group behind the solenoid; no topology is
used), with rational scalars $(r\cdot f)(q)=f(rq)$.

\begin{lemma}[A rational-line exact sequence]\label{trigonometry:per:lem:solenoid}
There is an exact sequence of $\Q$-vector spaces
$0\to\R\xrightarrow{j}\Qv\xrightarrow{\cl}\Zhat/\Z\to0$ with $j(a)(q)=qa+\Z$.
\end{lemma}
\begin{proof}
$j$ is injective. For $f\in\Qv$ choose $\theta\in\R$ representing $f(1)$ and $g=f-j(\theta)$; then
$g(1/n)=a_n/n+\Z$ with compatible $a_n\in\Z/n\Z$, i.e.\ $\alpha=(a_n)\in\Zhat$; put $\cl f=\alpha+\Z$,
independent of $\theta$ and additive, hence linear. If $\cl f=0$ then $\alpha=k\in\Z$ and
$f=j(\theta+k)$. For any $\alpha\in\Zhat$ the rule
\begin{equation}\label{trigonometry:per:eq:falpha}
 f_\alpha(k/n)=\frac{ka_n}{n}+\Z
\end{equation}
is well defined and additive (common denominators use exactly the compatibility), with $f_\alpha(1)=0$
and $\cl f_\alpha=\alpha+\Z$.
\end{proof}
By choice between sets fix a linear section $\varsigma:\Zhat/\Z\to\Qv$ of $\cl$; this is a choice
between set-sized vector spaces even when the defect below is defined on all of $\Pi$.

\begin{lemma}[The defect is read from the ordinary character]\label{trigonometry:per:lem:hdefect}
For $\psi:\Pi_F\to\R/\Z$ put $H_\psi(p)(q)=\psi(qp)$. Then $H_\psi:\Pi_F\to\Qv$ is $\Q$-linear and
$\pd_{\Psi_{\psi,\cL}}=\cl\circ H_\psi$ for every $\cL$; in particular the profinite defect does not
depend on $\cL$.
\end{lemma}
\begin{proof}
Linearity is additivity of $\psi$. Choose $\theta\in\R$ representing $\psi(p)$; then
$a=p-\theta-\cL(p)\in\cI$ with $\pr(a)=p$, and for every $n$, since $\cL$ is linear,
$\Psi(2\pi a/n)=e^{2\pi i(\psi(p/n)-\theta/n)}$: the infinitesimal factors cancel. These are the residues
defining $\cl(H_\psi(p))$.
\end{proof}

\begin{theorem}[Prescribed-defect realization; source 21]\label{trigonometry:per:thm:realize}
Every $\Q$-linear $d:\Pi_F\to\Zhat/\Z$ is the profinite defect of a phase, also for a given
class-linear $d$ on $\Pi$ when $F=\No$, without a class Hamel basis. After fixing a realizing
ordinary character $\psi$, the infinitesimal character $\cL$ may be chosen freely.
\end{theorem}
\begin{proof}
Put $H=\varsigma\circ d$ and $\psi(p)=H(p)(1)$. Then $\psi(qp)=H(qp)(1)=(q\cdot H(p))(1)=H(p)(q)$, so
$H_\psi=H$ and $\pd_{\Psi_{\psi,\cL}}=\cl\,\varsigma\,d=d$ for every linear $\cL$, by
\cref{trigonometry:per:lem:hdefect}. The only choice is the set map $\varsigma$.
\end{proof}

\begin{theorem}[Counts on $F_{\Q}$; source 21, ZFC]\label{trigonometry:per:thm:cardinality}
On $F_{\Q}=\R((t^{\Q}))$: every cardinal $\kappa\le2^{\aleph_0}$ is $\dim_{\Q}\pd_\Psi(\Pi_F)$ for some
phase; there are exactly $2^{2^{\aleph_0}}$ phases with \emph{reduced} period group
($\cI^{\mathrm{div}}=0$) and exactly $2^{2^{\aleph_0}}$ with split period group; the same counts hold
for their distinct normalized period groups.
\end{theorem}
\begin{proof}
All supports lie in the countable $\Q$, so $|F_{\Q}|=2^{\aleph_0}$; $\Pi_F\supseteq\R\omega$ gives
$\dim_{\Q}\Pi_F=2^{\aleph_0}=\dim_{\Q}\Zhat/\Z$ (\cref{trigonometry:per:lem:zhat}), so linear maps
of every rank exist and are realized by \cref{trigonometry:per:thm:realize}. A linear isomorphism
$d$ gives a reduced group by \cref{trigonometry:per:thm:defect}; multiplying the images of a basis
independently by $1$ or $2$ gives $2^{2^{\aleph_0}}$ of them, with distinct defects, phases and (by
\cref{trigonometry:per:thm:dictionary}) groups. For split phases take $\Q$-linear $A:\Pi_F\to\R$ and
$\Psi_A(2\pi(p+f))=\cis(2\pi(f+A(p)))$; distinct $A$ give distinct phases, their difference mapping the
divisible $\Pi_F$ into $\Z$. There are at most $(2^{\aleph_0})^{2^{\aleph_0}}=2^{2^{\aleph_0}}$ maps
$F_{\Q}\to F_{\Q}[i]$.
\end{proof}
No continuum hypothesis or regularity of $2^{\aleph_0}$ is used. A reduced group still has infinite
elements (every additive integer part of a non-Archimedean field has them), and these are counts of
embedded period groups in a fixed field, not a claim that their abstract groups are pairwise
nonisomorphic.

\begin{example}[Two explicit nonliftable profinite characters]\label{trigonometry:per:ex:profinite}
\emph{Source 21.} Let $\alpha_2\in\Zhat$ have $2$-adic component $0$ and every odd component $1$; its
residue $a_n$ for $n=2^em$, $m$ odd, solves $a_n\equiv0\pmod{2^e}$, $a_n\equiv1\pmod m$:
$a_2=0$, $a_3=1$, $a_6=4$, $a_{12}=4$, $a_{24}=16$, $a_{40}=16$. It is not a diagonal integer, so
$[\alpha_2]\ne0$, and $f_{\alpha_2}$ vanishes on integers and dyadic rationals but not at $1/3$. On $\No$,
or on $F_\Gamma$ with $1\in\Gamma$, choose a $\Q$-linear $\nu:\R\to\Q$ fixing $\Q$ (a Hamel-basis choice
on which nothing displayed depends) and put $\psi(p)=f_{\alpha_2}(\nu(\ell(p)))$, $\cL=0$. Then
\[
 \Psi(2\pi\omega)=1,\quad \Psi(2\pi\omega/2^e)=1,\quad \Psi(2\pi\omega/3)=e^{2\pi i/3},\quad
 \pd_\Psi(\omega)=[\alpha_2]\ne0 ,
\]
so a period need not bring its rational divisions with it. With $u=(\omega-1)/3$ one has
$\Psi(2\pi u)=\Psi(2\pi\omega)=1$ but $\Psi\bigl(2\pi u\omega\bigr)=e^{-2\pi i/3}\ne1$, since
$u\omega=(\omega^2-\omega)/3$: this period group is not a ring, and it is not split. The example
establishes neither a minimal multiplier ring nor reducedness.

\emph{Source 20.} Let $\alpha_3\in\Zhat$ have $3$-adic component $1$ and every other component $0$:
for $n=3^km$ with $3\nmid m$, $a_n\equiv0\pmod m$ and $a_n\equiv1\pmod{3^k}$. The first residues are
\[
\begin{array}{c|rrrrrrrrrr}
 n&2&3&4&6&8&9&12&18&24&27\\\hline
 a_n&0&1&0&4&0&1&4&10&16&1
\end{array}
\]
and $f_{\alpha_3}(1)=0$, $f_{\alpha_3}(1/2^k)=0$, $f_{\alpha_3}(1/3^k)=1/3^k+\Z$. \emph{It does not lift
to an additive map $\Q\to\R$}: a lift is $q\mapsto qf(1)$ with $f(1)\in\Z$ divisible by every $2^k$,
hence zero, which cannot represent $f_{\alpha_3}(1/3)$. (In general $f_\alpha$ lifts exactly when
$[\alpha]=0$, by \cref{trigonometry:per:lem:solenoid}.) The roots are chosen coherently --- the cube
of the value at $1/3$ is the value at $1$ --- but not by dividing one additive angle.
\end{example}

\subsection{Divisible periods are unavoidable and cofinal on \texorpdfstring{$\No$}{No}}
\label{trigonometry:per:sec:cofinal}
\begin{theorem}[Cofinal rationally divisible periods; source 21]\label{trigonometry:per:thm:cofinal}
Let $\Psi$ be any phase on $\No$. For every $B\in\No$ there is $a>B$ with $\Psi(2\pi qa)=1$ for every
$q\in\Q$. Equivalently, $\cI(\Psi)^{\mathrm{div}}$ is cofinal in $\No$.
\end{theorem}
\begin{proof}
We may take $B\ge0$. Choose $g>0$ with $\omega^g>2B+\omega$, let $\kappa=(2^{\aleph_0})^+$ and
$p_\beta=\omega^{g+\beta}$ for ordinals $\beta<\kappa$. The values $\pd_\Psi(p_\beta)$ lie in the set
$\Zhat/\Z$ of cardinality $2^{\aleph_0}$, so $\pd_\Psi(p_\beta)=\pd_\Psi(p_\gamma)$ for some
$\beta<\gamma$. Since $\gamma-\beta\ge1$, $p_\gamma\ge\omega p_\beta$ and
$p=p_\gamma-p_\beta\ge(\omega-1)\omega^g>B+\omega$. As $p\in\ker\pd_\Psi$,
\cref{trigonometry:per:thm:defect} gives $a\in\cI^{\mathrm{div}}$ with $\pr(a)=p$; $a-p$ is finite,
so $a>B$.
\end{proof}
The argument is one set-sized pigeonhole, not cardinal arithmetic with a purported cardinality of
$\No$. It strengthens \cref{trigonometry:thm:infiniteperiods}: not only does every coset of $\OO_{\R}$
contain a period, but whole rational lines of periods reach beyond every bound. It does not assert
that strong sums are fine limits.

\begin{corollary}\label{trigonometry:per:cor:noreduced}
No phase on $\No$ has a reduced period group. The reduced examples on $F_{\Q}$ of
\cref{trigonometry:per:thm:cardinality} therefore cannot stay reduced in any extension to $\No$
(this is conditional only on an extension being given; nothing is asserted about extending an
arbitrary class character).
\end{corollary}
The defect rank is nevertheless unrestricted on $\No$: for $\kappa\le2^{\aleph_0}$ take a rank-$\kappa$
linear $A:\R\to\Zhat/\Z$ and $d(p)=A(\ell(p))$; \cref{trigonometry:per:thm:realize} realizes it, and
its divisible periods are still cofinal. \emph{A set/class contrast}: on the set field $F_{\Q}$ period
groups can be reduced; on the class $\No$ they must contain cofinally large rational lines.

\subsection{Multiplier rings and the integer-part criterion}\label{trigonometry:per:sec:multiplier}
For an additive integer part $\cI$ put $\Mult(\cI)=\{a\in F:a\cI\subseteq\cI\}$. It is a unital subring
containing $\Z$ and contained in $\cI$ (apply it to $1$), hence discretely ordered with least positive
element $1$. For an exponential $E$ as in \cref{trigonometry:per:prop:exponentials} the
\emph{kernel-multiplier ring} of Ehrlich and Kaplan is
\begin{equation}\label{trigonometry:per:eq:ZE}
 Z_E=\{z\in F[i]:\ E(y)=1\Rightarrow E(yz)=1\ \text{for every }y\in F[i]\}.
\end{equation}

\begin{proposition}[The kernel-multiplier ring]\label{trigonometry:per:prop:ZE}
For $E=E_\chi$ with $\cI=\cI(\Psi_\chi)$, $Z_E=\Mult(\cI)\subseteq F$.
\end{proposition}
\begin{proof}
$2\pi i\in\ker E$, so $z=u+iv\in Z_E$ gives $E(2\pi iz)=1$, whose modulus $\exp_F(-2\pi v)$ forces
$v=0$; for real $u$, preserving every $2\pi ia$ with $a\in\cI$ is $u\cI\subseteq\cI$.
\end{proof}

\begin{theorem}[The integer-part criterion]\label{trigonometry:per:thm:criterion}
For an additive integer part $\cI$ of $F$ the following are equivalent:
(i) $\Mult(\cI)$ is an integer part of $F$; (ii) $\Mult(\cI)=\cI$; (iii) $\cI$ is a subring.
If $\cI=\cI(\Psi_\chi)$ and $F$ has a real exponential, they are also equivalent to (iv) $Z_{E_\chi}$
is an integer part, and to (v) the binary kernel condition
\begin{equation}\label{trigonometry:per:eq:binary}
 \forall u,v\in F\qquad \bigl(E_\chi(2\pi iu)=E_\chi(2\pi iv)=1\bigr)\Longrightarrow E_\chi(2\pi iuv)=1 .
\end{equation}
If $F$ is initial in $\No$ ($F=\No$ or $\No(\lambda)$) and $Z_{E_\chi}$ is initial, these conditions
hold if and only if $E_\chi$ is the canonical exponential $\Exp$ on $F[i]$.
\end{theorem}
\begin{proof}
If (i), take $a\in\cI$ and a $\Mult(\cI)$-floor $m$ of $a$; $m\in\cI$ gives $a-m\in\cI\cap[0,1)=\{0\}$,
so $\cI\subseteq\Mult(\cI)$, which is (ii). A ring equal to $\cI$ makes $\cI$ a ring. If $\cI$ is a ring,
each of its elements preserves it, so $\Mult(\cI)=\cI$, an integer part by the floor property.
(iv) is (i) by \cref{trigonometry:per:prop:ZE}, and (v) is (iii) written with $E_\chi$. If $F$ and
$Z_{E_\chi}$ are initial and the conditions hold, the uniqueness of the initial integer part
\cite[Lemma 11.1]{EK} gives $\cI=Z_{E_\chi}=\Oz\cap F$, and the kernel dictionary gives $\Psi_\chi=\Phi$.
Conversely $\Exp$ has $\cI=\Oz\cap F$, an initial integer part equal to its multiplier ring.
\end{proof}
This is source~20's Lemma 5.1 and Theorem 5.2 and source~21's Propositions 7.1, 8.2 and equation
(8.7); the printed proof is source~20's. In source~21's formulation: with
$P_E=\{a\in F:E(2\pi ia)=1\}$, $Z_E$ is an integer part if and only if $P_E$ is closed under
multiplication. \emph{The criterion is an exact reduction, not a classification of every multiplier
ring, and not a substitute for independent analytic or geometric hypotheses implying its right-hand
side}; merely assuming that side would restate the missing content rather than explain it.

\begin{remark}[Ring defects of a linear phase; source 20]\label{trigonometry:per:rem:ringdefect}
For a $\Q$-linear $A:\Pi_F\to\R$ and the split phase with coordinates $(A+\Z,0)$ the period group is
$\cI_A=\{p-A(p)+n:p\in\Pi_F,\ n\in\Z\}$. For $d_p=p-A(p)$ and $d_q=q-A(q)$ the product has principal
part $pq-A(p)q-A(q)p$ and constant part $A(p)A(q)$, so
\[
 d_pd_q\in\cI_A\iff A\bigl(pq-A(p)q-A(q)p\bigr)+A(p)A(q)\in\Z .
\]
This holds for merely $\Q$-linear $A$; the generally irrational numbers $A(p),A(q)$ must not be pulled
through $A$.
\end{remark}

\begin{example}[A non-ring kernel whose multiplier ring is not initial; source 20]
\label{trigonometry:per:ex:noninitial}
Let $A_2(p)$ be half the real coefficient of $\omega^2$ in $p$, and take the phase of
\cref{trigonometry:per:rem:ringdefect}. Then $\Psi(2\pi\omega)=1$ and $\Psi(2\pi\omega^2)=-1$, so
$\cI_{A_2}\ni\omega$ but $\omega^2\notin\cI_{A_2}$: the kernel is not a ring. Multiplication by
$\omega^3$ sends every element $p-A_2(p)+n$ to an element with no finite part and zero coefficient at
$\omega^2$, which lies in $\cI_{A_2}$, so $\omega^3\in\Mult(\cI_{A_2})$ while $\omega^2\notin
\cI_{A_2}\supseteq\Mult(\cI_{A_2})$. The ordinal $\omega^2$ is a sign prefix of $\omega^3$, so this
multiplier ring is \emph{not initial}. Initiality cannot be tested from one infinite-angle example;
the constructions below prevent every nonintegral multiplier.
\end{example}

\subsection{Minimal multiplier rings}\label{trigonometry:per:sec:minimal}
\emph{Source 21's route: fix the ordinary character, perturb the infinitesimal one.} Fix an additive
$\psi:\Pi_F\to\R/\Z$ and let $\bar\psi(p)\in[0,1)$ be the ordinary representative of $\psi(p)$ (not
additive). The infinitesimal character $\cL$ is built linearly.

\begin{lemma}[One extension step]\label{trigonometry:per:lem:onestep}
Let $U\subseteq\Pi_F$ be a $\Q$-subspace, $\cL_U:U\to\mm_F$ linear, and $a\in F$ positive infinite.
Suppose $p\in\Pi_F$ is such that $p$ and $y:=\pr\bigl(a(p-\bar\psi(p))\bigr)$ are linearly
independent modulo $U$. Then $\cL_U$ extends to $U+\Q p+\Q y$ so that for \emph{every} later linear
extension $\cL$ the phase $\Psi_{\psi,\cL}$ has a normalized period $b$ with $ab$ not a normalized
period.
\end{lemma}
\begin{proof}
Fix $0\ne\eps_*\in\mm_F$ ($F$ is non-Archimedean), let $b=p-\bar\psi(p)$ and $f=\fin(ab)$, and
prescribe $\cL(p)=0$, $\cL(y)=\eps_*-(f-\st f)$, consistent by independence. For every later
extension, $\Psi(2\pi b)=e^{2\pi i(\psi(p)-\bar\psi(p))}=1$, while $ab=y+f$ gives
$\Psi(2\pi ab)=e^{2\pi i(\psi(y)+\st f)}\Ex(2\pi i\eps_*)$, whose infinitesimal angle $2\pi\eps_*$ is
nonzero, so it is not $1$. Both conclusions depend only on the two prescribed values.
\end{proof}

\begin{lemma}[Disjoint blocks]\label{trigonometry:per:lem:blocks}
If a $\Q$-vector space contains two-dimensional subspaces $(W_j)_{j\in J}$ whose sum is direct and
$\dim_{\Q}U<|J|$, then some $W_j$ meets $U$ only in $0$.
\end{lemma}
\begin{proof}
Otherwise nonzero $u_j\in W_j\cap U$ would be $|J|$ independent vectors of $U$.
\end{proof}

\begin{lemma}[Fresh pairs in a continuum-sized field]\label{trigonometry:per:lem:setpairs}
Let $F$ be a non-Archimedean truncation-closed subfield of $\No$ containing $\R$ with
$|F|=2^{\aleph_0}$, let $a\in F$ be positive infinite, and let $U\subseteq\Pi_F$ have
$\dim_{\Q}U<2^{\aleph_0}$. Then some $p$ satisfies \cref{trigonometry:per:lem:onestep}.
\end{lemma}
\begin{proof}
Truncation closure gives a positive infinite monomial $X\in F$. Put $Y=\pr(aX)$ and $P_a=\pr(a)$; if the
leading $\omega$-exponents of $a$ and $X$ are $\alpha,\gamma>0$, that of $Y$ is $\alpha+\gamma$, above
those of $P_a$ and $X$. For a $\Q$-basis $(r_j)_{j\in J}$ of $\R$, $|J|=2^{\aleph_0}$, put $p_j=r_jX$ and
$y_j=\pr(a(p_j-\bar\psi(p_j)))=r_jY-\bar\psi(p_j)P_a$. In a finite relation
$(\sum v_jr_j)Y-(\sum v_j\bar\psi(p_j))P_a+(\sum u_jr_j)X=0$ the leading exponent of $Y$ forces
$\sum v_jr_j=0$, so all $v_j=0$, and then all $u_j=0$. Apply
\cref{trigonometry:per:lem:blocks} to $W_j=\Q p_j+\Q y_j$.
\end{proof}

\begin{theorem}[Minimal multiplier ring on a continuum-sized field; source 21, ZFC]
\label{trigonometry:per:thm:minimalset}
Let $F$ be as in \cref{trigonometry:per:lem:setpairs}, with the finite phase. For every additive
$\psi:\Pi_F\to\R/\Z$ there is a $\Q$-linear $\cL:\Pi_F\to\mm_F$ with
$\Mult\bigl(\cI(\Psi_{\psi,\cL})\bigr)=\Z$. This applies to $F_{\Q}$ and to $\No(\eps_0)$.
\end{theorem}
\begin{proof}
Let $\lambda_{\mathfrak c}$ be the initial ordinal of $2^{\aleph_0}$; enumerate $\Pi_F$ as
$(u_\xi)_{\xi<\lambda_{\mathfrak c}}$ and the positive infinite elements of $F$ as
$(a_\xi)_{\xi<\lambda_{\mathfrak c}}$, with repetitions allowed. Build increasing partial linear maps
$\cL_\xi:U_\xi\to\mm_F$: at stage $\xi$ add $u_\xi$ with value $0$ if it is new (the domain then has
dimension at most $\max(\aleph_0,|\xi|)<2^{\aleph_0}$), find a pair for $a_\xi$ by
\cref{trigonometry:per:lem:setpairs}, extend by \cref{trigonometry:per:lem:onestep}; at limits take
unions. The union $\cL$ is total. Every positive infinite $a$ has a period $b$ with $ab$ not a period;
a negative infinite multiplier would give a positive one; a finite multiplier lies in
$\cI\cap\OO_F=\Z$; and $\Z$ always multiplies $\cI$ into itself. Below $\lambda_{\mathfrak c}$ every
stage has $|\xi|<2^{\aleph_0}$ even if $2^{\aleph_0}$ is singular: no regularity is needed.
\end{proof}

\begin{lemma}[Fresh pairs outside a set-sized domain]\label{trigonometry:per:lem:classpairs}
Let $a>\R$ be surreal, $\psi:\Pi\to\R/\Z$ a class character, and $U\subseteq\Pi$ a set-sized
subspace. Then a pair as in \cref{trigonometry:per:lem:onestep} exists.
\end{lemma}
\begin{proof}
Let $\alpha>0$ be the leading exponent of $a$ and $\kappa>\max(|U|,\aleph_0)$ a set cardinal. The
exponents $\gamma_\beta=1+\alpha/(4(\beta+1))$, $\beta<\kappa$ (ordinals read as surreals), are
distinct and lie in $(1,1+\alpha/2)$. Put $p_\beta=\omega^{\gamma_\beta}$ and
$y_\beta=\pr(a(p_\beta-\bar\psi(p_\beta)))$, whose leading exponent is $\alpha+\gamma_\beta$, since
$a\bar\psi(p_\beta)$ has exponents at most $\alpha$. These leading exponents are distinct and exceed
every $\gamma_{\beta'}$ and $\alpha$, so the blocks $\Q p_\beta+\Q y_\beta$ are in direct sum;
\cref{trigonometry:per:lem:blocks} finishes.
\end{proof}

\begin{theorem}[Minimal multiplier ring on $\No$; source 21, GBC]\label{trigonometry:per:thm:minimalclass}
In NBG with global choice, for every class character $\psi:\Pi\to\R/\Z$ there is a class-linear
$\cL:\Pi\to\mm_{\R}$ with $\Mult\bigl(\cI(\Psi_{\psi,\cL})\bigr)=\Z$.
\end{theorem}
\begin{proof}
Global choice gives set-like ordinal enumerations of $\Pi$ and of the positive infinite surreals. Run
the recursion of \cref{trigonometry:per:thm:minimalset} along the ordinals with
\cref{trigonometry:per:lem:classpairs}. Every partial domain is a set (finitely many vectors at a
successor, a set of earlier stages at a limit); the choices of sets and set functions use the fixed
global well-order, so this is set-valued transfinite recursion with a class parameter; no class-valued
stage or class Hamel basis occurs. The finite and negative cases are as before.
\end{proof}

\begin{corollary}[Any defect together with the minimal ring; source 21]\label{trigonometry:per:cor:simultaneous}
On the fields of \cref{trigonometry:per:thm:minimalset} (ZFC) and on $\No$ (GBC), every $\Q$-linear
$d:\Pi_F\to\Zhat/\Z$ is the profinite defect of a phase with multiplier ring $\Z$; in particular
$\pd_\Psi=0$ and $\Mult(\cI(\Psi))=\Z$ occur together. Throughout the construction the standard part
$\st\Psi(2\pi(p+f))=e^{2\pi i(\psi(p)+\st f)}$ of the phase is fixed at every argument.
\end{corollary}
\begin{proof}
Realize $d$ by $\psi$ (\cref{trigonometry:per:thm:realize}) and apply the minimal-multiplier theorem
to that $\psi$; the defect does not depend on $\cL$ (\cref{trigonometry:per:lem:hdefect}), nor does the
displayed standard part. Take $\psi=0$ for the last example.
\end{proof}
For $\psi=0$ the phase is $\cis\circ2\pi\mathfrak r_{\cL}$ with $\mathfrak r_{\cL}(p+f)=f+\cL(p)$, and
its period group is the split group $\{p-\cL(p)+n\}$, whose divisible part is the whole graph of $-\cL$;
yet no nonintegral scalar preserves it. Division-compatible periods are abundant while universal
multiplicative symmetries are minimal.

\emph{Source 20's route: prescribe a set-sized part of the ordinary character, keep $\cL=0$.} This
route works on $\No$ and gives ordinary-circle values, which
\cref{trigonometry:per:sec:reflection} needs.

\begin{lemma}[Fresh ordinal monomial]\label{trigonometry:per:lem:fresh}
Let $W\subseteq\Pi$ be a set-sized $\Q$-subspace and $x\in\No\setminus\Q$. There is an ordinal
$\beta$ such that $\mu=\omega^\beta$ satisfies $\mu,x\mu\in\Pi$ with $q\mu+rx\mu\in W$
($q,r\in\Q$) only for $q=r=0$.
\end{lemma}
\begin{proof}
Let $S$ be the union of the supports of elements of $W$ and $G=\supp(x)\cup\{0\}$; the set
$\{s-g:s\in S\cup\{0\},\ g\in G\}$ is bounded above by an ordinal $\beta$ (a set of birthdays has an
ordinal bound). Then each $\beta+g$ (surreal addition) is positive and above all of $S$, so $\mu$ and
$x\mu$ are purely infinite with all exponents above $S$. If $q\mu+rx\mu\in W$, uniqueness of normal
forms makes both sides vanish, so $q+rx=0$, and $x\notin\Q$ gives $r=q=0$.
\end{proof}
No bound on the exponents of $x$ and no finiteness of its support is assumed; only that single supports
and the union of the supports of a set are sets.

\begin{theorem}[Relative minimal-multiplier extension; source 20, GBC]\label{trigonometry:per:thm:relative}
Let $W_0\subseteq\Pi$ be a set-sized $\Q$-subspace and $\psi_0:W_0\to\R/\Z$ a homomorphism. There is a
class homomorphism $\psi:\Pi\to\R/\Z$ extending $\psi_0$ such that the phase $\Psi=\Psi_{\psi,0}$ has,
for every $x\in\No\setminus\Q$, a positive purely infinite monomial $\mu_x$ with
\[
 \Psi(2\pi\mu_x)=1,\qquad \Psi(2\pi x\mu_x)=-1 .
\]
Hence $\Mult(\cI(\Psi))=\Z$. Likewise, a $\Q$-linear $A_0:W_0\to\R$ extends to a $\Q$-linear
$A:\Pi\to\R$ whose split phase (coordinates $(A+\Z,0)$) has $\Mult=\Z$ and period group
$\{p-A(p)+n\}$.
\end{theorem}
\begin{proof}
Fix a set-like global enumeration $(x_\alpha)_{\alpha\in\mathrm{On}}$ of $\No$ and build set-sized
subspaces $W_\alpha$ with compatible characters $\psi_\alpha$. At stage $\alpha$ adjoin
$\pr(x_\alpha)$ if new, with value $0$ on its rational line. If $x_\alpha\notin\Q$, apply
\cref{trigonometry:per:lem:fresh} to obtain $\mu_\alpha$, adjoin the direct summand
$\Q\mu_\alpha\oplus\Q x_\alpha\mu_\alpha$ and set
$\psi(w+q\mu_\alpha+rx_\alpha\mu_\alpha)=\psi(w)+r/2+\Z$. At limits take unions (set many stages). The
union is defined on all of $\Pi$ and keeps every stage assignment. If $x\notin\Q$ then $\mu_x\in\cI$
and $x\mu_x\notin\cI$, so $\Mult(\cI)\subseteq\Q\cap\cI=\Z$. For the split variant run the same
recursion with real values $A(\mu_\alpha)=0$, $A(x_\alpha\mu_\alpha)=1/2$; the phase at $2\pi(p+r+\eps)$
is $e^{2\pi i(A(p)+r)}\Ex(2\pi i\eps)$, which is $1$ exactly when $\eps=0$ and $A(p)+r\in\Z$, and
$p\mapsto p-A(p)$ is a section.
\end{proof}
The character on a new rational direction is specified on \emph{all} its rational multiples, and the
whole prescribed $\psi_0$ is retained; this is decisive below. Source~20 states the first part as its
Theorem 6.2 and the split variant as Theorem 6.4.

\begin{theorem}[A nonsplit minimal example; source 20, GBC]\label{trigonometry:per:thm:nonsplit}
There is a phase $\Psi_{\mathrm n}$ on $\No$ with ordinary-circle values, multiplier ring $\Z$, and
nonsplit period sequence, satisfying
\[
 \Psi_{\mathrm n}(2\pi\omega)=1,\qquad\Psi_{\mathrm n}(2\pi\omega/2^k)=1,\qquad
 \Psi_{\mathrm n}(2\pi\omega/3^k)=e^{2\pi i/3^k}\quad(k\ge1).
\]
\end{theorem}
\begin{proof}
Apply \cref{trigonometry:per:thm:relative} to $W_0=\Q\omega$ and $\psi_0(q\omega)=f_{\alpha_3}(q)$
(\cref{trigonometry:per:ex:profinite}). A splitting would give, by
\cref{trigonometry:per:thm:split}(5), a real additive lift of $f_{\alpha_3}$, which does not exist.
\end{proof}
\emph{Comparison of the routes.} Source~21's theorems keep $\psi$ fixed on all of $\Pi_F$, hence the
profinite defect and the standard part of the phase at every argument, and hold in ZFC on
$\No(\eps_0)$ and $F_{\Q}$; with \cref{trigonometry:per:cor:simultaneous} they also give split
($d=0$) and nonsplit ($d\ne0$) minimal examples. Their phase values have nonzero infinitesimal
angles, and nothing is claimed about their restriction to subfields. Source~20's theorem prescribes
$\psi$ only on a set-sized subspace, changes it elsewhere, keeps every value in the ordinary circle, and
excludes every nonrational multiplier, finite ones included, by an explicit witness; that is what makes
the restriction and reflection of \cref{trigonometry:per:sec:reflection} work.

\subsection{Surcomplex exponentials and the Ehrlich--Kaplan question}\label{trigonometry:per:sec:ek}
\begin{theorem}[Initial kernel-multiplier ring $\Z$]\label{trigonometry:per:thm:robustness}
\begin{enumerate}[label=(\alph*)]
\item \emph{(Source 21, ZFC.)} There is an exponential $E$ on $\No(\eps_0)[i]$ with all properties of
 \cref{trigonometry:per:prop:exponentials,trigonometry:per:prop:local} and $Z_E=\Z$. It can be chosen
 with
 \begin{equation}\label{trigonometry:per:eq:invisible}
  \frac{E(z)}{\Exp(z)}\in\TT\bigl(\No(\eps_0)\bigr)\cap(1+\mm[i])\qquad\text{for every }z ,
 \end{equation}
 and with zero profinite defect.
\item \emph{(Source 21, GBC.)} The same holds on $\No[i]$.
\item \emph{(Source 20, GBC.)} On $\No[i]$, for every set-sized $W_0\subseteq\Pi$ and homomorphism
 $\psi_0:W_0\to\R/\Z$, there is $E=E_\chi$ with ordinary-circle phase values, coordinates extending
 $(\psi_0,0)$, and $Z_E=\Z$; examples of this kind with split and with nonsplit period sequence exist
 (\cref{trigonometry:per:thm:relative,trigonometry:per:thm:nonsplit}).
\end{enumerate}
In every case $Z_E=\Z$ is an initial subring of the field but not an integer part.
\end{theorem}
\begin{proof}
For (a) and (b) apply \cref{trigonometry:per:thm:minimalset} or \cref{trigonometry:per:thm:minimalclass}
with $\psi=0$ and take $E=E_\chi$ for $\Psi_\chi=\Psi_{0,\cL}$. Then $Z_E=\Z$ by
\cref{trigonometry:per:prop:ZE}; the canonical phase is $\Psi_{0,0}$ because $\Oz=\Pi\oplus\Z$, so
$E(a+ib)/\Exp(a+ib)=\Ex\bigl(2\pi i\cL(\pr(b/2\pi))\bigr)$, of modulus one and standard part one; the
defect vanishes by \cref{trigonometry:per:lem:hdefect}. Part (c) is
\cref{trigonometry:per:thm:relative,trigonometry:per:thm:nonsplit} with
\cref{trigonometry:per:prop:ZE}. Every proper sign prefix of an ordinary integer is an ordinary integer,
so $\Z$ is initial; no ordinary integer lies in $[\omega-1,\omega]$, so it is not an integer part.
\end{proof}
This is source~21's Theorem 8.4 and source~20's Theorem 1.1(i) with Corollary 6.5. The obstruction is
not a failure of the group law, of the finite-angle normalization, of the modulus or conjugation laws,
or of differentiability: all are built in.

\emph{What this answers.} Ehrlich and Kaplan ask, for an initial trigonometric-exponential subfield
$K\subseteq\No$ and any exponential on $K[i]$ extending their strip map and satisfying the modulus law,
under which conditions --- including initiality of $Z_E$ --- $Z_E$ is an integer part, equivalently
the exponential is the canonical one \cite[\S11.1, first open question; arXiv v3, printed
p.~34]{EK}. \Cref{trigonometry:per:thm:robustness} shows that initiality of $Z_E$ is \emph{not}
sufficient, even together with the strip values, the modulus and conjugation laws, the local Taylor
identities of every order, and agreement of the standard part of the phase at every argument; $Z_E$
can be as small as possible. With \cref{trigonometry:per:thm:criterion} this identifies the missing
algebraic content exactly: multiplicative closure of the normalized periods.

\emph{What it does not answer.} It is a partial negative answer to that one sufficiency possibility.
It does not classify the additional conditions that would suffice; it does not assert that Ehrlich and
Kaplan conjectured initiality to be sufficient; and it does not touch their separate second question,
on how different the exponentials induced by different initial embeddings of one
trigonometric-exponential field can be, whether they are isomorphic, and whether they have the same
first-order theory \cite[\S11.1, second open question]{EK}. The first-order separation of
\cref{trigonometry:per:sec:reflection} is proved in the broader class of all exponentials with the
strip and modulus properties, whose members need not arise from initial embeddings; it answers
nothing about that second question. The robustness question therefore remains open, re-scoped to
conditions beyond initiality.

\begin{remark}[The strip map in the source]\label{trigonometry:per:rem:ekstrip}
Source~21 observes that the strip map displayed on printed p.~33 of arXiv v3 of \cite{EK} is
$E(x+iy)=(\exp x)(\sin y+i\cos y)$, with sine and cosine interchanged, whereas the global formula
immediately after it uses $\cos b+i\sin b$; the former gives $E(0)=i$ and cannot literally be the
asserted group isomorphism. We checked this against the arXiv v3 PDF for this merge. Throughout,
``strip exponential'' means $\exp(x)(\cos y+i\sin y)$, as in \cref{trigonometry:thm:stripexp},
consistent with the group law and the adjacent formula; no argument uses the typographical discrepancy.
\end{remark}
$E'=E$ here is the fine derivative of a function of $z$, not the Berarducci--Mantova derivation of
the scalar field \cite{BM}. No claim is made that the noncanonical exponentials of
\cref{trigonometry:per:thm:robustness}(a),(b) are isomorphic or nonisomorphic to each other as
unrestricted exponential fields (source~21); their relation to the canonical one is settled by
\cref{trigonometry:per:cor:combined} below.

\subsection{Reflection to set-sized fields and first-order separation}
\label{trigonometry:per:sec:reflection}
This subsection is source~20's. Restricting a class exponential with $Z_E=\Z$ to a subfield need not
keep $Z_E=\Z$: a nonintegral element might fail to be a multiplier only because of a witness outside
the subfield. The witnesses of \cref{trigonometry:per:thm:relative} let one close the field under
them, and ordinary-circle values let the map itself restrict.

\begin{lemma}[Restriction to epsilon fragments]\label{trigonometry:per:lem:restriction}
Let $\Psi=\Psi_{\psi,0}$ be a phase on $\No$ with ordinary-circle character values. For every epsilon
number $\lambda$, $E_\chi$ maps $\No(\lambda)[i]$ onto $\No(\lambda)[i]^\times$ and is an exponential
there with the properties of \cref{trigonometry:per:prop:exponentials}.
\end{lemma}
\begin{proof}
For $x\in F=\No(\lambda)$, initiality and truncation closure put $\pr(x),\fin(x)$ in $F$;
$e^{2\pi i\psi(p)}$ is an ordinary complex number; $\cis$ of a finite element of $F$ lies in $F[i]$ by
trigonometric closure, and $\exp_{\No}$ preserves $F$ by exponential closure. Surjectivity follows as
in \cref{trigonometry:per:prop:exponentials} from the real logarithm of $F$ and surjectivity of
$\cis$ onto $\TT(F)$.
\end{proof}

\begin{theorem}[Unbounded set-sized reflection]\label{trigonometry:per:thm:reflection}
Let $E_{\mathrm s}$ and $E_{\mathrm n}$ be the split and nonsplit examples of
\cref{trigonometry:per:thm:relative,trigonometry:per:thm:nonsplit} (constructed in GBC), with witness
maps $x\mapsto\mu^{\mathrm s}_x,\mu^{\mathrm n}_x$. For every ordinal $\alpha$ there is an epsilon
number $\lambda>\alpha$ of cofinality $\omega$ such that on $F=\No(\lambda)$ both restrictions have
$Z_E=\Z$, the first restricted period sequence splits, the second does not, and their period groups are
not isomorphic.
\end{theorem}
\begin{proof}
Take an epsilon number $\lambda_0>\max\{\alpha,\omega\}$; given $\lambda_n$, the set of witnesses
$\mu^{\mathrm s}_x,\mu^{\mathrm n}_x$ for $x\in\No(\lambda_n)\setminus\Q$ (a set by replacement) has
birthdays bounded by an ordinal; choose an epsilon number $\lambda_{n+1}>\lambda_n$ above that bound,
and let $\lambda=\sup_n\lambda_n$, an epsilon number of cofinality $\omega$. For $x\in F\setminus\Q$
both witnesses and their products with $x$ lie in $F$, so $x$ is not a multiplier in $F$; nonintegral
rationals are excluded by $\cI\cap\OO_F=\Z$; the maps restrict by
\cref{trigonometry:per:lem:restriction}. The split section $p\mapsto p-A(p)$ stays in $F$. For the
nonsplit example, $\Q\omega\subseteq F$ and its displayed values persist, so an additive angle lift on
$\Pi_F$ would restrict to a forbidden lift on $\Q\omega$. \Cref{trigonometry:per:thm:homz} gives
nonisomorphism.
\end{proof}
\emph{Set-theoretic strength.} The reflection uses only arbitrarily large epsilon numbers and ordinal
bounds for sets of birthdays; no large cardinal is hidden. Global choice organizes the construction on
$\No$, and conservativity of GBC over ZFC for statements about sets means the resulting set-sized
existence statements need no additional strength. The theorem gives unboundedly many suitable
fragments, \emph{not} that every fragment works for a fixed constructed phase (closure under the
witnesses was part of the proof), and no computable description of the least suitable $\lambda$.

Consider the language of fields with one unary function symbol $E$, naming neither $i$, $\pi$, the real
axis, the order, standard part, nor the normal-form projections, and the sentence
\begin{equation}\label{trigonometry:per:eq:ringker}
 \mathsf{RingKer}:\quad \exists\tau\Bigl[\tau\ne0\wedge E(\tau)=1\wedge\forall x\,\forall y\bigl(
 (E(x)=1\wedge E(y)=1)\Rightarrow\exists w\,(xy=\tau w\wedge E(w)=1)\bigr)\Bigr].
\end{equation}

\begin{lemma}\label{trigonometry:per:lem:ringker}
If an exponential field satisfies $\mathsf{RingKer}$ with witness $\tau$, then $\tau^{-1}\ker E$ is a
unital subring contained in $Z_E$.
\end{lemma}
\begin{proof}
$\tau^{-1}\ker E$ is additive and contains $1$. For $x,y\in\ker E$, $xy/\tau\in\ker E$, so
$(x/\tau)(y/\tau)\in\tau^{-1}\ker E$ and $(x/\tau)y\in\ker E$.
\end{proof}

\begin{theorem}[First-order separation]\label{trigonometry:per:thm:firstorder}
Let $F$ be $\No$ or $\No(\lambda)$ with $\Pi_F\ne0$. The canonical exponential satisfies
$\mathsf{RingKer}$; every exponential as in \cref{trigonometry:per:prop:exponentials} with $Z_E=\Z$
fails it. Consequently, on each field of \cref{trigonometry:per:thm:reflection} the canonical
exponential field is not elementarily equivalent to either minimal example, even in the bare language of
fields with $E$.
\end{theorem}
\begin{proof}
For $\Exp$ take $\tau=2\pi i$: its kernel is $2\pi i(\Oz\cap F)$ and $\Oz\cap F$ is a ring. If $Z_E=\Z$
and $\tau$ witnesses $\mathsf{RingKer}$, \cref{trigonometry:per:lem:ringker} gives a unital ring
$\tau^{-1}\ker E\subseteq\Z$, hence equal to $\Z$, so $\ker E=\tau\Z$ is cyclic, contrary to
\cref{trigonometry:per:thm:dictionary}.
\end{proof}

\begin{corollary}[Three structures; source 20]\label{trigonometry:per:cor:three}
For arbitrarily large epsilon numbers $\lambda$ of countable cofinality, $\No(\lambda)[i]$ carries at
least three pairwise nonisomorphic exponentials agreeing on the real exponential and on every
finite-imaginary strip; $\mathsf{RingKer}$ separates the canonical one from the other two, and
\cref{trigonometry:per:thm:homz} separates those two.
\end{corollary}

\begin{corollary}[Combining the sources; stated in neither]\label{trigonometry:per:cor:combined}
In ZFC: (i) the exponential $E$ of \cref{trigonometry:per:thm:robustness}(a) on $\No(\eps_0)[i]$ fails
$\mathsf{RingKer}$, so $(\No(\eps_0)[i],E)$ is not elementarily equivalent, in particular not
isomorphic, to $(\No(\eps_0)[i],\Exp)$; (ii) $\No(\eps_0)[i]$ carries three pairwise nonisomorphic
exponentials agreeing on the real exponential and every finite-imaginary strip: $\Exp$, one with $Z_E=\Z$
and zero profinite defect, and one with $Z_E=\Z$ and nonzero profinite defect.
\end{corollary}
\begin{proof}
(i) $\No(\eps_0)$ is initial with $\Pi_F\ne0$, and $Z_E=\Z$, so \cref{trigonometry:per:thm:firstorder}
applies; its proof is elementary. (ii) By \cref{trigonometry:per:cor:simultaneous} on $\No(\eps_0)$ take
$d=0$ and some $d\ne0$ (possible since $\dim_{\Q}\Zhat/\Z=2^{\aleph_0}$); both have $Z_E=\Z$, so both
fail $\mathsf{RingKer}$; their period sequences split and do not split
(\cref{trigonometry:per:thm:split}), so \cref{trigonometry:per:thm:homz} distinguishes their kernels, and
an isomorphism of exponential fields restricts to an isomorphism of kernels.
\end{proof}
Part (i) settles source~21's non-claim only as far as comparison with the canonical exponential goes;
whether distinct noncanonical examples of source~21 are isomorphic to one another remains open.

\begin{corollary}[Definable integers and an omitted type; source 20]\label{trigonometry:per:cor:definable}
In each set-sized example with $Z_E=\Z$, the parameter-free formula
$\varphi_{\Z}(z):\ \forall y\,(E(y)=1\Rightarrow E(yz)=1)$ defines exactly the ordinary integers, with
their inherited ring operations, and the exponential field is not countably saturated: the type
$\{\varphi_{\Z}(x)\}\cup\{x\ne n\cdot1:n\in\Z\}$ is finitely satisfiable and omitted.
\end{corollary}
The failure of saturation concerns the exponential expansion, not the underlying field. By compactness
an elementary extension realizes the type, and there the definable multiplier ring has elements beyond
the numerals: $Z_E=\Z$ does not axiomatize standardness across the elementary class. $\mathsf{RingKer}$
gives non-elementary-equivalence; the invariant $\Hom(\cI,\Z)$ gives nonisomorphism of the split and
nonsplit examples but \emph{no} claim about their elementary equivalence, since existence of an
external homomorphism to $\Z$ is not a first-order sentence of the exponential field.

\subsection{Valued-field automorphisms act freely on phases}\label{trigonometry:per:sec:shifts}
In Hahn notation $t^g=\omega^{-g}$ with well-ordered supports, let $[g]_0$ be the real coefficient of
$\omega^0$ in the normal form of the exponent $g$; it is additive, also on infinite exponents, with
$[1]_0=1$, and it is not a standard part of an infinite number. On $F_{\Q}$ it is the identity. For real
$c$ put
\begin{equation}\label{trigonometry:per:eq:shift}
 \Sh_c(t^g)=t^g(1+ct)^{-[g]_0},\qquad \Sh_c\Bigl(\sum_g a_gt^g\Bigr)=\sum_g a_g\Sh_c(t^g),
\end{equation}
the powers of $1+ct$ being their binomial series.

\begin{lemma}[Support control]\label{trigonometry:per:lem:shiftsupport}
\Cref{trigonometry:per:eq:shift} is well defined on every Hahn series and every set-supported normal
form, and it preserves strongly summable families.
\end{lemma}
\begin{proof}
The new support lies in $S+\N$, which is well ordered; at a fixed exponent only finitely many
$(g,n)\in S\times\N$ contribute, since infinitely many would give a decreasing sequence in $S$. For a
strongly summable family apply this to the union of supports.
\end{proof}

\begin{theorem}[The shift family]\label{trigonometry:per:thm:shifts}
The maps $\Sh_c$ are strongly $\R$-linear field automorphisms of $F_{\Q}$ and of $\No$ that fix every
valuation and leading coefficient, with
\[
 \Sh_c\Sh_d=\Sh_{c+d},\qquad \Sh_c(t)=\frac{t}{1+ct},\qquad \Sh_c(\omega)=\omega+c .
\]
They preserve order, $\OO$, $\mm$, standard part and the finite phase:
$\Sh_c(\cis f)=\cis(\Sh_cf)$ for finite $f$, with $\Sh_c(i)=i$.
\end{theorem}
\begin{proof}
Additivity follows from \cref{trigonometry:per:lem:shiftsupport}; multiplicativity on monomials from
additivity of $[\,\cdot\,]_0$, and in general by strong products. Each monomial image is the monomial
times a unit of standard part one, so leading terms, order, $\OO$, $\mm$ and standard parts are kept.
Strong substitution gives $\Sh_c\Sh_d(t^g)=t^g(1+ct)^{-[g]_0}(1+dt/(1+ct))^{-[g]_0}=t^g(1+(c+d)t)^{-[g]_0}$,
identities among units $1+{}$infinitesimal, with no branch choice. The maps fix real coefficients and
preserve the strong sums defining $\Ex$, hence the finite phase.
\end{proof}
These are the fixed-shift flows of \cite[Theorem 5.2]{Saut} with $\delta=1$, $\lambda=[\,\cdot\,]_0$
(its $\ell$ of equation \texttt{saut:eq:ell}) and $\Sh_c=T_{-c}$ there; they are strong
$1$-automorphisms in the sense of \cite[\S4]{KKS}. Source~21 (Theorem 9.2) credits both and re-proves
them to keep the action theorem independent of any decomposition theory; source~20 (Theorem 11.1) proves
them without either citation. They are \emph{not} asserted to preserve $\exp_{\No}$, Conway's omega map
or surreal simplicity; ``value-preserving'' means $v(\Sh_cx)=v(x)$ in the fixed value-group copy, not
that the automorphism fixes the surreals labelling it.

The shifts act on phases by $(\Sh_c\cdot\Psi)(x)=\Sh_c\bigl(\Psi(\Sh_{-c}x)\bigr)$, which preserves the
finite restriction.

\begin{theorem}[Free action; source 21]\label{trigonometry:per:thm:free}
On $F_{\Q}$ or $\No$, $\Sh_c\cdot\Psi=\Psi$ for a phase $\Psi$ implies $c=0$. The real translations act
freely on phases and on their normalized period groups.
\end{theorem}
\begin{proof}
Invariance means $\Sh_c(\Psi(x))=\Psi(\Sh_cx)$. At $x=2\pi\omega/n$, since $\Sh_c$ fixes reals,
$\Sh_c x=x+2\pi c/n$, and taking standard parts gives $\st\Psi(x)=\st\Psi(x)\,e^{2\pi ic/n}$. The
standard part is nonzero, so $c/n\in\Z$ for every $n$, whence $c=0$. The acted phase has period group
$\Sh_c(\cI)$, and the kernel dictionary is injective.
\end{proof}

\begin{corollary}[Valued-field naturality fails]\label{trigonometry:per:cor:naturality}
No phase on $F_{\Q}$ or on $\No$ is invariant under all strongly $\R$-linear valued-field
automorphisms; no global phase can be singled out by an automorphism-invariant characterization using
only that structure and the real coefficients.
\end{corollary}

\begin{remark}[Source 20's form]\label{trigonometry:per:rem:halfshift}
Source~20 proves the case $c=1/2$ directly: equivariance under $\Sh_{1/2}$ would give
$\Sh_{1/2}\bigl(\Psi(2\pi\omega)\bigr)=\Psi(2\pi\omega+\pi)=-\Psi(2\pi\omega)$, and since $\Sh_{1/2}$
fixes standard parts of finite elements, $\st\Psi(2\pi\omega)=-\st\Psi(2\pi\omega)$, impossible for a
unit. The finite phase itself commutes with every $\Sh_c$;
the conflict arises only in extending it to infinite arguments. The obstruction concerns the
\emph{weaker} valued-field data. It does not contradict the canonical construction: initiality,
truncation and the monomial section are additional structure, and $\Sh_{1/2}$ moves the monomial
$\omega$ to $\omega+1/2$. Nor is a failure of invariance under exponential-field automorphisms
asserted; that is a strictly stronger preservation requirement which the present witnesses do not
satisfy in general.
\end{remark}

\emph{The orbit of the canonical phase.} Let $\Phi_c=\Sh_c\cdot\Phi$ on $F_{\Q}$ or $\No$ (on $F_{\Q}$,
$\Phi$ means $\cis\circ\fin$). Its period group $\Sh_c(\Pi_F\oplus\Z)$ is an integer-part ring with
zero profinite defect, the groups are distinct for distinct $c$, and the retraction is
$\Sh_c\circ\fin\circ\Sh_{-c}$.

\begin{proposition}[The standard-part shadow; source 21]\label{trigonometry:per:prop:shadow}
For a normal form $x=\sum_\alpha a_\alpha\omega^\alpha$ in $F_{\Q}$ or $\No$,
\[
 \st\Phi_c(2\pi x)=\exp\Bigl(2\pi i\sum_{n\in\N}a_n(-c)^n\Bigr),
\]
a finite sum.
\end{proposition}
\begin{proof}
The coefficient of $t^0$ in $\Sh_{-c}(x)$ comes only from terms $t^{-n}=\omega^n$ with $n\in\N$, since
the binomial corrections raise $t$-exponents by nonnegative integers; $\Sh_{-c}(t^{-n})=t^{-n}(1-ct)^n$
contributes $(-c)^n$. A reverse well-ordered support contains only finitely many nonnegative
integers. Applying $\Sh_c$ after the finite phase does not change the standard part.
\end{proof}
Translating the canonical phase changes its ordinary shadow while keeping a ring kernel; the avoidance
construction of \cref{trigonometry:per:thm:robustness}(a),(b) keeps the shadow while destroying every
nonintegral multiplier.

\subsection{What each test detects}\label{trigonometry:per:sec:tests}
\begin{center}\small
\begin{tabular}{L{0.26\textwidth}L{0.30\textwidth}L{0.34\textwidth}}
\toprule
Test (source 21) & Detects & Does not force\\
\midrule
Finite values and local Taylor identities & The canonical germ on every infinitesimal neighbourhood & The global phase or its period ring\\
Standard part of the phase at every input & The ordinary character $\psi$ & The infinitesimal character $\cL$ or the multiplier ring\\
Rational-division residues & The profinite defect and $\cI^{\mathrm{div}}$ & Multiplicative closure of periods\\
Closure of normalized periods under products & $\Mult(\cI)=\cI$ & Initiality of this integer part\\
Initiality of the multiplier ring alone & A simplicity-closed discrete subring & The integer-part property\\
Invariance under all valued translations & Impossible for every phase & Nothing against extra simplicity structure\\
\bottomrule
\end{tabular}
\end{center}
\begin{center}\small
\begin{tabular}{L{0.30\textwidth}L{0.17\textwidth}L{0.17\textwidth}L{0.19\textwidth}}
\toprule
Invariant (source 20; $F$ reflected) & Canonical & Minimal split & Minimal nonsplit\\
\midrule
$Z_E$ & $\Oz\cap F$ & $\Z$ & $\Z$\\
Normalized kernel a ring & yes & no & no\\
Period sequence splits & yes & yes & no\\
$\Hom(\cI,\Z)$ & $\Z$ & $\Z$ & $0$\\
$\mathsf{RingKer}$ & true & false & false\\
Finite-angle agreement; $E'=E$ & yes & yes & yes\\
\bottomrule
\end{tabular}
\end{center}
The second table does not equate the first-order theories of the two minimal examples; that is not
proved. Three kinds of splitting must not be conflated (source 20): the domain always splits as
$F=\Pi_F\oplus\OO_F$; the circle always splits as $\TT(F)\cong\TT(\R)\times\mm_F$; neither makes the
period sequence split, and a profinite character sits exactly between them. Conversely a split kernel
need not be a ring: the minimal split example is additively isomorphic to the canonical kernel, yet its
multiplier ring is $\Z$; ring closure depends on the embedding, not on the abstract group. Two
failures are logically independent (source 21): \cref{trigonometry:per:ex:profinite} has nonzero
defect and no ring closure; \cref{trigonometry:per:thm:robustness}(a),(b) has zero defect and no ring
closure; the canonical orbit has zero defect and ring kernels. ``All periods split into a divisible
part plus integers'' does not imply ``the periods form a ring.''

On $F_{\Q}$ a reduced period group embeds in $\Zhat$ by its residue map, so distinct periods are
separated by congruences: the classical ``Leibnizian'' phenomenon for $\Z$-groups
\cite[Theorem 3.3]{Jerabek}, not a new theory of Presburger arithmetic. Nothing is claimed about
tameness, decidability, o-minimality or elementary equivalence of the phase expansion with the ordinary
complex exponential field. The avoidance constructions are existence arguments using linear choices;
they give no finite names, computable coefficient oracles, algorithms for the phases, or effective
recognition of periods, and finitely many rational divisions cannot determine an element of $\Zhat$.

\subsection{Checks, formal status, and remaining questions}\label{trigonometry:per:sec:questions}
Source~20's script (\texttt{code/20-exponential-kernels-verify.py}) passed $47{,}939$ exact finite
checks: compatible CRT residues and additivity of $f_{\alpha_3}$, finite portions of the $2$- and
$3$-power root towers, coefficient identities for $\Sh_a\Sh_b=\Sh_{a+b}$ through a fixed order, and
finite Laurent-polynomial instances of \cref{trigonometry:per:lem:fresh} including a pair of periods
whose product is not a period. Source~21's script (\texttt{code/21-period-arithmetic-verify\_examples.py})
passed checks on $128$ CRT moduli, $645$ residue-compatibility pairs, $58{,}081$ additivity pairs of
$f_{\alpha_2}$, $15$ displayed example assertions, $1{,}575$ translation-composition coefficients,
$24$ one-step multiplier witnesses and $20$ shadow checks. All arithmetic is exact. \emph{These checks
do not establish the transfinite recursions, arbitrary surreal supports, the reflection theorem, the
abstract vector-space sections or novelty.}

Neither source was verified in Lean, and neither built or modified the repository. Both propose a
formalization order --- the kernel dictionary over an ordered field with a surjective phase of kernel
$\Z$; $\Mult$ and the criterion; the splitting, the CRT characters and $\Zhat/\Z$ as abstract groups;
the one-step or fresh-support lemma with an explicit set-sized parameter; then the set-sized recursion,
with the class statements in a universe-relative form --- and identify the delicate obligations: support
control for \cref{trigonometry:per:eq:shift}, compatibility of residues, and the cardinal bookkeeping of
the recursion. These are proposed architectures, not existing declarations.

\emph{Remaining questions}, as the sources state them: which discrete subrings occur as $Z_E$ on a fixed
initial field, and which rings between $\Z$ and an additive integer part occur for a fixed ordinary
character and profinite defect (sources 20, 21; source 21 realizes $\Z$ for every defect and exhibits
maximal rings on the canonical orbit, and classifies no intermediate ring); whether an independent
compatibility with simplicity, truncation or multiplication forces the normalized periods to be a ring
(source 21); whether the split and nonsplit minimal examples can have the same first-order theory
(source 20); which definability requirements in a specified expansion exclude the diagonal
constructions without naming the canonical projection (source 20); the orbits of phases under the full
valued automorphism group, which the free action does not classify (source 21); and the isomorphism
types of the resulting exponential fields when no base structure is fixed (source 21; the comparison
with the canonical exponential is \cref{trigonometry:per:cor:combined}). The second Ehrlich--Kaplan
question is untouched.

\section{Conclusions and dependency audit}\label{trigonometry:sec:conclusion}

\subsection{What the theory establishes}
The central geometric conclusion is exact: \emph{arbitrary surreal lengths do not require infinite
circular angles.} The finite-angle quotient $\OO_{\R}/2\pi\Z\cong\TT(\No)\cong\TT(\R)\times\mm_{\R}$
describes every unit direction and remembers both an ordinary direction and its full infinitesimal
displacement, and ordinary finite triangle constructions give the same identities at all scales, in
the plane, on the sphere and in the hyperbolic disk. The infinitesimal part is not negligible data:
it controls chord precision, nearly straight angles, the order of the triangle-inequality slack,
the relative flatness threshold, and tangency splitting. Valuations then extract information with
no ordinary numerical counterpart --- angle hierarchies, flattening thresholds, and condition loss
near collisions.

A second conclusion concerns the complex plane. The canonical strip $\mathscr V$ is large enough
for a surjective complex sine with completely determined fibres and for the complete algebraic
solution theory of every finite trigonometric polynomial; the passage $U=E(iz)$ turns periodic root
questions into finite-degree polynomial questions. Positive-support coefficient lifting
distinguishes stable finite algebras from unstable individual roots, and the three stability
certificates of \cref{trigonometry:sec:threecertificates} cover three genuinely different
questions. Positivity, too, has to be tested on the full surreal circle and not only on its
ordinary real shadow.

Finally, full real-argument oscillation is a different problem. The Ehrlich--Kaplan integer-part
normalization supplies one explicit answer; the character classification shows it is one member of
a family all of whose members are fine-locally analytic, obey the same addition laws and the same
differential equations, necessarily acquire infinite periods, and fail common-domain coherence at
infinite frequency. \emph{Thus neither nonexistence nor canonical uniqueness is an appropriate
blanket claim about ``surreal sine.''} The answer depends on the domain and on the analytic
structure specified.

Read through its rotation group, the circle is a class group with the classical finite algebra and
an exact, very large infinitesimal factor (\cref{trigonometry:rot:sec:main}). Isometries fixing the
origin are linear; torsion and finite subgroups are the standard ones; each valuation layer of
infinitesimal rotations is one real line; algebraic endomorphisms and weights are ordinary integers
while fine-continuous nonalgebraic symmetries exist; and the fine, standard-part and semialgebraic
geometries of the same circle answer different questions and must not be interchanged.

The periods of the global extensions have an arithmetic of their own (\cref{trigonometry:per:sec:main}).
Normalized period groups are exactly the additive integer parts; their rational-division residues form a
profinite defect, which can be prescribed arbitrarily and vanishes exactly when a linear finite-part
choice suffices; on $\No$ divisible periods are cofinal, on $\R((t^{\Q}))$ they can be absent. Neither
abundance of divisible periods nor perfect local exponential behaviour forces multiplicative closure:
the kernel-multiplier ring can be the initial ring $\Z$, so initiality alone does not select the
Ehrlich--Kaplan exponential, while closure of the normalized periods under products, together with
initiality, does. The
canonical choice therefore depends on global structure beyond the local identities and beyond
valued-field symmetry, which acts freely on all phases.

\subsection{A dependency audit}
\begin{center}
\small
\begin{longtable}{L{0.28\textwidth}L{0.62\textwidth}}
\toprule
Result family & Inputs actually used\\
\midrule\endhead
Finite sine, cosine, circle logarithm & Ordinary analytic germs, normal forms, positive-support Neumann lemma\\
Triangle, circle and concurrence geometry & Ordered-field algebra, positive square roots, finite-angle circle parametrization\\
Valuation estimates, defect and flatness & Exact algebraic identities and unit-leading infinitesimal series\\
Tangency, cosine fold, coupled collision & Strong composition of square-root and inverse-sine germs at infinitesimal arguments\\
Conditioned root existence and stability & A finite rational tangent chart, polynomial intermediate values, a nonsingular residue root; no Newton-sequence convergence\\
Cluster lifting & Coprime polynomial factorization, a finite Sylvester map, transfinite coefficient recursion over $S^*$\\
Strip functions, hyperbolic disk, spherical laws & Gonshor exponential and logarithm, Gram determinants, plus the finite-angle theory\\
Trigonometric root bounds & Algebraic closedness and the strip exponential quotient\\
Fej\'er--Riesz factorization & Real closedness, rational circle parametrization, finite polynomial factorization\\
Finite Fourier identities & Finite geometric sums; the coefficientwise integral of the ring (R1) where an integral is used\\
Arc and sector formulas & The explicit primitive-based coefficientwise integral convention, ring (R1)\\
Infinite phase classification & The additive normal-form split $\No=\Pi\oplus\OO_{\R}$ and circle surjectivity\\
Coherence obstructions & Ring (R1) only; normal-form uniqueness and one well-ordered support\\
Rotation group: isometries, torsion, dihedral subgroups, invariants & Polarization, the splitting of \cref{trigonometry:thm:polar}, ordinary roots of unity, one-variable Laurent independence\\
Norm-one torus, Lie algebra, representations & Finite algebra over a real closed field, dual numbers, Laurent coactions; the torus language of \cite{Milne}\\
Fine, standard-part and semialgebraic circle & The lower-bound cut, open infinitesimal subgroups, the chord comparison; $\pi_1^{\mathrm{sa}}\cong\Z$ imported from \cite{DK,BaroOtero}\\
Period arithmetic: dictionary, defect, splitting & \Cref{trigonometry:thm:characters}, floors, profinite residues; classical $\Z$-group facts \cite{Jerabek}; a set-sized linear section\\
Minimal multiplier rings, reflection & Two-vector avoidance or fresh ordinal monomials; ZFC recursion on a continuum-sized field, or GBC recursion on $\No$ with conservativity; initial integer parts \cite[Lemma 11.1]{EK}\\
Free action of translations & The fixed-shift flows of \cite[Theorem 5.2]{Saut}, support control, standard parts\\
\bottomrule
\end{longtable}
\end{center}
No result uses a proper-class sum, assumes that a positive valuation increment is cofinal, or
identifies fine differentiation with a chosen derivation of the surreal field \cite{BM}. The
several-variable Noetherianity, division and residue theorems of
\cite{SourceGeometry,SourceIntersections} are not assumed as black boxes anywhere.

\subsection{Reproducible checks and their limits}
Sources 16--18, the three editions of the trigonometric core, carried accompanying Python scripts using exact SymPy
algebra and formal power series, which between them check the Gram, rational-circle, Heron,
circumradius, Euler and Ptolemy identities; the Cayley group law; Chebyshev formulas; the spread
and triple-spread identities; flat-triangle, tangency and inverse-cosine coefficients; the coupled
angular algebra after sine coordinates; the algebraization of trigonometric equations; line--circle
formulas and the collision algebra; finite quadrature and Fourier convolution; a spectral-factor
example; and the rational identities behind the spherical and hyperbolic laws. Those runs reported
$67$, $43$ and $68$ passing checks respectively. Source 19's script checks, in $41$ exact identities,
the Cayley norm, inverse chart, affine and homogeneous group laws and associativity, the exceptional
chart cases and \cref{trigonometry:rot:eq:halfturn}, the invariant differential
\cref{trigonometry:rot:eq:cayleyform}, rotation-matrix products, orthogonality and determinant, the
displacement norm identity and the rational corrections of \cref{trigonometry:rot:ex:cayleydisp}, the
Clifford calculation \cref{trigonometry:rot:eq:spin}, and truncated Taylor coefficients of the chart,
of \cref{trigonometry:rot:eq:arctanbridge}, of $\Ex(ih)\Ex(-ih)=1$ and of the chord--angle ratio.

\emph{These checks validate the finite algebra and selected Taylor coefficients occurring in the
proofs.} They do not implement the full surreal class, verify transfinite support arguments or
recursions, establish real closedness, or machine-check any mathematical proof. The support lemma,
the branch arguments and the proper-class distinctions are justified in the text, not inferred from
symbolic experiments, and no proof-assistant verification is claimed. The same holds for source
19's checks: they cannot check the surreal normal-form theorem, a class-sized topology, Neumann
summability for arbitrary supports, or the universal quantifiers of the structure theorem, and no Lean
compilation or independent kernel verification of that manuscript was performed.

Sources 20 and 21 carried standard-library Python scripts with exact rational arithmetic; their scope
and counts ($47{,}939$ checks, and $128$ moduli, $645$ compatibility pairs, $58{,}081$ additivity pairs,
$1{,}575$ translation coefficients, $24$ witnesses and $20$ shadow checks, respectively) are stated in
\cref{trigonometry:per:sec:questions}. They test CRT residues, root towers, shift coefficients and
finite instances of the avoidance step, not the transfinite recursions, the reflection theorem or
arbitrary supports; no Lean proof of either manuscript was built.

\subsection{What is not proved}\label{trigonometry:sec:notproved}
The following non-claims are made in the body of the article and are collected here. They are
specific to individual results and are not pooled into a single blanket disclaimer.

\begin{enumerate}\itemsep0.35em
\item \emph{(\cref{trigonometry:prop:lift}.)} No general transfer of arbitrary infinite families of
 inequalities is asserted; in particular, coefficients that are themselves infinitesimal are not
 automatically covered by an ordinary real-coefficient sign argument. The uniqueness assertion is
 deliberately about the Taylor rule, not about arbitrary functions satisfying the same differential
 equation in the fine topology.
\item \emph{(\cref{trigonometry:prop:order}.)} The proof does not infer global monotonicity from
 the sign of a fine derivative; such an inference would require additional hypotheses on a general
 surreal function.
\item \emph{(\cref{trigonometry:prop:leading}, \cref{trigonometry:cor:phaseisometry}.)} The
 asymptotic equivalences are statements about a single arbitrary infinitesimal. They are stronger
 than standard-part equalities and do not mean that a sequence of real arguments converges in the
 fine topology.
\item \emph{(\cref{trigonometry:rem:noendpointderivative}.)} The endpoint derivative of inverse
 sine does not exist as a surreal-valued derivative: its difference quotient grows beyond every
 fixed surreal bound as the increment becomes sufficiently small, and allowing infinitely large
 elements as derivative values does not repair this, because there is always a still smaller
 increment. Throughout, these statements are about differentiation of functions of a variable, not
 about a derivation on surreal numbers interpreted as transseries.
\item \emph{(\cref{trigonometry:thm:sss}.)} The usual side--side--angle ambiguity remains: a
 ray--circle intersection is a quadratic problem and may have zero, one or two admissible
 solutions; no additional infinite family of solutions appears merely because the coordinates are
 surreal.
\item \emph{(\cref{trigonometry:eq:spreadlaws}.)} Quadrances and spreads forget orientation and
 distinguish neither an angle from its supplement nor a directed angle from its negative. That loss
 is useful for rational formulas but must be remembered.
\item \emph{(\cref{trigonometry:thm:ceva}.)} A directed Menelaus or Ceva formula can be obtained in
 the same way after fixing signs, but is not needed for the present development.
\item \emph{(\cref{trigonometry:thm:ptolemy}.)} Cyclic order is the order induced by finite
 representatives around the circle; it is not defined by fine-continuous motion of a real
 parameter. More generally, algebraic identities among finitely many vectors often survive
 unchanged, while any analytic interpretation of the associated angles must still be supplied.
\item \emph{(\cref{trigonometry:thm:metric}.)} Although cosine loses first-order information near
 zero, the full unit direction does not. At other angles one should use both coordinates, or a
 nonsingular tangent chart, rather than infer angular error solely from cosine error.
\item \emph{(\cref{trigonometry:thm:valuationtriangle}.)} The conclusion about side valuations does
 not say that all three sides have equal valuation: a very short third side is possible.
\item \emph{(\cref{trigonometry:thm:flatrelative}.)} The denominator $bc/(b+c)$ is essential. A
 defect that is small compared with the perimeter need not make the opposite angle close to $\pi$
 when one adjacent side is much shorter than the other.
\item \emph{(\cref{trigonometry:thm:tangency}, \cref{trigonometry:thm:conditioned}.)} The condition
 $e\prec\tau$ cannot be omitted from a uniform branch-matching theorem: $e=-\tau$ merges the roots
 and $e<-\tau$ destroys the two real intersections. Correspondingly the strict threshold
 $\sigma>2\kappa$ cannot be weakened to equality uniformly near an extremum.
\item \emph{(\cref{trigonometry:sec:coupled}.)} The dimension four of $\mathcal B_{s,u}$ is a
 \emph{total} multiplicity of the entire finite zero set in the monad, not a claim that each
 individual simple root has a four-dimensional local algebra.
\item \emph{(\cref{trigonometry:thm:cluster}.)} The support bound belongs to the factor
 coefficients, not to individual roots, which may require fractional scales; and no reality
 conclusion is asserted for an arbitrary multiple zero. The construction is a coefficient
 recursion, not a convergence claim for a Newton sequence.
\item \emph{(\cref{trigonometry:cor:chebyshev}.)} The derivative statement bounds stationary
 angles; it does not invoke a general compactness or extreme-value theorem on a surreal interval.
\item \emph{(\cref{trigonometry:thm:parseval}, \cref{trigonometry:eq:fejerkernel}.)} The finite
 Fourier identities establish no convergence or completeness for an unrestricted infinite Fourier
 series; infinitely many ordinary Fourier summands at one Hahn exponent do not form a strongly
 summable family merely because the corresponding real series converges. The Dirichlet and
 Fej\'er kernel identities are exact algebraic statements, not a fine-topological approximation
 theorem.
\item \emph{(\cref{trigonometry:thm:fejer}.)} The theorem is a real-closed-field extension of a
 classical result, not a claim of a new spectral-factorization theorem. Multivariable
 trigonometric positivity is a different problem; no multivariable single-square theorem is
 asserted.
\item \emph{(\cref{trigonometry:ex:invisible}.)} The example does not contradict
 \cref{trigonometry:prop:lift}: its coefficients depend on a non-real infinitesimal parameter, so
 the polynomial is not a single ordinary real-coefficient analytic function.
\item \emph{(\cref{trigonometry:sec:arcs}.)} The coefficientwise arc-length convention is not an
 assertion that a nonconstant ordinary real-parameter path is continuous in the full fine topology,
 nor a definition of an unrestricted Riemann integral; other surreal integration frameworks have
 substantially different scope \cite{CE}.
\item \emph{(\cref{trigonometry:thm:spherical}, \cref{trigonometry:thm:hyperbolic}.)} The spherical
 and hyperbolic sections give metric identities under explicit nondegeneracy hypotheses, not a
 complete global geometric or measure-theoretic foundation, and not a theory of hyperbolic area or
 of non-Archimedean geodesic completeness.
\item \emph{(\cref{trigonometry:def:globaltrig}, \cref{trigonometry:thm:globalexp}.)} The global
 normalization agrees with the established Ehrlich--Kaplan construction and is not claimed as a new
 definition; the published characterization of $\Oz$ is used to explain the word canonical, and its
 tree-theoretic proof is not reproduced. The local power-series argument establishes fine-local
 analyticity without asserting uniqueness from a differential equation.
\item \emph{(\cref{trigonometry:thm:characters}.)} The classification concerns group homomorphisms
 with the specified finite restriction; it is not a classification of all solutions of a
 differential equation, and by itself it does not answer the separate robustness questions
 formulated in \cite[\S11]{EK} for more general initial fields and their integer parts
 (\cref{trigonometry:per:sec:ek} records what sources 20 and 21 add, and what stays open). It does not
 follow that
 defining a full-plane convention is impossible: the trivial-character convention is perfectly
 explicit. What fails is uniqueness from the stated analytic and algebraic axioms.
\item \emph{(\cref{trigonometry:thm:allscale}, \cref{trigonometry:thm:infinitefrequency}.)} Both
 are restrictions on uniform analytic descriptions over the ring (R1) of
 \cref{trigonometry:sec:rings}, not failures of the algebraic identities or of arbitrary-radius
 circle geometry, and neither is asserted over the radius-free ring (R3) or the formal ring (R4).
 It is not claimed that a unique canonical value of circular sine at every infinite real argument
 follows from the finite-angle theory.
\item \emph{(\cref{trigonometry:rot:sec:main}, source 19.)} The following non-claims are source 19's,
 each also stated at its point of use.
 \begin{enumerate}[label=(\alph*),itemsep=0.2em]
 \item Source 19 is a self-contained exposition organized around the rotation group, not a claim of
  priority or a machine-checked formalization; the finite-angle splitting is not presented as a new
  discovery, and it is compatible with, but logically prior to, the global construction of
  \cite[\S11]{EK}.
 \item ``Rotation'' means the algebraic isometry; continuity is an extra property in a specified
  topology. Bounded matrix entries do not make the group set-sized or ordinarily compact.
 \item All statements about $\SO(2,\No)$ are class-group statements: no proper-class Hahn sum, no
  quantification over a set of all class functions, no treatment of the group as a set-sized
  manifold, and topological statements use ball neighbourhoods with explicit quantifiers.
 \item $\mm_{\R}$ is not an $\No$-vector space, and the $\R$-vector structure of $\Tm$ does not extend
  to surreal scalars (\cref{trigonometry:rot:rem:notvector}).
 \item Preserving norms alone does not force linearity, and a field automorphism changing distances
  is not an isometry; the semidirect products \cref{trigonometry:rot:eq:orthogonal} are algebraic, not
  statements about connected components.
 \item The rational chart parametrizes $F$-points; it is not an isomorphism of $\PP^1_F$ with the
  affine torus (\cref{trigonometry:rot:rem:projective}).
 \item Formal evaluation at infinitesimals is asserted for real or complex coefficients only; an
  arbitrary series with surreal coefficients need not be evaluable at every infinitesimal. Strong sums
  are not limits of partial sums in the fine topology.
 \item A branch of $\Arg$ is not an additive homomorphism; the group law belongs to angle classes.
  $\No/2\pi\Z$ is not the finite-angle model compatible with the finite phase; source 19 does not
  assert that some unrelated abstract class-group isomorphism $\No/2\pi\Z\cong\TT(\No)$ is impossible,
  the compatibility requirement being essential.
 \item Torsion means ordinary finite order; infinitesimal rotations give elements of infinite order,
  not new regular polygons. The square-root formula \cref{trigonometry:rot:eq:sqrt} is a local branch,
  not a multiplicative section; the $\operatorname{Spin}(2)$ calculation is finite algebra, not path
  lifting; finite orthogonal groups are conjugate to standard ones, which does not make their matrices
  real.
 \item \Cref{trigonometry:rot:thm:graded} is a filtration, not a direct-sum decomposition; $\sim$ is an
  algebraic relation, not a sequence limit; \cref{trigonometry:rot:eq:hahnproduct} gives no value to
  infinitely many identical turns, and its partial products do not converge finely.
 \item Cyclic order and the absence of fixed points are algebraic, not defined by ordinary real paths;
  \cref{trigonometry:rot:eq:invariants} concerns polynomial functions, not all class functions on
  orbits.
 \item The algebraic group is applied at surreal points, without treating $\No$ as a set-sized
  coefficient field of scheme theory; the exponential series is not summable at infinite $T$ and the
  canonical exponential lives on $\OO_{\R}\mathsf J$; the invariant form defines no integration over
  arbitrary full-class contours; a connected-component argument cannot replace the tangent
  calculation.
 \item The representation classification supplies no Fourier convergence theorem, no Haar measure on
  the full class, and no classification of arbitrary abstract representations or of all abstract
  automorphisms; the automorphisms $\Xi_a$ use no basis choice.
 \item The surreal exponent \cref{trigonometry:rot:eq:surrealpower} is an analytic scaling on a
  restricted domain, not a transfinite iteration; beyond it a phase extension supplies further choices.
 \item Small-set discreteness is special to the full class, not a theorem about set-sized
  non-Archimedean fields with their intrinsic topology; it does not forbid a properly class-indexed net
  approaching zero. There is no ordinary real Lie-group interpretation of the fine circle, and
  classical compactness or Haar integration cannot be imported.
 \item The fine and standard-part uses of $\st$ must not be conflated. The semialgebraic fundamental
  group \cref{trigonometry:rot:eq:sapi} is imported from \cite{DK,BaroOtero}, not proved; the sequence
  $2\pi\Z\to\OO_{\R}\to\TT(\No)$ is not a universal-covering claim.
 \item Solving the local rotation equation with $\Psi(0)=1$ does not choose a global normalization;
  the classification of phase homomorphisms is not a classification of all solutions of the fine
  differential equation; mean-value and initial-value uniqueness arguments need further hypotheses
  (\cref{trigonometry:rot:ex:locallyconstant}); no compatibility with a scalar derivation is assumed.
 \item No representation of arbitrary surreal class elements by finite strings is asserted, and
  equality of effective Hahn-series names is not claimed decidable.
 \item Source 19 claims no novel construction of global surreal sine, no resolution of an
  Ehrlich--Kaplan robustness question (sources 20 and 21 later answer one part of the first one
  negatively, \cref{trigonometry:per:sec:ek}), no classification of all abstract automorphisms, no Haar or
  Fourier theory on the full class group, and no global mean-value or initial-value theorem for fine
  differentiable functions; it does not assume that an infinitesimal Taylor sum is a fine limit. Its
  repository search covered the root and documentation overviews and the opening of this article, and
  no exhaustive search of the repository or of the literature is claimed; its arguments do not
  presuppose verification of the rest of this report, whose component manuscripts are unrefereed
  AI-assisted drafts.
 \end{enumerate}
\item \emph{(\cref{trigonometry:per:sec:main}, sources 20 and 21.)} The following non-claims are those
 of sources 20 and 21, each also stated at its point of use; the source is named when only one makes it.
 \begin{enumerate}[label=(\alph*),itemsep=0.2em]
 \item \Cref{trigonometry:per:thm:robustness} is a partial negative answer to one sufficiency
  possibility in the first question of \cite[\S11.1]{EK}; it does not classify additional axioms that
  would suffice, does not assert that Ehrlich and Kaplan conjectured initiality to be sufficient, and
  does not answer their second question on exponentials from different initial embeddings (both). The
  first-order separation is proved in the broader class of all exponentials with the strip and modulus
  properties (source 20).
 \item Finite-angle uniformization, the canonical phase, the character classification
  \cref{trigonometry:thm:characters}, unavoidable infinite periods \cref{trigonometry:thm:infiniteperiods}
  and the local law are prerequisites from this article, not new; the canonical integer-part
  exponential is published \cite[Proposition 11.2, Theorem 11.1]{EK}; not every elementary lemma is
  claimed new (both).
 \item The residue map, the maximal divisible subgroup and the splitting and homomorphism facts are
  classical $\Z$-group theory \cite{Jerabek}; the Leibnizian phenomenon is not a new theory of
  Presburger arithmetic (source 21; source 20 does not cite this literature).
 \item The shift automorphisms are a prior type of construction \cite[Theorem 5.2]{Saut},
  \cite[\S4]{KKS}, re-proved; they are not asserted to preserve $\exp_{\No}$, Conway's omega map or
  simplicity; ``value-preserving'' refers to the fixed value-group copy; the naturality obstruction
  concerns only the weaker valued-field data, is compatible with the Ehrlich--Kaplan construction, and
  no failure of invariance under exponential-field automorphisms is asserted (both).
 \item The integer-part criterion is an exact reduction, not a classification of every multiplier ring
  and not a substitute for independent hypotheses implying multiplicative closure; intermediate
  multiplier rings are not classified; the free action does not classify orbits under the full valued
  automorphism group (both).
 \item No claim is made that distinct noncanonical exponentials of source 21 are isomorphic or
  nonisomorphic to one another (source 21); their non-isomorphism with the canonical one is the merge's
  \cref{trigonometry:per:cor:combined}. The split and nonsplit minimal examples are nonisomorphic, with
  no claim that their first-order theories differ (source 20).
 \item Reflection gives unboundedly many suitable fragments $\No(\lambda)$, not every fragment for a
  fixed phase, and no computable least $\lambda$; no large-cardinal strength is hidden, the set-sized
  consequences resting on conservativity of global choice (source 20).
 \item The failure of countable saturation concerns the exponential expansion; $Z_E=\Z$ does not
  axiomatize standardness across the elementary class (source 20). Nothing is claimed about tameness,
  decidability, o-minimality or elementary equivalence of the phase expansion with the complex
  exponential field; no finite names, computable oracles or effective recognition of periods
  (source 21).
 \item $E'=E$ is the fine derivative, not the Berarducci--Mantova derivation; the local law is
  inherited, not a new uniqueness theorem; strong sums, including those in the cofinal-period theorem,
  are not asserted to be fine limits of partial sums (both).
 \item $\No(\eps_0)$ is not asserted to be a full Hahn field; $\Zhat/\Z$ is an abstract group, not a
  topological quotient; the counts on $\R((t^{\Q}))$ are counts of embedded period groups, not of
  abstract isomorphism types, and use no continuum hypothesis; the explicit profinite examples
  establish neither a minimal multiplier ring nor reducedness (source 21).
 \item Both are AI-assisted, unrefereed drafts, not Lean-verified; neither modified or built the
  repository; their proposed Lean routes are architectures, not existing declarations; their finite
  checks ($47{,}939$ and the counts of \cref{trigonometry:per:sec:questions}) test examples and finite
  identities only, not transfinite recursions, arbitrary supports, reflection or novelty. Novelty is
  provisional and rests on targeted audits (source 20: the root and documentation READMEs, this
  report's guide and its Section 17, and targeted searches; source 21: the repository tree, root and
  documentation READMEs, this report and the automorphism report's guide); failure to find a prior
  statement is not proof of priority.
 \end{enumerate}
\item \emph{(Global.)} No proper-class sum, no arbitrary infinite polygon, no full-class Riemann
 integral, no ordinary sequential Fourier convergence theorem, and no uniqueness theorem for all
 fine-analytic differential equations is asserted anywhere. The three supplied manuscripts are
 identified as unpublished supplied documents, not as published literature. The literature search
 behind this report was targeted rather than exhaustive: a proof of an analogue is not, by itself,
 evidence that its exact formulation is absent from all previous literature, and no priority claim
 and no claim to have resolved a named published conjecture is made; the partial answer of
 \cref{trigonometry:per:sec:ek} refutes one sufficiency possibility and leaves the Ehrlich--Kaplan
 questions open.
\end{enumerate}

\appendix
\section{A compact formula sheet with domain reminders}\label{trigonometry:app:formulas}
For finite real $\theta$, nonzero $z,w\in\SC$, $u,v\in\TT(\No)$ and an ordinary integer $n$:
\begin{align*}
 z&=|z|\cis\Arg z&&\text{with }\Arg z\in\OO_{\R}/2\pi\Z,\\
 \cos\ang zw&=\frac{\re(\bar zw)}{|z||w|},&
 \sin\ang zw&=\frac{\im(\bar zw)}{|z||w|},\\
 \cis(n\theta)&=(\cis\theta)^n,&
 \tan(\theta/2)&=\frac{\sin\theta}{1+\cos\theta},\\
 |u-v|&=2\sin\bigl(\ang uv/2\bigr),&
 \arctan t&=\Arg(1+it)\in(-\pi/2,\pi/2)\quad(t\in\No).
\end{align*}
The first argument is an angle class, and the sine in the directed-angle formula may be negative.
In a triangle of positive area the interior sines are positive and
\begin{gather*}
 \alpha+\beta+\gamma=\pi,\qquad a^2=b^2+c^2-2bc\cos\alpha,\qquad a=2R\sin\alpha,\\
 \Delta^2=s(s-a)(s-b)(s-c),\qquad r=\Delta/s,\qquad R=abc/(4\Delta),\\
 \tan(\alpha/2)=r/(s-a),\qquad \ell_A=\tfrac{2bc}{b+c}\cos(\alpha/2),\qquad
 |O-I|^2=R(R-2r),\\
 v(\Delta)=v(b)+v(c)+v(d(\alpha)),\qquad v(R)=v(a)+v(b)+v(c)-v(\Delta).
\end{gather*}
For $z=x+iy$ with finite $x$ and arbitrary $y\in\No$,
\[
 \Sin z=\sin x\cosh y+i\cos x\sinh y,\qquad E(iz)=\exp_{\No}(-y)\cis x .
\]
For infinite real $x$, circular sine and cosine require an explicitly selected extension; the
lowercase finite-angle notation alone does not specify their values, and every such extension has
an infinite period --- on $\No$, cofinally many rationally divisible ones
(\cref{trigonometry:per:thm:cofinal}).

\section{Worked formal series}\label{trigonometry:app:series}
All of the following are strongly summable for infinitesimal $h$ or $\tau$, with the indicated
positive square root in the last line:
\begin{align*}
 \arctan h&=h-\frac{h^3}3+\frac{h^5}5-\frac{h^7}7+O(h^9),\\
 \arcsin h&=h+\frac{h^3}6+\frac{3h^5}{40}+\frac{5h^7}{112}+O(h^9),\\
 \Logm(1+h)&=h-\frac{h^2}2+\frac{h^3}3-\frac{h^4}4+O(h^5),\\
 \sqrt{1+h}&=1+\frac h2-\frac{h^2}8+\frac{h^3}{16}-\frac{5h^4}{128}+O(h^5),\\
 \arccos(1-\tau)&=\sqrt{2\tau}\left(1+\frac\tau{12}+\frac{3\tau^2}{160}+\frac{5\tau^3}{896}
 +\frac{35\tau^4}{18432}+O(\tau^5)\right).
\end{align*}
For example, if $h=t^\lambda+t^\mu$ with $0<\lambda<\mu$, the first terms of $\sin h$ are obtained
by substituting this finite sum into its power series, mixed terms such as $t^{2\lambda+\mu}$ being
collected with every other term of the same exponent. \emph{The displayed order in a polynomial
truncation is not necessarily the order of all Hahn monomials}: if $\mu\gg\lambda$, arbitrarily
many powers $t^{n\lambda}$ precede $t^\mu$. The support convention handles this automatically and
does not replace it by a single real small parameter.

\section{Notation and domain guide}\label{trigonometry:app:notation}
\begin{center}\small
\begin{longtable}{@{}L{0.26\textwidth}L{0.67\textwidth}@{}}
\toprule
Symbol & Meaning and domain\\
\midrule\endhead
$\No$, $\SC=\No[i]$ & Surreal ordered class field and its complexification.\\[0.25em]
$t^\gamma=\omega^{-\gamma}$ & Conway normal-form monomial; $v$ the least support exponent, $\lc$ the leading complex coefficient, $v(0)=+\infty$.\\[0.25em]
$\OO_{\R},\mm_{\R};\OO_{\C},\mm_{\C}$ & Finite elements and infinitesimals of $\No$ and of $\SC$.\\[0.25em]
$\st$, $\fin$ & Standard part of a finite element; the additive projection retaining all nonnegative-exponent terms of any real surreal.\\[0.25em]
$a\prec b$, $a\sim b$ & $a/b$ infinitesimal; respectively $a/b-1$ infinitesimal.\\[0.25em]
$\TT(\No)$, $\TT(\R)$ & The algebraic unit circle $\{u:u\bar u=1\}$ and its ordinary counterpart; every point has finite coordinates.\\[0.25em]
$\sin,\cos,\cis$ & Taylor-lifted functions at finite real arguments; $\cis\theta=\cos\theta+i\sin\theta$.\\[0.25em]
$C(t)$ & The rational Cayley parametrization of the circle, never the cosine.\\[0.25em]
$\ang uv$, $d(\alpha)$ & Angular metric on $\TT(\No)$; defect $\min\{\alpha,\pi-\alpha\}$ of an interior angle.\\[0.25em]
$h,\eps,\tau,\delta$ & Generic infinitesimal increment; generic positive infinitesimal; ramification/tangency parameter; triangle side gap $b+c-a$.\\[0.25em]
$\arctan,\arcsin,\arccos$ & Real inverse branches with ranges $(-\pi/2,\pi/2)$, $[-\pi/2,\pi/2]$, $[0,\pi]$; the domain of $\arctan$ is all of $\No$.\\[0.25em]
$a,b,c;\alpha,\beta,\gamma$ & Triangle sides and opposite interior angles; the angles are finite even when the sides are not.\\[0.25em]
$\Delta,R,r,s$ & Area, circumradius, inradius, semiperimeter. In \cref{trigonometry:sec:coupled} only, $s$ and $u$ denote the two coupled parameters instead.\\[0.25em]
$\exp_{\No},\log_{\No}$ & The Gonshor real exponential and logarithm.\\[0.25em]
$\Ex,\Logm$ & Infinitesimal exponential and logarithm between $\mm_{\C}$ and $1+\mm_{\C}$.\\[0.25em]
$\mathscr H$, $\mathscr V$, $E$ & The two class strips and the strip exponential $E(x+iy)=\exp_{\No}(x)\cis y$.\\[0.25em]
$\Sin,\Cos,\Tan$ & Canonical strip functions; after \cref{trigonometry:def:globaltrig} also the canonical global functions, with which they agree.\\[0.25em]
$\Pi$, $\Oz=\Pi\oplus\Z$ & Purely infinite normal forms; the omnific integer part, whose finite elements are exactly $\Z$.\\[0.25em]
$\chi$, $\Psi_\chi$, $E_\chi$ & A phase character on $\Pi$, the resulting full-real-class phase, and the resulting full-plane exponential.\\[0.25em]
$\HInt$ & Coefficientwise Hahn integral over the ring (R1), with analytic endpoint corrections; not a fine-topological Riemann integral.\\[0.25em]
(R1)--(R4) & Fixed-domain, common-domain germ, radius-free and formal coefficient rings of \cref{trigonometry:sec:rings}.\\[0.25em]
$\SO(2,\No)$, $\Rot(a,b)$ & The rotation group and its matrix $\left(\begin{smallmatrix}a&-b\\ b&a\end{smallmatrix}\right)$, $a^2+b^2=1$ (\cref{trigonometry:rot:sec:main}).\\[0.25em]
$\ellm$, $\Tm$ & Infinitesimal angle $-i\Logm(u/\st u)$ of a unit; the identity monad $\{u:\st u=1\}\cong(\mm_{\R},+)$.\\[0.25em]
$\mathsf J$, $\mathsf c$, $\Xi_a$, $\mathcal S_\rho$ & Generator of $\mathfrak{so}(2)$; complex conjugation as a reflection; the rescaling automorphism \cref{trigonometry:rot:eq:nonalgauto}; the circle of radius $\rho$.\\[0.25em]
$\cI$, $\cI^{\mathrm{div}}$, $\Mult$, $Z_E$ & Normalized period group $\{a:\Psi(2\pi a)=1\}$ of a phase, its maximal divisible subgroup, the multiplier ring, and the kernel-multiplier ring of an exponential (\cref{trigonometry:per:sec:main}).\\[0.25em]
$(\psi,\cL)$, $\pd_\Psi$, $\Zhat$ & Ordinary and infinitesimal coordinates of a character; the profinite defect $\Pi\to\Zhat/\Z$; the profinite integers.\\[0.25em]
$\Sh_c$, $[g]_0$, $F_{\Q}$ & Strong shift automorphism with $\Sh_c(\omega)=\omega+c$; the real coefficient of $\omega^0$ in an exponent; the Hahn field $\R((t^{\Q}))$.\\[0.25em]
$\Z$, $n,N,M$ & Ordinary integers and ordinary finite degrees or sample sizes, never an unspecified class of ``surreal integers.''\\
\bottomrule
\end{longtable}
\end{center}

\begingroup
\small
\setlength{\parskip}{0pt}
\begin{thebibliography}{99}
\setlength{\itemsep}{0.35em}
\raggedright
\addcontentsline{toc}{section}{References}

\bibitem{SourceAnalysis}
\emph{Surcomplex Analysis: Infinitesimal Calculus, Hahn-Coherent Holomorphy, and Contour Theory}.
User-supplied unpublished manuscript, September 21, 2026,
\texttt{surcomplex\_analysis(2).tex}.
Used for normal-form conventions, strong summability, finite Taylor lifting, coherent integrals,
the fine-local logarithm, and the distinction between fine-local and common-domain coherent
analyticity. Its global fine-analytic exponentials are the explicit predecessors of the
infinite-phase family of \cref{trigonometry:ex:phases}.

\bibitem{SourceGeometry}
\emph{Finite Geometry of Hahn-Coherent Surcomplex Zeros: Conservation of Multiplicity, Normal
Forms, and Collision-Stable Residues}. User-supplied unpublished manuscript, September 21, 2026,
in \texttt{surcomplex\_finite\_geometry.zip}.
Background for positive-support conservation of finite zero clusters and collision-stable residue
pairings; its several-variable theorems are not prerequisites for the geometry here.

\bibitem{SourceIntersections}
\emph{Finite Intersections in Surcomplex Analysis: Support-Controlled Division, Residue Duality,
and Sharp Root Stability}. User-supplied unpublished manuscript, September 21, 2026, in
\texttt{surcomplex\_intersections.zip}.
Its \S9 treats conditioned root stability with the threshold $\sigma>2\kappa$; the trigonometric
specializations in \cref{trigonometry:thm:stability} and \cref{trigonometry:thm:conditioned} are
proved independently here.

\bibitem{Conway}
John H. Conway, \emph{On Numbers and Games}, Academic Press, 1976; second edition, A K Peters,
2001. Normal forms, the ordered-field structure, and omnific integers.

\bibitem{Gonshor}
Harry Gonshor, \emph{An Introduction to the Theory of Surreal Numbers},
London Mathematical Society Lecture Note Series \textbf{110}, Cambridge University Press, 1986.
\href{https://doi.org/10.1017/CBO9780511629143}{doi:10.1017/CBO9780511629143}.
Normal forms, real closedness, and the real exponential.

\bibitem{Alling}
Norman L. Alling, \emph{Foundations of Analysis over Surreal Number Fields},
North-Holland Mathematics Studies \textbf{141}, North-Holland, 1987.

\bibitem{vdDE}
Lou van den Dries and Philip Ehrlich, ``Fields of surreal numbers and exponentiation,''
\emph{Fundamenta Mathematicae} \textbf{167} (2001), 173--188.
\href{https://doi.org/10.4064/fm167-2-3}{doi:10.4064/fm167-2-3}.
Erratum, \emph{Fundamenta Mathematicae} \textbf{168} (2001), 295--297.
The uses here are restricted analytic evaluation, real closedness and the exponential group law,
not the ordinal estimates addressed by the erratum.

\bibitem{EK}
Philip Ehrlich and Elliot Kaplan, ``Surreal ordered exponential fields,''
\emph{The Journal of Symbolic Logic} \textbf{86}, no.~3 (2021), 1066--1115.
\href{https://doi.org/10.1017/jsl.2021.59}{doi:10.1017/jsl.2021.59};
\href{https://arxiv.org/abs/2002.07739}{arXiv:2002.07739}.
In particular \S11: canonical trigonometric functions from an initial integer part, and canonical
surcomplex exponentiation; its two open questions on robustness are in \S11.1, on printed p.~34 of
arXiv v3 (June 22, 2021), whose page and section numbers are those cited in
\cref{trigonometry:per:sec:main}. The paper credits an earlier idea of using $\Oz$ for trigonometric
functions to Martin Kruskal.

\bibitem{Neumann}
B.~H. Neumann, ``On ordered division rings,''
\emph{Transactions of the American Mathematical Society} \textbf{66} (1949), 202--252.
The ordered-support summability lemma.

\bibitem{Poonen}
Bjorn Poonen, ``Maximally complete fields,''
\emph{L'Enseignement Math\'ematique} (2) \textbf{39} (1993), 87--106.
\href{https://math.mit.edu/~poonen/papers/amsval.pdf}{Author's full text}.
The algebraic-closure input is Corollary~4.

\bibitem{DLMF}
F.~W.~J. Olver et al., eds., \emph{NIST Digital Library of Mathematical Functions}, Chapter 4,
especially \S\S4.14, 4.21, 4.23--4.24. \href{https://dlmf.nist.gov/4}{dlmf.nist.gov/4}.
Consulted September 21, 2026, for ordinary trigonometric identities and inverse-function
expansions; the surreal domains and support arguments are supplied here.

\bibitem{DR}
Michael A. Dritschel and James Rovnyak, ``The operator Fej\'er--Riesz theorem,'' in
\emph{A Glimpse at Hilbert Space Operators}, Operator Theory: Advances and Applications
\textbf{207}, Birkh\"auser, 2010, 223--254.
\href{https://arxiv.org/abs/0903.3639}{arXiv:0903.3639}.
Cited for the established factorization problem and its classical context; the scalar
real-closed-field proof used here is given in full.

\bibitem{Wildberger}
Norman J. Wildberger, ``Rotor coordinates and vector trigonometry,''
\emph{Journal for Geometry and Graphics} \textbf{21} (2017), no.~1, 89--106.
\href{https://www.heldermann.de/JGG/JGG21/JGG211/jgg21009.htm}{Publisher record}.

\bibitem{Gunn}
Charles Gunn, ``Rational trigonometry via projective geometric algebra,'' 2014.
\href{https://arxiv.org/abs/1401.2371}{arXiv:1401.2371}.

\bibitem{CE}
Ovidiu Costin and Philip Ehrlich, ``Integration on the surreals,''
\emph{Advances in Mathematics} \textbf{452} (2024), article 109823.
\href{https://arxiv.org/abs/2208.14331}{arXiv:2208.14331}.
A separate integration framework, not identified with the elementary coefficientwise functional of
\cref{trigonometry:sec:arcs}.

\bibitem{BM}
Alessandro Berarducci and Vincenzo Mantova, ``Surreal numbers, derivations and transseries,''
\emph{Journal of the European Mathematical Society} \textbf{20} (2018), 339--390.
\href{https://doi.org/10.4171/JEMS/769}{doi:10.4171/JEMS/769}.
Cited to distinguish a derivation on the scalar field from the function derivatives used here.

\bibitem{BMgerms}
Alessandro Berarducci and Vincenzo Mantova, ``Transseries as germs of surreal functions,''
\emph{Transactions of the American Mathematical Society} \textbf{371} (2019), 3549--3592.
\href{https://doi.org/10.1090/tran/7428}{doi:10.1090/tran/7428};
\href{https://arxiv.org/abs/1703.01995}{arXiv:1703.01995}. Cited by source 19, especially \S3, for
the surreal analytic use of Neumann's lemma and local Taylor calculus.

\bibitem{Milne}
J.~S. Milne, \emph{Algebraic Groups: The Theory of Group Schemes of Finite Type over a Field},
Cambridge Studies in Advanced Mathematics \textbf{170}, Cambridge University Press, 2017,
especially Chapters 10 and 12. Background for the norm-one torus and its representations in
\cref{trigonometry:rot:sec:lie,trigonometry:rot:sec:representations}.

\bibitem{DK}
Hans Delfs and Manfred Knebusch, \emph{Locally Semialgebraic Spaces}, Lecture Notes in Mathematics
\textbf{1173}, Springer, 1985. See also their survey ``An introduction to locally semialgebraic
spaces,'' \emph{Rocky Mountain Journal of Mathematics} \textbf{14} (1984), 945--963,
\href{https://doi.org/10.1216/RMJ-1984-14-4-945}{doi:10.1216/RMJ-1984-14-4-945}.

\bibitem{BaroOtero}
El\'ias Baro and Margarita Otero, ``Locally definable homotopy,''
\href{https://arxiv.org/abs/0812.2112}{arXiv:0812.2112} (2008), especially Corollary 6.7 on
comparison with classical homotopy. With \cite{DK}, the source of the imported
\cref{trigonometry:rot:eq:sapi}.

\bibitem{Jerabek}
Emil Je\v{r}\'abek, ``Rigid models of Presburger arithmetic,''
\emph{Mathematical Logic Quarterly} \textbf{65} (2019), no.~1, 108--115.
\href{https://doi.org/10.1002/malq.201800019}{doi:10.1002/malq.201800019};
\href{https://arxiv.org/abs/1803.05797}{arXiv:1803.05797}. Residue maps into $\Zhat$, maximal divisible
subgroups and Leibnizian $\Z$-groups (\S\S2--4), cited by source 21.

\bibitem{KKS}
Elliot Kaplan, Lothar Sebastian Krapp and Michele Serra, ``Decomposing the automorphism group of the
surreal numbers,'' preprint, \href{https://arxiv.org/abs/2509.22374v3}{arXiv:2509.22374v3}
(April 23, 2026). Class conventions and strongly linear automorphisms (\S4).

\bibitem{Saut}
\emph{Automorphisms of the Surcomplex Numbers: Real Forms, Hahn Symmetries, Valuation, and Rigidity}.
Report of this collection, \texttt{docs/surcomplex/surcomplex-field-automorphisms/}; Theorem 5.2
(label \texttt{saut:thm:shiftflow}) is the fixed-shift flow used in \cref{trigonometry:per:sec:shifts}.
Unrefereed.

\end{thebibliography}\endgroup
\end{document}
